% Les pays au décollage économique difficile — Géographie Tle E — Cours signé OnBuch+
% Programme officiel MINESEC. Compiler avec : tectonic N15-pays-decollage-economique-difficile.tex
\newcommand{\DOCMATIERE}{Géographie}
\newcommand{\DOCNIVEAU}{Tle E}
\newcommand{\DOCMODULE}{Module 2 : La libéralisation des échanges (la mondialisation)}
\newcommand{\DOCLECON}{N15}
\newcommand{\DOCTITRE}{Les pays au décollage économique difficile}
% OnBuch+ — préambule « cours signés » ENRICHI (compilé avec tectonic / XeLaTeX)
% Système visuel « Pop » d'OnBuch+ : fond crème (jamais blanc), bordures encre
% épaisses, ombres « dures » décalées, titres Archivo Black en capitales.
% Version MATHÉMATIQUES : graphes (pgfplots/tikz), boîtes pédagogiques
% (définition, propriété, méthode, attention, le savais-tu, exercice résolu,
% exercice, TP, bilan), page de garde et signature OnBuch+ ancrée en bas.
%
% Macros attendues dans main.tex AVANT \input{preamble} :
%   \DOCMATIERE  \DOCNIVEAU  \DOCMODULE  \DOCLECON (numéro)  \DOCTITRE
\documentclass[11pt]{article}

\usepackage{fontspec}
\usepackage[french]{babel}
\usepackage[a4paper,margin=2cm,top=3.4cm,bottom=2.8cm,headheight=1.4cm,footskip=1.3cm]{geometry}
\usepackage{amsmath,amssymb,amsfonts}
\usepackage{mathtools} % \xmapsto, \coloneqq... souvent utilisés par les modèles
\usepackage{cancel} % simplifications barrées (bilans de forces)
% \abs{z} (module, valeur absolue) et \norm{u} (norme), utilisés par les modèles
\providecommand{\abs}{}\let\abs\relax\DeclarePairedDelimiter{\abs}{\lvert}{\rvert}
\providecommand{\norm}{}\let\norm\relax\DeclarePairedDelimiter{\norm}{\lVert}{\rVert}
\usepackage[table]{xcolor} % [table] : fournit aussi \rowcolors (alternance de lignes)
\usepackage{pagecolor}
\usepackage{tikz}
\usetikzlibrary{arrows.meta,positioning,calc,shapes.geometric,shapes.misc,shapes.arrows,decorations.pathmorphing,decorations.pathreplacing,decorations.markings,patterns,shadows,mindmap,trees,fit,backgrounds,matrix,intersections,angles,quotes}
\usepackage{pgfplots}
\usepackage[version=4]{mhchem} % équations nucléaires \ce{^{235}_{92}U}
\usepackage[european,siunitx]{circuitikz} % schémas électriques
\pgfplotsset{compat=1.18}
\usepgfplotslibrary{fillbetween,groupplots} % aires entre courbes (calcul intégral)
\usepackage{enumitem}
\usepackage{fancyhdr}
% [most] charge toutes les bibliothèques tcolorbox (listings, minted, raster,
% documentation...) et gonfle inutilement la mémoire TeX sur un document long
% (~300 boîtes) : on ne charge que ce qui est réellement utilisé.
\usepackage[skins,breakable]{tcolorbox}
\usepackage{graphicx}
\graphicspath{{/home/user/t-t/geo-onbuch/images/}}
\usepackage{colortbl}
\usepackage{booktabs}
\usepackage{tabularx}
\usepackage{multirow}
\usepackage{diagbox} % case d'en-tête diagonale des tableaux à double entrée
% Colonnes extensibles centrée / gauche / droite (C, L, R), souvent utilisées par les modèles
\newcolumntype{C}{>{\centering\arraybackslash}X}
\newcolumntype{L}{>{\raggedright\arraybackslash}X}
\newcolumntype{R}{>{\raggedleft\arraybackslash}X}
\usepackage{array}
\usepackage{multicol}
\usepackage{siunitx}
% parse-numbers=false : accepte aussi \SI{2,5 \times 10^{-3}}{...}
\sisetup{locale=FR,output-decimal-marker={,},group-minimum-digits=4,parse-numbers=false}
\DeclareSIUnit\ohm{\ensuremath{\symup{\Omega}}} % U+2126 absent de la police
\DeclareSIUnit\division{div} % réglages d'oscilloscope (ms/div, V/div)
\DeclareSIUnit\tr{tr} % tours (vitesse de rotation, ex. \SI{1500}{\tr\per\minute})
\DeclareSIUnit\molar{M} % concentration molaire (ex. \SI{3}{\molar}, \SI{1}{\micro\molar})
\DeclareSIUnit\an{an} % année (le modèle confond parfois avec \year, un compteur TeX) — ex. \SI{5}{\kilo\meter\per\an}
\DeclareSIUnit\FCFA{FCFA} % franc CFA (ex. \SI{500}{\FCFA\per\kilo\gram})
\DeclareSIUnit\degreeBrix{\textdegree Brix} % teneur en sucre d'un sirop/jus (réfractométrie)
\DeclareSIUnit\month{mois} % le modèle utilise parfois \month, un compteur TeX, comme unité de durée
\DeclareSIUnit\microliter{\micro\liter} % le modèle écrit parfois \microliter en un seul token
\DeclareSIUnit\million{millions} % dénombrement (ex. \SI{15}{\million\per\mL} spermatozoïdes)
\usepackage{qrcode}
% \degree : commande fréquemment utilisée par les modèles pour un angle ou une
% température en dehors de \SI{}{} (non fournie par les paquets chargés).
\providecommand{\degree}{^{\circ}}
% \qed (fin de démonstration), utilisé par les modèles sans amsthm
\providecommand{\qed}{\hfill$\square$}
% \barre : double barre des valeurs interdites dans un tableau de variations (tkz-tab)
\providecommand{\barre}{\ensuremath{\|}}
% environnement proof (amsthm non chargé) utilisé par les modèles
\ifdefined\proof\else\newenvironment{proof}[1][Démonstration]{\par\noindent\textit{#1.}\ }{\qed\par}\fi
\usepackage{lastpage}
\usepackage{hyperref}
\hypersetup{hidelinks}

% ============================================================
% Palette « Pop » OnBuch+ — mêmes hex que la classe P (plus_ui.dart)
% ============================================================
\definecolor{poporange}{HTML}{F59321}
\definecolor{popdark}{HTML}{E07A0C}
\definecolor{popcream}{HTML}{FFF6E8}
\definecolor{popink}{HTML}{1C1714}
\definecolor{poppurple}{HTML}{7A5AE0}
\definecolor{popgreen}{HTML}{2E9E6A}
\definecolor{poppink}{HTML}{E0457B}
\definecolor{popblue}{HTML}{2F7DE1}
\definecolor{popgold}{HTML}{C98A12}
\definecolor{popmuted}{HTML}{8A7F76}
\definecolor{popline}{HTML}{EADFD2}
\definecolor{popgray}{HTML}{8A7F76}
\colorlet{popgrey}{popgray}
% Alias tolérants (le modèle nomme parfois les couleurs différemment)
\colorlet{popwhite}{white}
\colorlet{popblack}{popink}
\colorlet{popred}{poppink}
\colorlet{popyellow}{popgold}
% Teintes claires (fonds de tableaux/encadrés) — le modèle nomme parfois
% les couleurs avec un suffixe L (ex. popblueL) sans les avoir définies.
\colorlet{poporangeL}{poporange!20}
\colorlet{popdarkL}{popdark!20}
\colorlet{poppurpleL}{poppurple!20}
\colorlet{popgreenL}{popgreen!20}
\colorlet{poppinkL}{poppink!20}
\colorlet{popblueL}{popblue!20}
\colorlet{popgoldL}{popgold!20}
\colorlet{popredL}{popred!20}
\colorlet{popgrayL}{popgray!20}
% Teintes claires pour fonds de tableaux / zones
\colorlet{popblueL}{popblue!10!white}
\colorlet{poporangeL}{poporange!14!white}
\colorlet{popgreenL}{popgreen!12!white}
\colorlet{poppurpleL}{poppurple!12!white}
\colorlet{poppinkL}{poppink!12!white}
\colorlet{popgoldL}{popgold!14!white}

\pagecolor{popcream}

% ============================================================
% Typographie
% ============================================================
\defaultfontfeatures{Ligatures=TeX}
\setmainfont{PlusJakartaSans-Medium.ttf}[Path=../../../fonts/,BoldFont=PlusJakartaSans-Bold.ttf]
% Maths en sans-serif (Fira Math) avec chiffres et lettres droites de Plus
% Jakarta, pour que les formules se fondent dans le texte.
\usepackage[math-style=ISO,bold-style=ISO]{unicode-math}
\setmathfont{FiraMath-Regular.otf}[Path=../../../fonts/]
\setmathfont{PlusJakartaSans-Medium.ttf}[Path=../../../fonts/,range={up/num,up/{Latin,latin},\mathup}]
\usepackage{icomma} % virgule décimale sans espace en mode math
% Caractères mathématiques Unicode absents des polices de texte (Plus Jakarta,
% Archivo Black) : rendus par leur équivalent mathématique, en texte comme en
% formule — sinon ils s'affichent en carré vide (ex. « Arithmétique dans ℤ »).
\usepackage{newunicodechar}
\newunicodechar{ℝ}{\ensuremath{\mathbb{R}}}
\newunicodechar{ℤ}{\ensuremath{\mathbb{Z}}}
\newunicodechar{ℂ}{\ensuremath{\mathbb{C}}}
\newunicodechar{ℕ}{\ensuremath{\mathbb{N}}}
\newunicodechar{ℚ}{\ensuremath{\mathbb{Q}}}
\newunicodechar{𝔻}{\ensuremath{\mathbb{D}}}
\newunicodechar{α}{\ensuremath{\alpha}}
\newunicodechar{β}{\ensuremath{\beta}}
\newunicodechar{γ}{\ensuremath{\gamma}}
\newunicodechar{δ}{\ensuremath{\delta}}
\newunicodechar{ε}{\ensuremath{\varepsilon}}
\newunicodechar{θ}{\ensuremath{\theta}}
\newunicodechar{λ}{\ensuremath{\lambda}}
\newunicodechar{μ}{\ensuremath{\mu}}
\newunicodechar{σ}{\ensuremath{\sigma}}
\newunicodechar{τ}{\ensuremath{\tau}}
\newunicodechar{φ}{\ensuremath{\varphi}}
\newunicodechar{ψ}{\ensuremath{\psi}}
\newunicodechar{ω}{\ensuremath{\omega}}
\newunicodechar{Σ}{\ensuremath{\Sigma}}
% majuscules produites par \MakeUppercase dans les titres (α-aminés -> Α)
\newunicodechar{Α}{\ensuremath{\alpha}}
\newunicodechar{Β}{\ensuremath{\beta}}
\newunicodechar{Γ}{\ensuremath{\Gamma}}
\newunicodechar{Δ}{\ensuremath{\Delta}}
\newunicodechar{Ε}{\ensuremath{\varepsilon}}
\newunicodechar{Θ}{\ensuremath{\Theta}}
\newunicodechar{Λ}{\ensuremath{\Lambda}}
\newunicodechar{Μ}{\ensuremath{\mu}}
\newunicodechar{Τ}{\ensuremath{\tau}}
\newunicodechar{Φ}{\ensuremath{\Phi}}
\newunicodechar{Ψ}{\ensuremath{\Psi}}
\newunicodechar{Ω}{\ensuremath{\Omega}}
\newunicodechar{↦}{\ensuremath{\mapsto}}
\newunicodechar{∈}{\ensuremath{\in}}
\newunicodechar{∉}{\ensuremath{\notin}}
\newunicodechar{∧}{\ensuremath{\wedge}}
\newunicodechar{⇔}{\ensuremath{\Leftrightarrow}}
\newunicodechar{⟺}{\ensuremath{\Longleftrightarrow}}
\newunicodechar{⇒}{\ensuremath{\Rightarrow}}
\newunicodechar{⟹}{\ensuremath{\Longrightarrow}}
\newunicodechar{∘}{\ensuremath{\circ}}
\newunicodechar{∪}{\ensuremath{\cup}}
\newunicodechar{∩}{\ensuremath{\cap}}
\newunicodechar{≡}{\ensuremath{\equiv}}
\newunicodechar{⊂}{\ensuremath{\subset}}
\newunicodechar{⊥}{\ensuremath{\perp}}
\newunicodechar{∃}{\ensuremath{\exists}}
\newunicodechar{∀}{\ensuremath{\forall}}
\newunicodechar{∛}{\ensuremath{\sqrt[3]{\,}}}
\newunicodechar{∜}{\ensuremath{\sqrt[4]{\,}}}
\newunicodechar{⃗}{\ensuremath{{}^{\to}}}
\newunicodechar{ⁿ}{\ifmmode^{n}\else\textsuperscript{\ensuremath{n}}\fi}
\newunicodechar{ˣ}{\ifmmode^{x}\else\textsuperscript{\ensuremath{x}}\fi}
\newunicodechar{ᵇ}{\ifmmode^{b}\else\textsuperscript{\ensuremath{b}}\fi}
\newunicodechar{ʳ}{\ifmmode^{r}\else\textsuperscript{\ensuremath{r}}\fi}
\newunicodechar{ᵘ}{\ifmmode^{u}\else\textsuperscript{\ensuremath{u}}\fi}
\newunicodechar{ʸ}{\ifmmode^{y}\else\textsuperscript{\ensuremath{y}}\fi}
\newunicodechar{ᵃ}{\ifmmode^{a}\else\textsuperscript{\ensuremath{a}}\fi}
\newunicodechar{ᵉ}{\ifmmode^{e}\else\textsuperscript{\ensuremath{e}}\fi}
\newunicodechar{ᵏ}{\ifmmode^{k}\else\textsuperscript{\ensuremath{k}}\fi}
\newunicodechar{ᵖ}{\ifmmode^{p}\else\textsuperscript{\ensuremath{p}}\fi}
\newunicodechar{ᴺ}{\ifmmode^{N}\else\textsuperscript{\ensuremath{N}}\fi}
\newunicodechar{ᶜ}{\ifmmode^{c}\else\textsuperscript{\ensuremath{c}}\fi}
\newunicodechar{ʰ}{\ifmmode^{h}\else\textsuperscript{\ensuremath{h}}\fi}
\newunicodechar{ᵠ}{\ifmmode^{q}\else\textsuperscript{\ensuremath{q}}\fi}
\newunicodechar{ₙ}{\ifmmode_{n}\else\textsubscript{\ensuremath{n}}\fi}
\newunicodechar{ₐ}{\ifmmode_{a}\else\textsubscript{\ensuremath{a}}\fi}
\newunicodechar{ᵢ}{\ifmmode_{i}\else\textsubscript{\ensuremath{i}}\fi}
\newunicodechar{ᵣ}{\ifmmode_{r}\else\textsubscript{\ensuremath{r}}\fi}
\newunicodechar{ₖ}{\ifmmode_{k}\else\textsubscript{\ensuremath{k}}\fi}
\newunicodechar{ᵦ}{\ifmmode_{\beta}\else\textsubscript{\ensuremath{\beta}}\fi}
\newunicodechar{ₓ}{\ifmmode_{x}\else\textsubscript{\ensuremath{x}}\fi}
\newfontfamily\popdisplay{ArchivoBlack-Regular.ttf}[Path=../../../fonts/]
\newfontfamily\poplogo{SpaceGrotesk-Bold.ttf}[Path=../../../fonts/]
\newfontfamily\popsemi{PlusJakartaSans-SemiBold.ttf}[Path=../../../fonts/]
\newfontfamily\popxbold{PlusJakartaSans-ExtraBold.ttf}[Path=../../../fonts/]
\newfontfamily\popxboldls{PlusJakartaSans-ExtraBold.ttf}[Path=../../../fonts/,LetterSpace=3.0]

\setlist[enumerate]{leftmargin=1.7em,itemsep=5pt,topsep=4pt}
\setlist[itemize]{leftmargin=1.7em,itemsep=4pt,topsep=4pt,label={\color{poporange}\textbullet}}
\setlength{\parindent}{0pt}
\setlength{\parskip}{5pt}
\renewcommand{\arraystretch}{1.25}

% pgfplots : style Pop par défaut
\pgfplotsset{
  every axis/.append style={axis line style={popink,line width=1pt},tick style={popink},
    grid style={popline,line width=0.5pt},label style={font=\small},tick label style={font=\footnotesize}},
  popcurve/.style={poporange,line width=1.8pt,no markers,smooth},
}
% Même style disponible dans un \draw TikZ simple (hors pgfplots)
\tikzset{popcurve/.style={poporange,line width=1.8pt}}
\tikzset{popgrid/.style={popline,very thin}}
\colorlet{popcurve}{poporange} % le nom du style est parfois utilisé comme couleur

% ============================================================
% Pill / squiggle
% ============================================================
\newcommand{\poppill}[3]{%
  \tikz[baseline=-0.55ex]\node[draw=popink,line width=1.2pt,rounded corners=2.6pt,%
    fill=#2,inner xsep=6.5pt,inner ysep=3pt]{%
    \popxboldls\fontsize{7}{7}\selectfont\color{#3}\MakeUppercase{#1}};%
}
\newcommand{\popsquiggle}[1]{%
  \tikz[baseline=-0.4ex]{
    \draw[#1,line width=2.6pt,line cap=round]
      plot [smooth] coordinates {(0,0.09) (0.34,0.01) (0.68,0.21) (1.02,0.01) (1.36,0.10)};
  }%
}
% Mot-clé surligné (marqueur orange sous le texte)
\newcommand{\cle}[1]{{\popxbold\color{popdark}#1}}

% ============================================================
% En-tête / pied
% ============================================================
\pagestyle{fancy}
\fancyhf{}
\renewcommand{\headrulewidth}{0pt}
\renewcommand{\footrulewidth}{0pt}
\lhead{\raisebox{-0.15ex}{\poplogo\fontsize{13}{13}\selectfont\color{popink}OnBuch\color{poporange}+}}
\rhead{\poppill{\DOCMATIERE\ \textbullet\ \DOCNIVEAU\ \textbullet\ Leçon \DOCLECON}{white}{popink}}
\lfoot{\popxbold\fontsize{7.6}{7.6}\selectfont\color{popmuted}\MakeUppercase{Cours signé OnBuch+}\ \textbullet\ {\popsemi\DOCTITRE}}
\rfoot{\tikz[baseline=-0.6ex]\node[rounded corners=8.5pt,fill=popink,minimum height=17pt,minimum width=17pt,inner xsep=5pt,inner sep=0]{\popxbold\fontsize{8}{8}\selectfont\color{popcream}\thepage\,/\,\pageref*{LastPage}};}

% ============================================================
% Titres
% ============================================================
\newcommand{\chaptitre}[2]{%
  \begin{tcolorbox}[enhanced,colback=poporange,colframe=popink,boxrule=2.6pt,arc=11pt,
      drop shadow=popink,left=15pt,right=15pt,top=12pt,bottom=11pt,before skip=3pt,after skip=18pt]
    \poppill{#2}{popink}{popcream}\par\vspace{9pt}
    {\raggedright\hyphenpenalty=10000\exhyphenpenalty=10000\popdisplay\fontsize{22}{24}\selectfont\color{popcream}\MakeUppercase{#1}\par}
  \end{tcolorbox}
}
% Section numérotée : pastille orange avec le numéro + titre Archivo
\newcounter{popsec}
\newcommand{\coursec}[1]{%
  \stepcounter{popsec}\setcounter{popsub}{0}%
  \par\vspace{16pt}\noindent
  \begin{minipage}{\linewidth}
  \tikz[baseline=-0.7ex]\node[draw=popink,line width=1.6pt,fill=poporange,rounded corners=4pt,minimum size=22pt,inner sep=0]{\popdisplay\fontsize{12}{12}\selectfont\color{popcream}\thepopsec};%
  \hspace{8pt}{\popdisplay\fontsize{15}{17}\selectfont\color{popink}\MakeUppercase{#1}}\par
  \vspace{4pt}\noindent{\color{poporange}\rule{36pt}{3.6pt}}
  \end{minipage}\par\nopagebreak\vspace{8pt}
}
\newcounter{popsub}
\newcommand{\courssub}[1]{%
  \stepcounter{popsub}%
  \par\vspace{9pt}\noindent{\popxbold\fontsize{11.5}{13}\selectfont\color{popink}{\color{poporange}\thepopsec.\thepopsub}\hspace{6pt}#1}\par\nopagebreak\vspace{4pt}
}

% ============================================================
% Boîtes pédagogiques — même recette : fond blanc, bordure épaisse colorée,
% ombre dure encre, badge plein. Titre optionnel [#1].
% ============================================================
\tcbset{popbase/.style={breakable,enhanced,colback=white,boxrule=2.3pt,arc=9pt,
  shadow={3pt}{-3pt}{0pt}{popink},left=14pt,right=14pt,top=11pt,bottom=11pt,before skip=11pt,after skip=15pt,
  fontupper=\popsemi\fontsize{10.6}{14.5}\selectfont\color{popink},
  fontlower=\popsemi\fontsize{10.6}{14.5}\selectfont\color{popink}}}
% Coche dessinée (le glyphe ✓ n'existe pas dans les polices du document)
\newcommand{\popcheck}[1]{\tikz[baseline=-0.5ex]\draw[#1,line width=1.6pt,line cap=round,line join=round](-0.13,0.0)--(-0.04,-0.09)--(0.13,0.1);}
\providecommand{\coche}{\popcheck{popgreen}}
\providecommand{\resultat}[1]{\par\smallskip\noindent\fcolorbox{popgreen}{popgreenL}{\parbox{\dimexpr\linewidth-2\fboxsep-2\fboxrule\relax}{\textbf{Résultat.} #1}}\par\smallskip}
\newcommand{\popboxhead}[3]{\poppill{#1}{#2}{white}\ifx\relax#3\relax\else\hspace{6pt}{\popxbold\fontsize{10.6}{12}\selectfont\color{popink}#3}\fi\par\vspace{7pt}}

\newtcolorbox{aretenir}[1][]{popbase,colframe=popgreen,before upper={\popboxhead{À retenir}{popgreen}{#1}}}
\newtcolorbox{exemplebox}[1][]{popbase,colframe=poporange,before upper={\popboxhead{Exemple}{poporange}{#1}}}
\newtcolorbox{definition}[1][]{popbase,colframe=popblue,before upper={\popboxhead{Définition}{popblue}{#1}}}
\newtcolorbox{propriete}[1][]{popbase,colframe=poppurple,before upper={\popboxhead{Propriété}{poppurple}{#1}}}
\newtcolorbox{methode}[1][]{popbase,colframe=popgold,before upper={\popboxhead{Méthode}{popgold}{#1}}}
\newtcolorbox{attention}[1][]{popbase,colframe=poppink,before upper={\popboxhead{Attention}{poppink}{#1}}}
\newtcolorbox{experience}[1][]{popbase,colframe=popblue,colback=popblueL,before upper={\popboxhead{Expérience}{popblue}{#1}}}
% « Le savais-tu ? » : carte violette pleine (contexte, culture, Cameroun)
\newtcolorbox{savaistu}[1][]{popbase,colback=poppurple,colframe=popink,
  fontupper=\popsemi\fontsize{10.4}{14.2}\selectfont\color{white},
  before upper={\poppill{Le savais-tu\,?}{white}{poppurple}\ifx\relax#1\relax\else\hspace{6pt}{\popxbold\color{white}#1}\fi\par\vspace{7pt}}}
% Exercice résolu : énoncé en haut, \tcblower puis la solution sur fond crème
\newcounter{popexo}\newcounter{popexr}
\newtcolorbox{exoresolu}[1][]{popbase,colframe=poporange,colbacklower=popcream,
  segmentation style={popink,line width=1.2pt,dashed},
  before upper={\stepcounter{popexr}\popboxhead{Exercice résolu \thepopexr}{poporange}{#1}},
  before lower={\poppill{Solution}{popink}{popcream}\par\vspace{6pt}}}
% Exercice (non résolu en place) — la correction est dans la partie Corrigés
\newtcolorbox{exercice}[1][]{popbase,colframe=popink,
  before upper={\stepcounter{popexo}\popboxhead{Exercice \thepopexo}{popink}{#1}}}
% Corrigé
\newtcolorbox{corrige}[1][]{popbase,colframe=popgreen,colback=popgreenL,
  before upper={\popboxhead{Corrigé}{popgreen}{#1}}}
% Objectifs / prérequis / bilan
\newtcolorbox{objectifs}{popbase,colframe=popink,colback=poporangeL,
  before upper={\popboxhead{Objectifs de la leçon}{popink}{}}}
\newtcolorbox{prerequis}{popbase,colframe=popink,colback=popblueL,
  before upper={\popboxhead{Prérequis}{popblue}{}}}
\newtcolorbox{bilan}{popbase,colframe=popink,colback=white,boxrule=2.8pt,
  before upper={\popboxhead{Fiche bilan}{poporange}{L'essentiel à connaître pour l'examen}}}
% Figure encadrée avec légende
\newtcolorbox{popfigure}[1][]{enhanced,breakable=false,colback=white,colframe=popline,boxrule=1.4pt,arc=8pt,
  left=10pt,right=10pt,top=10pt,bottom=8pt,before skip=10pt,after skip=12pt,halign=center,
  lower separated=false,halign lower=center,
  fontlower=\popsemi\fontsize{9}{11.5}\selectfont\color{popmuted},
  #1}
\newcommand{\legende}[1]{\par\vspace{6pt}{\popsemi\fontsize{9}{11.5}\selectfont\color{popmuted}#1}}

% ============================================================
% Signature OnBuch+
% ============================================================
\newcommand{\poplogoinline}[1]{{\poplogo\fontsize{#1}{#1}\selectfont\color{popink}OnBuch\color{poporange}+}}
\newcommand{\signatureonbuch}{%
  \begin{tcolorbox}[enhanced,colback=popink,colframe=popink,boxrule=2.2pt,arc=10pt,
      drop shadow=poporange,left=14pt,right=14pt,top=10pt,bottom=10pt,before skip=0pt,after skip=0pt]
    \tikz[baseline=-0.7ex]\node[circle,fill=popgreen,draw=popcream,line width=1.3pt,minimum size=22pt,inner sep=0]{\popcheck{white}};%
    \hspace{9pt}\begin{minipage}[c]{0.62\linewidth}
      {\popxbold\fontsize{12}{13}\selectfont\color{popcream}Signé\ \textbullet\ {\poplogo OnBuch\color{poporange}+}}\par\vspace{2pt}
      {\raggedright\popsemi\fontsize{8}{10.5}\selectfont\color{popline}\MakeUppercase{Équipe pédagogique OnBuch+}\\\MakeUppercase{Conforme au programme officiel MINESEC}\par}
    \end{minipage}\hfill
    {\poplogo\fontsize{22}{22}\selectfont\color{popcream}OnBuch\color{poporange}+}
  \end{tcolorbox}%
}

% ============================================================
% Page de garde
% #1 titre · #2 matière · #3 niveau · #4 module · #5 n° leçon · #6 durée ·
% #7 description / objectifs (texte ou itemize)
% ============================================================
\newcommand{\pagegarde}[7]{%
  \thispagestyle{empty}
  \newgeometry{a4paper,margin=1.7cm,top=1.6cm,bottom=1.4cm}%
  \begin{tikzpicture}[remember picture,overlay]
    \fill[poporange!18] ([xshift=-3.2cm,yshift=-3.6cm]current page.north east) circle (4.6cm);
    \fill[poppurple!14] ([xshift=2.4cm,yshift=7.6cm]current page.south west) circle (3.8cm);
  \end{tikzpicture}%
  {\poplogo\fontsize{30}{30}\selectfont\color{popink}OnBuch\color{poporange}+}\hfill
  \raisebox{0.6ex}{\poppill{Cours signé \textbullet\ Premium}{popink}{popcream}}\par
  \vspace{1pt}\popsquiggle{poppurple}\par
  \vspace{8pt}
  \begin{tcolorbox}[enhanced,colback=poporange,colframe=popink,boxrule=2.8pt,arc=13pt,
      drop shadow=popink,left=18pt,right=18pt,top=12pt,bottom=13pt,after skip=10pt]
    \poppill{#2 \textbullet\ #3}{popink}{popcream}\hspace{5pt}\poppill{#4}{white}{popink}\par\vspace{8pt}
    {\popdisplay\fontsize{14}{14}\selectfont\color{popink}LEÇON #5}\par\vspace{4pt}
    {\raggedright\hyphenpenalty=10000\exhyphenpenalty=10000\popdisplay\fontsize{25}{27}\selectfont\color{popcream}\MakeUppercase{#1}\par}
    \vspace{8pt}
    {\popxbold\fontsize{9.5}{9.5}\selectfont\color{popink}\MakeUppercase{Durée conseillée : #6}}
  \end{tcolorbox}
  \begin{tcolorbox}[enhanced,colback=white,colframe=popink,boxrule=2.2pt,arc=10pt,
      drop shadow=popink,left=16pt,right=16pt,top=10pt,bottom=10pt,after skip=8pt]
    \poppill{Dans cette leçon}{poppurple}{white}\par\vspace{5pt}
    \popsemi\fontsize{9.3}{12.6}\selectfont\color{popink}#7
  \end{tcolorbox}
  \noindent\begin{minipage}[t]{0.487\linewidth}
    \begin{tcolorbox}[enhanced,colback=popgreen,colframe=popink,boxrule=2.2pt,arc=10pt,
        drop shadow=popink,left=11pt,right=11pt,top=10pt,bottom=10pt,height=3.1cm,valign=center]
      \tcbox[colback=white,colframe=popink,boxrule=1.3pt,arc=5pt,boxsep=0pt,nobeforeafter,
        left=3pt,right=3pt,top=3pt,bottom=3pt,valign=center]{\qrcode[height=1.9cm]{https://play.google.com/store/apps/details?id=com.luvvix.onbuch&hl=fr}}%
      \hspace{8pt}\begin{minipage}[c]{0.5\linewidth}
        {\popdisplay\fontsize{11}{12.5}\selectfont\color{white}\MakeUppercase{Télécharge OnBuch}}\par\vspace{4pt}
        {\raggedright\popsemi\fontsize{8.4}{11}\selectfont\color{white}Gratuit \textbullet\ Google Play\\Tuteur IA, annales, concours, résultats\par}
      \end{minipage}
    \end{tcolorbox}
  \end{minipage}\hfill
  \begin{minipage}[t]{0.487\linewidth}
    \begin{tcolorbox}[enhanced,colback=poppurple,colframe=popink,boxrule=2.2pt,arc=10pt,
        drop shadow=popink,left=11pt,right=11pt,top=10pt,bottom=10pt,height=3.1cm,valign=center]
      \tikz[baseline=-0.7ex]\node[circle,fill=white,draw=popink,line width=1.3pt,minimum size=1.6cm,inner sep=0]{\popdisplay\fontsize{24}{24}\selectfont\color{poppurple}+};%
      \hspace{8pt}\begin{minipage}[c]{0.6\linewidth}
        {\popdisplay\fontsize{11}{12.5}\selectfont\color{white}\MakeUppercase{Passe à OnBuch+}}\par\vspace{4pt}
        {\raggedright\popsemi\fontsize{8.4}{11}\selectfont\color{white}Tous les cours signés, Tuteur IA illimité, fiches PDF — onglet OnBuch+ de l'app\par}
      \end{minipage}
    \end{tcolorbox}
  \end{minipage}
  \vspace{8pt}
  \popcontenu
  \vfill
  \signatureonbuch
  \restoregeometry
  \newpage
}

% Rangée « Ce cours contient » (page de garde)
\newcommand{\poptile}[3]{% #1 couleur #2 gros texte #3 légende
  \begin{tcolorbox}[enhanced,colback=#1,colframe=popink,boxrule=2pt,arc=9pt,shadow={2.5pt}{-2.5pt}{0pt}{popink},
      left=6pt,right=6pt,top=9pt,bottom=9pt,halign=center,valign=center,height=1.75cm,width=\linewidth,nobeforeafter]
    {\popdisplay\fontsize{12.5}{13}\selectfont\color{white}\MakeUppercase{#2}}\par\vspace{3pt}
    {\centering\popsemi\fontsize{7.8}{9.5}\selectfont\color{white}#3\par}
  \end{tcolorbox}}
\newcommand{\popcontenu}{%
  \par\noindent{\popxboldls\fontsize{7.5}{7.5}\selectfont\color{popmuted}CE COURS SIGNÉ CONTIENT}\par\vspace{7pt}
  \noindent\begin{minipage}[t]{0.235\linewidth}\poptile{popblue}{Cours}{complet, illustré, pas à pas}\end{minipage}\hfill
  \begin{minipage}[t]{0.235\linewidth}\poptile{poporange}{Méthodes}{et exercices résolus}\end{minipage}\hfill
  \begin{minipage}[t]{0.235\linewidth}\poptile{poppink}{Exercices}{type examen + corrigés}\end{minipage}\hfill
  \begin{minipage}[t]{0.235\linewidth}\poptile{popgreen}{Fiche bilan}{l'essentiel à retenir}\end{minipage}\par}

% Sommaire
\newtcolorbox{sommaire}{enhanced,colback=white,colframe=popink,boxrule=2.2pt,arc=10pt,
  drop shadow=popink,left=18pt,right=18pt,top=13pt,bottom=13pt,before skip=6pt,after skip=20pt,
  before upper={\poppill{Sommaire}{poppurple}{white}\par\vspace{9pt}\popsemi\fontsize{10.6}{14.5}\selectfont\color{popink}}}

% Signature de fin de cours — poussée tout en bas de la dernière page
\newcommand{\findecours}{%
  \par\vfill\nopagebreak
  \begin{center}\popsquiggle{poporange}\end{center}\vspace{4pt}
  \signatureonbuch
}
\AtBeginDocument{\def\llbracket{\mathopen{[\mkern-3.5mu[}}\def\rrbracket{\mathclose{]\mkern-3.5mu]}}} % intervalles d'entiers (glyphe ⟦ absent de la police)
% Glyphes absents de Fira Math : redéfinis à partir de glyphes disponibles
\AtBeginDocument{%
  \def\setminus{\mathbin{\backslash}}\let\smallsetminus\setminus
  \def\vdots{\mathord{\vbox{\baselineskip3.2pt\lineskiplimit0pt\kern4pt\hbox{.}\hbox{.}\hbox{.}}}}%
  \def\ddots{\mathinner{\mkern1mu\raise6.5pt\hbox{.}\mkern2mu\raise3.6pt\hbox{.}\mkern2mu\raise0.7pt\hbox{.}\mkern1mu}}%
  \def\triangle{\mathord{\scalebox{0.9}{$\symup{\Delta}$}}}%
  \def\nabla{\mathord{\raisebox{\depth}{\scalebox{1}[-1]{$\symup{\Delta}$}}}}%
  \def\checkmark{\popcheck{popdark}}%
  \def\star{\mathbin{\ast}}\def\bigstar{\mathord{\ast}}%
  \def\diamond{\mathbin{\scalebox{0.6}{$\Diamond$}}}%
  \def\bigcup{\mathop{\mathchoice{\raisebox{-0.3em}{\scalebox{1.7}{$\cup$}}}{\scalebox{1.25}{$\cup$}}{\cup}{\cup}}\displaylimits}%
  \def\bigcap{\mathop{\mathchoice{\raisebox{-0.3em}{\scalebox{1.7}{$\cap$}}}{\scalebox{1.25}{$\cap$}}{\cap}{\cap}}\displaylimits}%
  \def\lesseqqgtr{\mathrel{\lessgtr}}\def\bigoplus{\mathop{\oplus}\displaylimits}%
}
\providecommand{\ssi}{si et seulement si} % abréviation fréquente
\AtBeginDocument{\providecommand{\wideparen}{\overparen}} % arc de cercle

%% FIGFIX-BEGIN v4 — correctifs globaux des figures (docs/onbuch-generation/tools/figfix)
\usepackage{adjustbox}\usepackage{etoolbox}
% toute figure TikZ/pgfplots plus large que la ligne est réduite à la largeur disponible
\BeforeBeginEnvironment{tikzpicture}{\begin{adjustbox}{max width=\linewidth}}
\AfterEndEnvironment{tikzpicture}{\end{adjustbox}}
% légendes pgfplots lisibles par-dessus les courbes
\pgfplotsset{every axis/.append style={legend style={fill=white,fill opacity=0.92,text opacity=1,draw=black!15,inner sep=3pt}}}
% étiquettes d'arêtes (syntaxe edge["..."]) avec fond blanc
\tikzset{every edge quotes/.append style={fill=white,inner sep=1.2pt}}
%% FIGFIX-END
\begin{document}

\pagegarde{Les pays au décollage économique difficile}{Géographie}{Tle E}{Module 2 : La libéralisation des échanges (la mondialisation)}{N15}{2 h + TP 5 (2 h)}{Cette leçon étudie les pays au décollage économique difficile dans le système-monde : les pays du Golfe enrichis par la rente pétrolière, les pays d'Afrique fragilisés, les PPTE surendettés et les pays à revenus intermédiaires. Elle montre, à partir de cartes, de données chiffrées et de cas camerounais, comment ces pays s'insèrent inégalement dans la mondialisation.
\vspace{4pt}
\begin{multicols}{2}\small
\begin{itemize}
\item À la fin de cette leçon, je sais définir les notions de décollage économique, de PPTE et de pays à revenus intermédiaires
\item À la fin de cette leçon, je sais localiser sur une mappemonde les pays du Golfe, les PPTE et les pays à revenus intermédiaires
\item À la fin de cette leçon, je sais décrire les caractéristiques économiques et sociales des pays au décollage difficile
\item À la fin de cette leçon, je sais expliquer les facteurs qui freinent le décollage économique des pays d'Afrique, à partir du cas du Cameroun
\item À la fin de cette leçon, je sais interpréter des cartes, photographies et tableaux statistiques (IDH, PIB, dette) relatifs à ces pays
\item À la fin de cette leçon, je sais construire un diagramme et légender un croquis montrant l'insertion inégale des pays pauvres dans la mondialisation
\end{itemize}\end{multicols}}

\begin{sommaire}
\begin{multicols}{2}
\begin{enumerate}
\item Introduction : des pays en difficulté dans la mondialisation
\item Les pays du Golfe : un décollage fondé sur la rente pétrolière
\item Les pays d'Afrique : un décollage freiné
\item Les PPTE : le poids de la dette et de la pauvreté
\item Les pays à revenus intermédiaires : entre émergence et stagnation
\item Synthèse et TP 5 : cartographier les foyers de la mondialisation
\item Activité d'intégration
\item Méthodes et exercices résolus
\item Exercices
\item Corrigés des exercices
\item Fiche bilan
\end{enumerate}
\end{multicols}
\end{sommaire}

\begin{objectifs}
\begin{itemize}
\item À la fin de cette leçon, je sais définir les notions de décollage économique, de PPTE et de pays à revenus intermédiaires
\item À la fin de cette leçon, je sais localiser sur une mappemonde les pays du Golfe, les PPTE et les pays à revenus intermédiaires
\item À la fin de cette leçon, je sais décrire les caractéristiques économiques et sociales des pays au décollage difficile
\item À la fin de cette leçon, je sais expliquer les facteurs qui freinent le décollage économique des pays d'Afrique, à partir du cas du Cameroun
\item À la fin de cette leçon, je sais interpréter des cartes, photographies et tableaux statistiques (IDH, PIB, dette) relatifs à ces pays
\item À la fin de cette leçon, je sais construire un diagramme et légender un croquis montrant l'insertion inégale des pays pauvres dans la mondialisation
\item À la fin de cette leçon, je sais représenter sur un fond de carte mappemonde la Triade, les pays émergents et les principaux flux financiers et de marchandises (TP 5)
\item À la fin de cette leçon, je sais classer les pays du monde selon leur niveau de développement et leur place dans les échanges mondiaux
\end{itemize}
\end{objectifs}

\begin{prerequis}
\begin{itemize}
\item La notion de mondialisation et ses acteurs (leçons du module 2)
\item La Triade et les pôles de l'économie mondiale (Tle, début du module)
\item Les indicateurs de développement : PIB, PIB/habitant, IDH (Première)
\item Les pays émergents et les NPI (leçon précédente)
\item Les notions de flux commerciaux et financiers mondiaux
\end{itemize}
\end{prerequis}

\newpage

\coursec{Introduction : des pays en difficulté dans la mondialisation}

\courssub{Situation d'accroche et problématique}

À Douala, capitale économique du Cameroun, un jeune diplômé achète un smartphone fabriqué en Chine, financé par un crédit d'une banque internationale. Le même jour, dans son village de la région du Sud-Ouest, le cacao familial est chargé sur un camion : il partira à Kumba, puis vers le port de Douala, avant de rejoindre les chocolateries d'Europe. Le Cameroun est donc pleinement inséré dans la mondialisation : il importe des produits manufacturés, exporte des matières premières, emprunte sur les marchés financiers internationaux. Et pourtant, malgré son pétrole, son bois, son cacao et sa position de carrefour de l'Afrique centrale, le pays peine à \og décoller \fg{} économiquement : une part importante de sa population vit encore sous le seuil de pauvreté, son industrie reste embryonnaire et sa dette extérieure pèse lourdement sur le budget de l'État.

Pendant ce temps, à quelques milliers de kilomètres de là, les Émirats arabes unis, autrefois territoires désertiques pauvres vivant de la pêche aux perles, affichent des gratte-ciel vertigineux à Dubaï et Abou Dabi, des compagnies aériennes mondiales (Emirates, Etihad) et un PIB par habitant parmi les plus élevés du monde. Le contraste est saisissant : certains pays ont réussi leur insertion dans la mondialisation, d'autres restent bloqués au pied de la piste.

\begin{attention}[Problématique]
Pourquoi le décollage économique reste-t-il difficile pour de nombreux pays malgré leur insertion dans la mondialisation ? Quels sont les différents visages des pays au décollage économique difficile, et comment les classer ?
\end{attention}

Pour répondre à cette problématique, il faut d'abord définir précisément ce qu'est le décollage économique, puis identifier les catégories de pays concernées, enfin comprendre les critères qui permettent de les classer. Le cas du Cameroun, à la fois pays à revenus intermédiaires et pays pauvre très endetté, servira de fil conducteur.

\courssub{Qu'est-ce que le décollage économique ?}

\textbf{L'intuition d'abord.} Imagine un avion sur la piste de l'aéroport international de Douala. Tant qu'il roule au sol, il consomme beaucoup d'énergie sans vraiment avancer vite. Puis vient un moment précis où la vitesse est suffisante : l'avion quitte le sol et monte presque tout seul, sans effort supplémentaire considérable. Une économie fonctionne de la même manière : après une longue phase de préparation (accumulation de capitaux, formation des hommes, construction des infrastructures), elle peut atteindre un seuil à partir duquel la croissance s'entretient d'elle-même.

\begin{definition}[Le décollage économique]
Le \cle{décollage économique} (\emph{take-off} en anglais) est, selon l'économiste américain \cle{Walt Whitman Rostow} (ouvrage \emph{Les étapes de la croissance économique}, 1960), la phase de transition au cours de laquelle une économie traditionnelle, dominée par l'agriculture de subsistance, passe à une \cle{croissance auto-entretenue} fondée sur l'industrie. C'est la troisième des cinq étapes du développement décrites par Rostow : société traditionnelle, préconditions du décollage, \textbf{décollage}, marche vers la maturité, ère de la consommation de masse. Le décollage suppose en général un taux d'investissement dépassant environ 10 \% du revenu national, l'émergence de secteurs industriels moteurs et l'apparition d'entrepreneurs capables d'innover.
\end{definition}

\begin{exemplebox}[Des décollages réussis... et des décollages manqués]
Les \og dragons \fg{} d'Asie de l'Est (Corée du Sud, Taïwan, Singapour, Hong Kong) ont connu leur décollage dans les années 1960-1970 grâce à l'industrialisation orientée vers l'exportation. La Chine a décollé à partir des années 1980. À l'inverse, de nombreux pays africains, indépendants depuis les années 1960, n'ont jamais franchi ce seuil : leur croissance reste dépendante de l'exportation de matières premières brutes (cacao, pétrole, coton, minerais) dont les prix fluctuent sur les marchés mondiaux, et l'investissement productif demeure insuffisant.
\end{exemplebox}

\begin{savaistu}[Rostow, une théorie discutée]
La théorie des étapes de Rostow, publiée en pleine guerre froide, présentait le développement comme un chemin unique que tous les pays suivraient à leur rythme. Elle est aujourd'hui nuancée : les pays ne partent pas tous avec les mêmes atouts (ressources, situation géographique, héritage colonial), et la mondialisation crée des dépendances qui peuvent bloquer le décollage. Elle reste cependant un outil précieux pour comprendre la notion de seuil de croissance auto-entretenue.
\end{savaistu}

\courssub{La diversité des pays au décollage difficile}

Les pays au décollage économique difficile ne forment pas un ensemble homogène. On distingue classiquement quatre catégories, qui feront l'objet des sections suivantes de la leçon.

\begin{itemize}
\item \textbf{Certains pays du Golfe} (Arabie saoudite, Émirats arabes unis, Koweït, Qatar...) : leur décollage repose presque exclusivement sur la \cle{rente pétrolière et gazière}. Ils affichent des revenus par habitant très élevés, mais leur développement reste fragile car dépendant d'une ressource épuisable et d'une main-d'œuvre largement étrangère.
\item \textbf{Les pays d'Afrique} : insérés dans la mondialisation surtout comme fournisseurs de matières premières, ils connaissent un décollage \emph{freiné} par la faiblesse de l'industrialisation, l'endettement, l'instabilité politique et la forte croissance démographique.
\item \textbf{Les PPTE} (Pays Pauvres Très Endettés) : catégorie définie par le poids insoutenable de la dette extérieure ; la plupart sont africains.
\item \textbf{Les pays à revenus intermédiaires (PRI)} : pays situés entre les pays riches et les pays pauvres selon le revenu national brut par habitant ; certains sont en voie d'émergence (Brésil, Turquie), d'autres stagnent dans ce que les économistes appellent le \og piège du revenu intermédiaire \fg{}.
\end{itemize}

\begin{popfigure}
\includegraphics[width=\linewidth,height=9cm,keepaspectratio]{N15/revenus.png}
\legende{Les groupes de revenu de la Banque mondiale en 2023 (carte en anglais ; légende en bas de la carte). Violet foncé : pays à faible revenu (la plus grande partie de l'Afrique subsaharienne, l'Afghanistan, le Yémen) ; rose : revenu intermédiaire inférieur (Inde, Nigeria, Cameroun, Maroc\ldots) ; vert clair : revenu intermédiaire supérieur (Chine, Brésil, Afrique du Sud, Mexique\ldots) ; vert foncé : revenu élevé (Amérique du Nord, Europe, Japon, Australie, Russie, pays du Golfe\ldots) ; hachures : pas de données. Elle permet de localiser les catégories de pays étudiées : pays riches de la Triade, pays du Golfe, pays africains du Sud du Sahara en grande majorité à faible revenu. \\ Source : Wikimedia Commons, Our World in Data (données Banque mondiale, 2024), CC BY.}
\end{popfigure}

\begin{attention}[Ne pas confondre PPTE et PMA]
Erreur fréquente au baccalauréat : confondre \cle{PPTE} et \cle{PMA}. Les \textbf{PMA} (Pays les Moins Avancés) sont une catégorie des Nations unies définie par le \emph{niveau de développement} (faible revenu, faiblesse des ressources humaines, vulnérabilité économique). Les \textbf{PPTE} (Pays Pauvres Très Endettés) sont une catégorie définie par le FMI et la Banque mondiale selon le \emph{poids de la dette extérieure} jugée insoutenable. Un pays peut être PPTE sans être PMA : c'est précisément le cas du Cameroun.
\end{attention}

\courssub{Les critères de classement des pays}

Pour classer les pays et mesurer leur niveau de développement, les géographes et les économistes utilisent plusieurs indicateurs complémentaires. Aucun ne suffit à lui seul.

\begin{itemize}
\item Le \cle{PIB par habitant} (ou RNB par habitant) : il mesure la richesse moyenne produite par personne. La Banque mondiale classe les pays en quatre groupes : revenu faible, revenu intermédiaire inférieur, revenu intermédiaire supérieur, revenu élevé. C'est le critère qui définit les \textbf{pays à revenus intermédiaires}.
\item L'\cle{IDH} (Indice de Développement Humain), calculé par le PNUD : il combine le revenu, l'espérance de vie et le niveau d'instruction. Il varie de 0 à 1.
\item Le \cle{taux d'endettement extérieur} : rapport entre la dette extérieure et les exportations ou le PIB. C'est le critère qui définit les \textbf{PPTE}.
\item La \cle{place dans les échanges mondiaux} : part dans le commerce mondial, nature des produits échangés (matières premières brutes ou produits manufacturés à forte valeur ajoutée), attractivité pour les investissements directs étrangers.
\end{itemize}

\begin{methode}[Lire et comparer des indicateurs de développement]
Pour analyser un tableau d'indicateurs (PIB/habitant, IDH, dette), procède en quatre étapes :
\begin{enumerate}
\item \textbf{Identifier} précisément chaque indicateur : nature (stock ou flux), unité (dollars, indice de 0 à 1, \%), année et source (Banque mondiale, PNUD, FMI).
\item \textbf{Repérer les extrêmes} : le pays le plus riche et le plus pauvre du tableau ; calculer l'écart (rapport entre les deux valeurs).
\item \textbf{Croiser les indicateurs} : un PIB/habitant élevé avec un IDH moyen révèle une richesse mal répartie (cas des pays du Golfe) ; un PIB moyen avec une dette élevée révèle un PPTE.
\item \textbf{Conclure avec prudence} : les moyennes nationales masquent les inégalités internes ; précise toujours la date des données, car elles évoluent vite.
\end{enumerate}
\end{methode}

\begin{popfigure}
\begin{center}
\begin{tikzpicture}
\begin{axis}[
  ybar, bar width=0.55cm,
  width=\linewidth, height=6.2cm,
  ymin=0, ymax=50000,
  ylabel={PIB par habitant (dollars PPA, ordre de grandeur)},
  symbolic x coords={Tchad, Cameroun, Br\'esil, France, \'Emirats},
  xtick=data,
  x tick label style={font=\small},
  y tick label style={font=\small},
  nodes near coords, nodes near coords style={font=\scriptsize, /pgf/number format/fixed, /pgf/number format/precision=0},
  every node near coord/.append style={rotate=0},
  axis lines=left,
  enlarge x limits=0.15,
]
\addplot[fill=popblue, draw=popdark] coordinates {(Tchad,1600) (Cameroun,3900) (Br\'esil,16000) (France,46000) (\'Emirats,48000)};
\end{axis}
\end{tikzpicture}
\end{center}
\legende{PIB par habitant en parité de pouvoir d'achat (PPA) de quelques pays représentatifs des catégories étudiées. Ordres de grandeur arrondis (Banque mondiale, début des années 2020) : le Tchad (PPTE et PMA) est environ 30 fois plus pauvre que les Émirats arabes unis (pays du Golfe) ; le Cameroun (PRI et PPTE) se situe dans la tranche intermédiaire inférieure.}
\end{popfigure}

Le diagramme ci-dessus illustre l'ampleur des inégalités : le rapport entre les Émirats et le Tchad est d'environ 1 à 30 (en effet, $\num{48000} / \num{1600} = 30$). Mais il montre aussi que le PIB par habitant seul ne suffit pas : les Émirats affichent un niveau comparable à celui de la France, sans pour autant avoir la même diversification économique ni le même modèle de développement.

\begin{exoresolu}[Comparer deux indicateurs pour classer deux pays]
\textbf{Énoncé.} On donne, pour deux pays A et B (ordres de grandeur, début des années 2020) : pays A : PIB/habitant de \num{45000} dollars PPA, IDH de 0,91, dette extérieure faible ; pays B : PIB/habitant de \num{3900} dollars PPA, IDH de 0,58, dette extérieure longtemps jugée insoutenable. 1) Identifier la catégorie de chaque pays. 2) Expliquer pourquoi le seul PIB/habitant ne suffit pas à conclure.
\tcblower
\textbf{Solution détaillée.}
1) Le pays A présente un revenu élevé et un IDH très élevé : il s'agit d'un pays développé (type France). Le pays B présente un revenu intermédiaire inférieur, un IDH moyen et une dette insoutenable : il s'agit d'un pays à revenus intermédiaires qui est en même temps un PPTE (type Cameroun).

2) Le PIB/habitant est une moyenne : il ne dit rien de la répartition des revenus, ni de la santé et de l'éducation de la population, ni du poids de la dette. Deux pays au PIB/habitant proche peuvent avoir des IDH très différents. Il faut donc croiser plusieurs indicateurs : le PIB/habitant pour la richesse produite, l'IDH pour le bien-être réel, le taux d'endettement pour la soutenabilité financière, et la structure des échanges pour la qualité de l'insertion dans la mondialisation.
\end{exoresolu}

\courssub{Le Cameroun : un cas paradoxal, pays à revenus intermédiaires et PPTE}

Le Cameroun illustre à lui seul la complexité du classement des pays au décollage difficile. D'un côté, la Banque mondiale le range parmi les \textbf{pays à revenus intermédiaires} (tranche inférieure) : avec un PIB par habitant d'environ \num{3900} dollars PPA (ordre de grandeur), il est plus riche que la plupart de ses voisins d'Afrique centrale et de l'Ouest. Son économie est relativement diversifiée : pétrole (dont la production décline), bois, cacao, café, coton, aluminium (Alucam à Edéa), et un secteur tertiaire dynamique autour de Douala et Yaoundé.

De l'autre côté, le Cameroun a été classé \textbf{PPTE} : dans les années 1990, sa dette extérieure était devenue insoutenable, absorbant une part considérable des recettes de l'État au détriment des investissements dans la santé, l'éducation et les infrastructures.

\begin{definition}[Les PPTE]
Les \cle{PPTE} (Pays Pauvres Très Endettés, en anglais \emph{Heavily Indebted Poor Countries}, HIPC) sont des pays à faible ou moyen revenu dont la \cle{dette extérieure} est jugée \emph{insoutenable}, c'est-à-dire que son service (remboursement du capital et paiement des intérêts) dépasse durablement la capacité de paiement du pays. L'initiative PPTE, lancée en 1996 par le FMI et la Banque mondiale, prévoit l'allègement puis l'annulation d'une partie de cette dette, en échange de réformes économiques et de l'affectation des sommes dégagées à la lutte contre la pauvreté. La grande majorité des pays concernés (environ 37 à 39 pays) se trouve en Afrique subsaharienne.
\end{definition}

\begin{savaistu}[2006 : le Cameroun allégé d'une grande partie de sa dette]
Après avoir atteint le \og point de décision \fg{} en 2000, le Cameroun a atteint le \cle{point d'achèvement} de l'initiative PPTE en \textbf{avril 2006}. Cette étape a entraîné l'annulation d'une grande partie de sa dette extérieure (plusieurs milliards de dollars, notamment dans le cadre du Club de Paris et des institutions multilatérales). Les ressources ainsi libérées devaient financer l'éducation, la santé et les infrastructures. Depuis, la dette camerounaise a cependant recommencé à croître, ce qui rappelle que l'annulation de la dette ne suffit pas à déclencher le décollage.
\end{savaistu}

Ce paradoxe camerounais — assez riche pour être un PRI, assez endetté pour être un PPTE — montre que le décollage économique ne dépend pas seulement des ressources naturelles ni du niveau de revenu, mais aussi de la qualité de la gouvernance, de l'industrialisation, de la maîtrise de la dette et de la nature de l'insertion dans les échanges mondiaux. C'est tout l'objet des sections suivantes : comprendre, catégorie par catégorie, pourquoi le décollage reste difficile, avant de cartographier dans le TP 5 les véritables foyers de la mondialisation.

\begin{aretenir}[L'essentiel de l'introduction]
\begin{itemize}
\item Le \textbf{décollage économique} (Rostow) est le passage d'une économie traditionnelle à une croissance industrielle auto-entretenue.
\item Les pays au décollage difficile forment quatre catégories : pays du Golfe (rente pétrolière), pays d'Afrique (décollage freiné), PPTE (dette insoutenable), pays à revenus intermédiaires (entre émergence et stagnation).
\item On les classe grâce à des indicateurs croisés : PIB/habitant, IDH, endettement, place dans les échanges mondiaux.
\item Le Cameroun est un cas paradoxal : PRI selon la Banque mondiale, mais classé PPTE jusqu'au point d'achèvement de 2006.
\item Ne jamais confondre PPTE (critère : la dette) et PMA (critère : le niveau de développement).
\end{itemize}
\end{aretenir}

\coursec{Les pays du Golfe : un décollage fondé sur la rente pétrolière}

Nous venons de voir, en introduction, que la mondialisation n'intègre pas tous les pays de la même manière : certains en sont les moteurs, d'autres en sont marginalisés, et beaucoup occupent une position intermédiaire, inconfortable. Les pays du golfe Persique illustrent parfaitement cette catégorie des \cle{pays au décollage économique difficile} : riches à l'extrême grâce à leurs hydrocarbures, ils restent fragiles parce que leur richesse repose sur un seul produit. Comprendre leur trajectoire, c'est comprendre pourquoi la richesse ne suffit pas à faire le développement.

\courssub{Localisation : un espace stratégique au cœur du monde}

\courssub{Situer le golfe Persique}

Le \cle{golfe Persique} (appelé aussi golfe Arabo-Persique) est un bras de mer peu profond de l'océan Indien, situé entre la péninsule Arabique au sud-ouest et l'Iran au nord-est. Il communique avec le golfe d'Oman, puis avec l'océan Indien, par le \cle{détroit d'Ormuz}, passage obligé des tankers : environ un cinquième du pétrole consommé dans le monde y transite chaque jour (ordre de grandeur à retenir). C'est l'un des \cle{détroits stratégiques} majeurs de la planète, au même titre que Malacca ou Suez.

Six États riverains arabes forment le \cle{Conseil de coopération du Golfe} (CCG), organisation régionale créée en 1981 :

\begin{itemize}
\item l'\cle{Arabie Saoudite}, de loin le plus vaste et le plus peuplé (environ 35 millions d'habitants, ordre de grandeur), capitale Riyad ;
\item le \cle{Koweït}, à la tête du golfe ;
\item les \cle{Émirats arabes unis} (fédération de sept émirats, dont Abou Dabi, la capitale fédérale, et Dubaï) ;
\item le \cle{Qatar}, petite péninsule rattachée à l'Arabie Saoudite, capitale Doha ;
\item \cle{Bahreïn}, État insulaire relié à l'Arabie Saoudite par une chaussée ;
\item le sultanat d'\cle{Oman}, qui borde le golfe d'Oman et la mer d'Arabie.
\end{itemize}

\begin{methode}[Construire et légender un croquis de localisation]
Pour réussir un croquis de localisation au baccalauréat, respecte toujours quatre étapes :
\begin{enumerate}
\item \textbf{Tracer le fond} : contour simplifié mais reconnaissable de l'espace étudié (ici le golfe et les côtes environnantes).
\item \textbf{Placer les informations} : pays (limites), villes (points), objets économiques (champs pétroliers, pipelines, terminaux) avec des figurés simples et constants.
\item \textbf{Orienter et situer} : flèche du nord, échelle indicative, mers et pays voisins nommés.
\item \textbf{Rédiger titre et légende} : le titre annonce le sujet et la date ; la légende, organisée et numérotée, explique chaque figuré. Une carte sans légende est une carte sans valeur.
\end{enumerate}
\end{methode}

\begin{popfigure}
\includegraphics[width=\linewidth,height=10cm,keepaspectratio]{N15/golfe_petrole.jpg}
\legende{Le golfe Persique et ses richesses en hydrocarbures (carte en anglais). Les zones vert foncé sont les champs de pétrole, les zones roses les champs de gaz (grand champ gazier de South Pars/North Field entre l'Iran et le Qatar) ; les traits verts sont les oléoducs, les traits rouges les gazoducs (pointillés : en construction). On lit la concentration des gisements autour du golfe (Arabie saoudite, Koweït, Émirats arabes unis, Qatar, Bahreïn, Iran, Irak), les villes et capitales (Riyad, Koweït, Doha, Abou Dhabi, Manama, Mascate) et les oléoducs est-ouest qui joignent le Golfe à la mer Rouge. \\ Source : Wikimedia Commons, Energy Information Administration (États-Unis), domaine public.}
\end{popfigure}

\begin{attention}[Ne pas confondre golfe Persique et Moyen-Orient]
Le golfe Persique n'est qu'une partie du Moyen-Orient. L'Iran et l'Irak bordent aussi le golfe mais ne font pas partie du CCG. Dans une copie, précise toujours de quels pays tu parles : « les six monarchies du CCG » est la formulation la plus sûre.
\end{attention}

\courssub{La rente pétrolière et gazière : le moteur du décollage}

\courssub{Des réserves parmi les plus importantes du monde}

Le sous-sol du golfe Persique concentre une part considérable des richesses énergétiques de la planète. La région détient, à elle seule, près de la moitié des \cle{réserves prouvées} de pétrole du monde et plus d'un tiers des réserves de gaz naturel (ordres de grandeur à retenir ; les chiffres exacts varient selon les années et les sources, comme le \emph{Statistical Review of World Energy}). L'Arabie Saoudite possède le champ de \cle{Ghawar}, le plus grand champ pétrolier terrestre du monde ; le Koweït exploite le champ de \cle{Burgan}, l'un des tout premiers mondiaux ; le Qatar partage avec l'Iran le champ gazier de \cle{North Field/South Pars}, le plus grand gisement de gaz naturel connu.

Ces États jouent un rôle central au sein de l'\cle{OPEP} (Organisation des pays exportateurs de pétrole), créée en 1960. L'Arabie Saoudite en est le chef de file : en ajustant sa production, elle influence directement les cours mondiaux. Le Qatar, lui, s'est spécialisé dans le \cle{gaz naturel liquéfié} (GNL), dont il est l'un des premiers exportateurs mondiaux, grâce au terminal de Ras Laffan.

\begin{definition}[Rente pétrolière]
La \cle{rente pétrolière} désigne l'ensemble des revenus qu'un État tire de l'exportation de ses hydrocarbures (pétrole brut, produits raffinés, gaz). Elle est qualifiée de « rente » parce qu'elle provient de la vente d'une ressource naturelle et non de la production industrielle nationale : elle ne récompense ni le travail ni l'innovation du pays, mais la simple possession d'un sous-sol riche, exploité souvent avec l'aide de firmes étrangères.
\end{definition}

\courssub{Des revenus par habitant parmi les plus élevés du monde}

Grâce à cette rente, rapportée à des populations nationales souvent faibles, les pays du Golfe affichent des \cle{PIB par habitant} spectaculaires.

\begin{exemplebox}[Le Qatar, champion mondial du PIB par habitant]
Le Qatar affiche un PIB par habitant supérieur à \num{60000} USD (parfois plus de \num{80000} USD en parité de pouvoir d'achat selon les années), ce qui le place régulièrement parmi les tout premiers rangs mondiaux, aux côtés du Luxembourg et de Singapour. Les Émirats arabes unis et le Koweït se situent également dans le peloton de tête, avec des PIB par habitant généralement compris entre \num{25000} et \num{45000} USD (ordres de grandeur, variables selon le cours du baril).
\end{exemplebox}

Cette richesse se traduit concrètement : gratuité ou quasi-gratuité de l'éducation et de la santé pour les nationaux, absence d'impôt sur le revenu, infrastructures ultramodernes (autoroutes, aéroports géants, gratte-ciel comme la Burj Khalifa à Dubaï, la plus haute tour du monde avec 828 m).

\begin{popfigure}
\begin{center}
\begin{tikzpicture}
\begin{axis}[
 ybar, bar width=9pt,
 width=0.95\linewidth, height=6.5cm,
 symbolic x coords={Arabie Saoudite, Koweït, Qatar, Émirats A. U., Oman},
 xtick=data,
 x tick label style={font=\scriptsize, rotate=18, anchor=east},
 ylabel={Part en \% (ordres de grandeur)},
 ymin=0, ymax=100,
 legend style={at={(0.5,1.02)}, anchor=south, legend columns=2, font=\scriptsize},
 nodes near coords, every node near coord/.append style={font=\tiny},
]
\addplot[fill=popblue] coordinates {(Arabie Saoudite,75) (Koweït,90) (Qatar,85) (Émirats A. U.,35) (Oman,70)};
\addplot[fill=poporange] coordinates {(Arabie Saoudite,65) (Koweït,90) (Qatar,80) (Émirats A. U.,30) (Oman,75)};
\legend{Part des hydrocarbures dans les exportations, Part dans les recettes publiques}
\end{axis}
\end{tikzpicture}
\end{center}
\legende{Part approximative des hydrocarbures dans les exportations de marchandises et dans les recettes de l'État des principaux pays du Golfe (ordres de grandeur, moyennes des années 2010-2020, sources FMI et banques de données nationales). Les Émirats, grâce à Dubaï, apparaissent nettement plus diversifiés que leurs voisins.}
\end{popfigure}

\courssub{Les limites du modèle : une richesse fragile}

\courssub{Une économie peu diversifiée}

Le graphique ci-dessus le montre clairement : hors Émirats, les hydrocarbures représentent souvent entre 65 \% et 90 \% des exportations et des recettes publiques. L'industrie manufacturière et l'agriculture (quasiment impossible dans ce désert, où l'eau potable provient en grande partie du dessalement, très coûteux, de l'eau de mer) restent embryonnaires. Ces pays importent presque toutes leurs marchandises, y compris les produits alimentaires.

\courssub{Une main-d'œuvre massivement étrangère}

Deuxième limite majeure : les nationaux sont minoritaires dans leur propre pays au Qatar et aux Émirats (environ 10 à 15 \% de la population totale, ordre de grandeur). La main-d'œuvre vient d'Inde, du Pakistan, du Bangladesh, des Philippines ou d'Égypte. Cette dépendance pose des problèmes économiques (fuites de revenus par les transferts d'argent des migrants vers leurs pays d'origine), sociaux (conditions de travail parfois dénoncées par les organisations internationales, système de la \emph{kafala} qui lie le migrant à son employeur) et politiques (comment intégrer durablement ces populations sans droit à la nationalité ?).

\begin{attention}[PIB élevé $\neq$ développement achevé]
Erreur fréquente dans les copies : conclure qu'un PIB par habitant élevé signifie un pays « développé ». Faux ! Le développement se mesure aussi à l'\cle{IDH} (santé, éducation, niveau de vie), à la \cle{diversification} de l'économie, à la répartition des revenus, aux libertés politiques et à l'égalité hommes-femmes. Le PIB/habitant du Golfe est en outre gonflé par la rente et masque le faible niveau de vie d'une partie de la main-d'œuvre immigrée. Au baccalauréat, nuance toujours ton jugement.
\end{attention}

\courssub{Les stratégies de diversification : préparer l'après-pétrole}

Conscients de ces fragilités, les États du Golfe investissent massivement leur rente pour transformer leurs économies :

\begin{itemize}
\item \textbf{La finance} : Dubaï, Bahreïn et Doha sont devenues des places financières régionales.
\item \textbf{Le tourisme et les grands événements} : Dubaï attire des millions de visiteurs chaque année ; le Qatar a organisé la Coupe du monde de football en 2022, une première au Moyen-Orient.
\item \textbf{Les hubs aériens} : les compagnies \emph{Emirates} (Dubaï), \emph{Qatar Airways} (Doha) et \emph{Etihad} (Abou Dabi) font de leurs aéroports des plaques tournantes entre l'Europe, l'Asie et l'Afrique.
\item \textbf{Les fonds souverains} : l'argent de la rente est placé à l'étranger (fonds ADIA d'Abou Dabi, KIA du Koweït, QIA du Qatar), parmi les plus importants du monde, afin de garantir des revenus aux générations futures.
\item \textbf{Les plans de transformation} : la « Vision 2030 » saoudienne vise à sortir le royaume de la dépendance pétrolière (industrie, tourisme, divertissement, mégaprojets comme la ville futuriste de Neom).
\end{itemize}

\begin{savaistu}[Dubaï, l'émirat qui ne vit plus du pétrole]
Contrairement à une idée répandue, Dubaï ne vit plus principalement du pétrole : les hydrocarbures représentent moins de 5 \% de son PIB. L'émirat s'est reconverti dans le commerce (son port de Jebel Ali est le premier du Moyen-Orient), la finance, l'immobilier et le tourisme. C'est Abou Dabi, voisin et capitale fédérale, qui détient l'essentiel du pétrole des Émirats. Dubaï est ainsi devenu le modèle — discuté mais réel — de la diversification dans le Golfe.
\end{savaistu}

\courssub{Un décollage « difficile » car fragile}

Le décollage des pays du Golfe reste « difficile » pour deux raisons structurelles. D'abord, la \cle{volatilité des cours} : quand le baril chute (comme en 2014-2016 ou lors de la crise de 2020), les budgets de ces États basculent dans le déficit, car ils dépendent presque entièrement de la rente. Ensuite, la \cle{transition énergétique mondiale} : la lutte contre le réchauffement climatique pousse les grandes puissances à réduire leur consommation d'hydrocarbures, ce qui menace à terme la demande et donc la rente. Leur richesse actuelle est donc une course contre la montre : diversifier avant que le pétrole ne vaille plus grand-chose.

\begin{exoresolu}[Commentaire d'un tableau statistique]
\textbf{Énoncé.} On donne pour le Koweït (ordres de grandeur) : hydrocarbures = 90 \% des exportations et 90 \% des recettes publiques ; PIB/habitant $\approx$ \num{30000} USD. Pour les Émirats : hydrocarbures $\approx$ 35 \% des exportations ; PIB/habitant $\approx$ \num{40000} USD. Compare la situation des deux pays et montre en quoi leur décollage est « difficile ».
\tcblower
\textbf{Solution détaillée.}
\begin{enumerate}
\item \textbf{Décrire} : les deux pays ont un PIB/habitant élevé, signe d'une grande richesse monétaire. Mais le Koweït dépend presque totalement des hydrocarbures (9 dollars sur 10 exportés), tandis que les Émirats n'en tirent qu'environ un tiers de leurs exportations.
\item \textbf{Interpréter} : les Émirats, grâce à Dubaï (commerce, finance, tourisme, hub aérien), ont engagé avec succès leur diversification ; le Koweït reste une économie de rente quasi pure.
\item \textbf{Conclure} : le décollage est « difficile » car la richesse du Koweït est vulnérable à la baisse des cours et à la transition énergétique, tandis que les Émirats montrent qu'une reconversion est possible — mais elle exige des investissements massifs financés... par la rente pétrolière elle-même.
\end{enumerate}
\end{exoresolu}

\begin{aretenir}[L'essentiel de la section]
\begin{itemize}
\item Les six monarchies du CCG (Arabie Saoudite, Koweït, Émirats arabes unis, Qatar, Bahreïn, Oman) contrôlent une part majeure des réserves mondiales d'hydrocarbures et pèsent dans l'OPEP.
\item Leur richesse repose sur la \cle{rente pétrolière}, qui leur donne des PIB/habitant parmi les plus élevés du monde (Qatar $>$ \num{60000} USD).
\item Mais le modèle est fragile : économie peu diversifiée, dépendance aux cours mondiaux, main-d'œuvre étrangère massive, menace de la transition énergétique.
\item Leur réponse : la diversification (finance, tourisme, hubs aériens, fonds souverains, « Vision 2030 »), dont Dubaï est le symbole.
\end{itemize}
\end{aretenir}

Cette fragilité d'un décollage fondé sur une seule ressource n'est pas propre au Golfe : elle se retrouve, sous des formes encore plus graves, sur le continent africain, que nous étudions maintenant.

\coursec{Les pays d'Afrique : un décollage freiné}

Après avoir analysé le cas particulier des pays du Golfe, dont le décollage repose sur une rente pétrolière massive mais fragile, tournons-nous vers un continent aux réalités bien plus contrastées : l'Afrique. Si le Golfe illustre une rente gérée par des États souvent stables et riches, l'Afrique présente une mosaïque de situations où l'abondance des ressources ne se traduit pas mécaniquement en développement. Comprendre pourquoi le décollage y est freiné, c'est saisir les mécanismes structurels de la mondialisation inégale.

\courssub{Un continent aux visages multiples : localisation et diversité}

L'Afrique n'est pas un pays. C'est un continent de \textbf{54 États} (55 si l'on compte le Sahara occidental, membre de l'Union africaine), d'une superficie d'environ \textbf{30,4 millions de km\textsuperscript{2}} (soit environ 55 fois la France métropolitaine), peuplé de plus de \textbf{1,4 milliard d'habitants} en 2024, soit environ 17 à 18\,\% de l'humanité. Pour analyser ses dynamiques géographiques, il faut distinguer deux grands ensembles :

\begin{itemize}
    \item \textbf{L'Afrique du Nord} (Maroc, Algérie, Tunisie, Libye, Égypte) : rattachée au monde méditerranéen et arabe, plus urbanisée, plus industrialisée, avec des revenus par habitant plus élevés (ce sont souvent des pays à revenus intermédiaires de la tranche supérieure).
    \item \textbf{L'Afrique subsaharienne} (les pays au sud du Sahara) : une diversité extrême, des géants démographiques et économiques (Nigéria, Éthiopie, Afrique du Sud, RDC, Kenya) aux petits États enclavés (Burkina Faso, Mali, Niger, Tchad, République centrafricaine).
\end{itemize}

\begin{popfigure}
\includegraphics[width=\linewidth,height=11cm,keepaspectratio]{N15/afrique_pol.jpg}
\legende{Carte politique de l'Afrique (en anglais) servant de fond pour situer les ressources d'exportation et les pays enclavés. Elle permet de repérer l'enclavement du Sahel et de l'Afrique centrale et australe (Mali, Niger, Tchad, République centrafricaine, Soudan du Sud, Ouganda, Rwanda, Burundi, Zambie, Zimbabwe, Botswana, Éthiopie\ldots, sans accès à la mer) et les grandes métropoles (Lagos, Kinshasa, Le Caire, Johannesburg, Nairobi, Luanda, Douala). Les ressources ne sont pas figurées sur la carte ; le cours les localise : hydrocarbures sur le pourtour (golfe de Guinée, Sahara, Libye), minerais dans des zones précises (cuivre et cobalt en RDC et en Zambie, or et platine en Afrique du Sud). \\ Source : Wikimedia Commons, Central Intelligence Agency (États-Unis), domaine public.}
\end{popfigure}

\begin{attention}[Piège à éviter : l'homogénéisation de l'Afrique]
Au Bac, ne parlez jamais de « l'Afrique » comme d'un bloc unique. \textbf{Distinguez toujours} :
\begin{itemize}
    \item Les \textbf{pays émergents} (Afrique du Sud, Maroc, Égypte, Kenya, Nigéria) : industrie naissante, services financiers, pôles régionaux.
    \item Les \textbf{pays rentiers} (Algérie, Angola, Gabon, Guinée équatoriale, Libye) : forte dépendance aux hydrocarbures.
    \item Les \textbf{pays agricoles exportateurs} de coton ou de cacao (Côte d'Ivoire, Ghana, Burkina Faso, Bénin, Cameroun).
    \item Les \textbf{pays enclavés et fragiles} (Tchad, Niger, Mali, Burkina Faso, RCA, Soudan du Sud) : coûts de transport prohibitifs, insécurité.
    \item Les \textbf{pays post-conflit ou en conflit} (Soudan, est de la RDC, Somalie, Sahel central).
\end{itemize}
C'est cette diversité qui explique pourquoi les moyennes continentales masquent des réalités opposées.
\end{attention}

\courssub{Des atouts réels mais mal valorisés}

L'Afrique possède des atouts majeurs que la mondialisation devrait théoriquement favoriser.

\begin{definition}[Dividende démographique]
Situation dans laquelle la part de la population en âge de travailler (15--64 ans) devient nettement supérieure à la part des personnes à charge (moins de 15 ans et plus de 64 ans). C'est une \textbf{fénêtre d'opportunité} pour la croissance, à condition que l'emploi et la formation suivent.
\end{definition}

\begin{itemize}
    \item \textbf{Ressources naturelles :} l'Afrique détient environ \textbf{30\,\% des réserves mondiales de minerais} (environ 50 à 70\,\% du cobalt selon les estimations, l'essentiel du platine, des diamants, de l'or, du cuivre, de la bauxite), environ 8\,\% du pétrole et 7\,\% du gaz mondiaux, près de 60\,\% des terres arables non encore cultivées de la planète, et domine la production de cacao (environ 60 à 70\,\% de la production mondiale, dont environ 40\,\% pour la seule Côte d'Ivoire) et une part importante du café.
    \item \textbf{Jeunesse de la population :} âge médian d'environ \textbf{19 ans} (contre plus de 42 ans en Europe). La population active devrait fortement augmenter d'ici 2050 : c'est le futur réservoir de main-d'œuvre mondiale.
    \item \textbf{Marchés en croissance :} classe moyenne en expansion (quelques centaines de millions de personnes), urbanisation rapide (environ 43 à 45\,\% de population urbaine), adoption massive du téléphone mobile.
\end{itemize}

Pourtant, la part de l'Afrique dans le PIB mondial stagne autour de \textbf{3\,\%} (pour environ 17\,\% de la population mondiale) et sa part dans le commerce mondial reste inférieure à \textbf{3\,\%}. Pourquoi un tel décalage entre les atouts et les résultats ?

\courssub{Les obstacles structurels au décollage : la « malédiction » et la dépendance}

\begin{definition}[Malédiction des matières premières (ou « maladie hollandaise »)]
Paradoxe dans lequel l'abondance de ressources naturelles (pétrole, minerais) freine la diversification économique et le développement. Mécanismes : 1) l'afflux de devises lié aux exportations de matières premières apprécie le taux de change, ce qui renchérit et étouffe l'industrie manufacturière et l'agriculture d'exportation hors ressources ; 2) la volatilité des cours rend les recettes budgétaires instables ; 3) l'effet de rente : l'État se finance sans impôts, ce qui affaiblit le contrat social et la redevabilité ; 4) les conflits pour le contrôle de la rente.
\end{definition}

\begin{exemplebox}[Le Cameroun : « l'Afrique en miniature » face au défi de la diversification]
Le Cameroun résume les potentialités et les blocages du continent.
\begin{itemize}
    \item \textbf{Ressources :} pétrole (offshore, première source de devises), bois, cacao, aluminium (Alucam à Édéa), coton, café, bananes.
    \item \textbf{Structure des exportations :} les produits bruts ou peu transformés représentent l'écrasante majorité des ventes à l'étranger (pétrole brut en tête, puis bois en grumes ou sciés, fèves de cacao, coton fibre, aluminium primaire).
    \item \textbf{Structure des importations :} dominées par les produits manufacturés (machines, véhicules, produits pharmaceutiques, plastiques) et par les carburants raffinés — paradoxe : le pays exporte du brut et importe du raffiné, faute de capacité de raffinage suffisante depuis l'incendie de la SONARA en 2019.
    \item \textbf{Bilan :} croissance moyenne d'environ 3 à 4\,\% par an dans les années 2010, insuffisante pour absorber l'arrivée massive de jeunes sur le marché du travail ; taux de pauvreté proche de 40\,\% de la population (seuil national) ; dette publique de l'ordre de 40 à 45\,\% du PIB (risque modéré mais service de la dette croissant).
\end{itemize}
\emph{Le Cameroun exporte de la richesse brute et importe de la valeur ajoutée.}
\end{exemplebox}

\begin{exoresolu}[Analyse de la structure du PIB : Afrique et Triade]
\textbf{Document :} structure sectorielle du PIB (\% de la valeur ajoutée, ordres de grandeur 2021--2023, d'après la Banque mondiale).

\begin{center}
\begin{tabularx}{\linewidth}{|l|X|X|X|}\hline
\rowcolor{popblueL}\textbf{Secteur} & \textbf{Afrique subsaharienne} & \textbf{Afrique du Nord} & \textbf{Triade (USA, UE, Japon)} \\ \hline
Agriculture & 16\,\% & 12\,\% & 1\,\% \\ \hline
Industrie (dont manufacturière) & 28\,\% (11\,\%) & 38\,\% (16\,\%) & 26\,\% (15\,\%) \\ \hline
Services & 56\,\% & 50\,\% & 73\,\% \\ \hline
\end{tabularx}
\end{center}

\textbf{Questions :}
\begin{enumerate}
    \item Construire un diagramme en barres groupées comparant les trois ensembles pour les trois secteurs.
    \item Commenter les écarts et expliquer en quoi ils illustrent le « décollage freiné » africain.
\end{enumerate}

\tcblower
\textbf{Solution détaillée :}
\begin{enumerate}
    \item \textit{Voir le diagramme ci-dessous (réalisé avec pgfplots).}
    \item \textbf{Commentaire :}
    \begin{itemize}
        \item \textbf{Agriculture :} poids environ \textbf{16 fois supérieur} en Afrique subsaharienne à celui de la Triade. Signe d'une agriculture peu productive (agriculture vivrière, faible mécanisation, faible valeur ajoutée par travailleur). Dans la Triade, l'agriculture est capitalistique et hyper-productive : peu d'emplois, beaucoup de valeur.
        \item \textbf{Industrie :} l'Afrique subsaharienne affiche une part en apparence « correcte » (28\,\%), mais \textbf{le manufacturier n'y représente que 11\,\%} du PIB (contre 15\,\% dans la Triade et 25 à 30\,\% dans l'Asie émergente). L'industrie africaine est surtout extractive (pétrole, mines, ciment, agro-industrie légère) et non transformatrice. C'est le cœur du retard : \textbf{l'absence de révolution industrielle}.
        \item \textbf{Services :} part élevée (56\,\%) mais il s'agit d'une \textbf{« tertiarisation précoce »} : services de faible productivité (commerce informel, transport, administration) et non services à haute valeur ajoutée (finance, numérique, recherche, santé spécialisée) comme dans la Triade (73\,\%).
    \end{itemize}
    \textbf{Conclusion :} la structure du PIB africain est celle d'économies \textbf{extraverties} (tournées vers l'exportation de matières premières) et \textbf{peu intégrées} (chaînes de valeur locales courtes).
\end{enumerate}
\end{exoresolu}

\begin{popfigure}
\begin{tikzpicture}
\begin{axis}[
    ybar,
    bar width=7pt,
    width=\linewidth,
    height=8cm,
    ylabel={Part dans le PIB (\%)},
    symbolic x coords={Agriculture, Industrie, Manufacturière, Services},
    xtick=data,
    x tick label style={rotate=25, anchor=east, font=\small},
    ymin=0, ymax=85,
    ymajorgrids=true,
    grid style=dashed,
    legend style={at={(0.5,-0.3)}, anchor=north, legend columns=3, font=\small},
    enlarge x limits=0.18
]
\addplot[popblue, fill=popblue!60] coordinates {(Agriculture,1) (Industrie,26) (Manufacturière,15) (Services,73)};
\addplot[poporange, fill=poporange!60] coordinates {(Agriculture,12) (Industrie,38) (Manufacturière,16) (Services,50)};
\addplot[popgreen, fill=popgreen!60] coordinates {(Agriculture,16) (Industrie,28) (Manufacturière,11) (Services,56)};
\legend{Triade, Afrique du Nord, Afrique subsaharienne}
\end{axis}
\end{tikzpicture}
\legende{Structure sectorielle du PIB : comparaison Afrique subsaharienne, Afrique du Nord et Triade (ordres de grandeur 2021--2023, d'après la Banque mondiale). L'Afrique subsaharienne se singularise par le poids de l'agriculture et la faiblesse relative de l'industrie manufacturière.}
\end{popfigure}

\courssub{Handicaps géographiques et institutionnels : l'enclavement et l'informel}

Au-delà de la structure productive, des handicaps géographiques et institutionnels pèsent lourd.

\begin{propriete}[Le coût de l'enclavement]
Les pays enclavés (Tchad, RCA, Mali, Niger, Burkina Faso, Éthiopie depuis l'indépendance de l'Érythrée, etc.) supportent des \textbf{coûts de transport très supérieurs} — souvent de 50 à 100\,\% — à ceux des pays côtiers pour accéder aux marchés mondiaux.
\begin{itemize}
    \item Dépendance aux infrastructures des pays de transit : port de Douala pour le Tchad et la RCA, port de Mombasa pour l'Ouganda, le Rwanda, le Burundi et l'est de la RDC, corridor Abidjan--Ouagadougou.
    \item Délais aux frontières (paperasserie, corruption, tracasseries) : \textbf{le temps, c'est de l'argent}, surtout pour les marchandises périssables (mangues, produits frais) ou sensibles aux cours (cacao, coton).
    \item Ordre de grandeur : acheminer un conteneur de Douala à N'Djamena (environ \SI{1800}{km} par la route) prend souvent deux à quatre semaines et coûte plusieurs milliers de dollars — parfois davantage qu'un transport maritime Douala--Shanghai de plus de \SI{10000}{km}.
\end{itemize}
\end{propriete}

\begin{methode}[Analyser une photographie de paysage économique africain (marché, port, chantier)]
Pour réussir un commentaire de photo au Bac (par exemple le marché central de Yaoundé, le Port autonome de Douala ou un chantier routier) :
\begin{enumerate}
    \item \textbf{Identifier le lieu et l'échelle :} espace urbain ou rural ? Secteur formel ou informel ? Indices : enseignes, types de bâtiments, densité, tenue vestimentaire.
    \item \textbf{Décrire les activités observées :} vente à la sauvette (étalages au sol, parapluies), commerce de gros (sacs, chariots), transport (motos-taxis, camions, pousse-pousse), transformation artisanale (décorticage, couture, soudure).
    \item \textbf{Analyser l'organisation de l'espace :} occupation de la voirie, absence de zonage, promiscuité habitat/activité, insalubrité (ordures, eaux usées).
    \item \textbf{Interpréter en géographe :} \textbf{poids de l'informel} (emploi de survie, absence de protection sociale, fuite fiscale), \textbf{faiblesse des infrastructures} (routes non bitumées, eau et électricité peu fiables), \textbf{vitalité économique} (créativité, résilience, circuits courts), \textbf{lien avec la mondialisation} (produits importés d'occasion, riz asiatique, huile de palme locale concurrencée).
    \item \textbf{Conclure :} ce paysage illustre la \textbf{dualité structurelle} : une économie formelle étroite (exportateurs, banques, administration) coexiste avec une économie informelle massive qui absorbe l'essentiel de l'emploi urbain.
\end{enumerate}
\end{methode}

\begin{aretenir}[L'informel : frein ou amortisseur ?]
L'économie informelle représente environ \textbf{80\,\% de l'emploi} et \textbf{30 à 50\,\% du PIB} en Afrique subsaharienne.
\begin{itemize}
    \item \textbf{Frein :} évasion fiscale (manque à gagner pour l'État, donc moins d'écoles, de routes, de santé), faible productivité (pas d'accès au crédit bancaire, pas de formation), précarité des travailleurs.
    \item \textbf{Amortisseur social :} seul filet de sécurité pour des millions de jeunes et de femmes ; il absorbe le choc démographique.
    \item \textbf{Enjeu :} non pas « éradiquer » mais \textbf{formaliser progressivement} (simplification fiscale, guichets uniques, paiement mobile, protection sociale adaptée).
\end{itemize}
\end{aretenir}

\courssub{Des signes d'espoir : croissance, intégration (ZLECAf) et saut technologique}

Le tableau n'est pas tout noir. Depuis les années 2000, l'Afrique a connu l'une de ses plus longues périodes de croissance soutenue (hors chocs de la Covid-19 et de la guerre en Ukraine).

\begin{itemize}
    \item \textbf{Croissance :} plusieurs économies africaines figurent régulièrement parmi les plus dynamiques du monde (Niger, Sénégal, Rwanda, Côte d'Ivoire, Éthiopie, Bénin selon les années), avec une tendance continentale de l'ordre de 3,5 à 4\,\% par an.
    \item \textbf{ZLECAf (Zone de libre-échange continentale africaine) :} entrée en vigueur commerciale en \textbf{janvier 2021}. \textbf{54 pays signataires} (tous sauf l'Érythrée). C'est la plus grande zone de libre-échange du monde \textbf{par le nombre de pays} (environ 1,4 milliard d'habitants, PIB cumulé de l'ordre de 3\,000 à 3\,500 milliards de dollars). Objectif : \textbf{supprimer les droits de douane sur environ 90\,\% des lignes tarifaires}, faciliter la circulation des personnes et des capitaux, créer des chaînes de valeur régionales (exemple : coton du Burkina Faso $\rightarrow$ textile en Côte d'Ivoire $\rightarrow$ confection au Maroc).
    \item \textbf{Révolution mobile et fintech :} M-Pesa (Kenya), Orange Money, MTN Mobile Money (Cameroun, Côte d'Ivoire, etc.). Dans plusieurs pays, la majorité des adultes possède un compte de mobile money, bien plus qu'un compte bancaire classique : inclusion financière, paiement des factures, épargne, microcrédit, transferts transfrontaliers. L'Afrique \textbf{saute l'étape de la « banque physique »}.
    \item \textbf{L'urbanisation comme opportunité :} explosion des villes intermédiaires ; demande massive en logement, transport, eau, énergie $\rightarrow$ chantiers, emplois et marché intérieur.
\end{itemize}

\begin{savaistu}[La ZLECAf : le « game changer » africain]
Lancée opérationnellement le \textbf{1\textsuperscript{er} janvier 2021} (reportée de six mois à cause de la Covid-19), la ZLECAf est le projet phare de l'Agenda 2063 de l'Union africaine.
\begin{itemize}
    \item \textbf{Chiffres clés :} 54 pays, environ 1,4 milliard d'habitants, PIB cumulé de l'ordre de 3\,400 milliards de dollars. Marché unique le plus vaste \textbf{en nombre de pays} (devant l'UE, 27 membres).
    \item \textbf{Enjeu majeur :} le commerce intra-africain ne représente qu'environ \textbf{15\,\% du commerce total de l'Afrique} (contre environ 60\,\% en Europe et en Asie). La ZLECAf vise à le porter nettement au-delà à l'horizon 2035.
    \item \textbf{Gains attendus (Banque mondiale, 2020) :} hausse du revenu régional de l'ordre de 7\,\% (environ 450 milliards de dollars), sortie de l'extrême pauvreté pour 30 à 50 millions de personnes, hausse des salaires.
    \item \textbf{Défis :} infrastructures de transport insuffisantes, barrières non tarifaires (normes, corruption, paperasserie), règles d'origine complexes, diversité monétaire (zones franc CFA, rand, autres), volonté politique réelle.
\end{itemize}
\end{savaistu}

\begin{exercice}[Entraînement Bac : le Cameroun face à la ZLECAf]
\textbf{Sujet :} « Le Cameroun, par sa position géographique et sa structure économique, a tout à gagner de la ZLECAf, mais doit lever des obstacles majeurs. »
À l'aide de vos connaissances et du cours, discutez cette affirmation en structurant votre réponse en deux parties (atouts / obstacles) et une conclusion prospective.
\textit{Indications : pensez au port de Douala, au corridor Douala--N'Djamena/Bangui, à l'industrie naissante (aluminium, ciment, agro-industrie), à la concurrence nigériane, aux infrastructures routières et énergétiques, à l'informel.}
\end{exercice}

\begin{corrige}[Exercice : le Cameroun face à la ZLECAf]
\textbf{Introduction :} situé au carrefour de l'Afrique centrale et de l'Afrique de l'Ouest, le Cameroun (« l'Afrique en miniature ») est un acteur pivot de la ZLECAf. Son décollage freiné pourrait trouver dans l'intégration continentale un levier puissant, sous conditions.

\textbf{I. Des atouts majeurs pour tirer profit de la ZLECAf}
\begin{itemize}
    \item \textbf{Porte d'entrée naturelle :} le port de Douala, premier port de la CEMAC (environ 12 millions de tonnes par an), est le débouché du Tchad, de la RCA, et dessert aussi le Congo et la Guinée équatoriale ; le corridor routier Douala--Ngaoundéré--Garoua mène vers le Tchad.
    \item \textbf{Offre exportable diversifiée :} produits bruts (bois, cacao, coton, aluminium, bananes) et premiers produits transformés (ciment, huiles et savons, boissons, textiles). Marché CEMAC captif (environ 60 millions d'habitants) et Nigéria voisin (plus de 220 millions d'habitants) à portée de main.
    \item \textbf{Main-d'œuvre et marché intérieur :} environ 28 à 30 millions d'habitants, urbanisation supérieure à 55\,\%, classe moyenne émergente, coût salarial compétitif.
    \item \textbf{Stabilité relative :} malgré les crises (régions du Nord-Ouest et du Sud-Ouest, Boko Haram à l'Extrême-Nord), l'État tient, les institutions fonctionnent, la monnaie est stable (franc CFA arrimé à l'euro).
\end{itemize}

\textbf{II. Des obstacles structurels à lever d'urgence}
\begin{itemize}
    \item \textbf{Infrastructures défaillantes :} réseau routier majoritairement non bitumé ; électricité coûteuse et instable (coupures fréquentes) ; logistique portuaire lente (temps de séjour des conteneurs long, coûts de manutention élevés).
    \item \textbf{Concurrence asymétrique :} le Nigéria (géant de plus de 220 millions d'habitants, agro-industrie, textile, ciment, pétrochimie) pourrait inonder le marché camerounais si les règles d'origine sont mal négociées ; risque de triangulation avec des produits asiatiques réexportés via Lagos ou Lomé.
    \item \textbf{Environnement des affaires :} fiscalité complexe, lourdeurs administratives, corruption perçue, justice lente ; investissements directs étrangers faibles (quelques centaines de millions de dollars par an).
    \item \textbf{Faible capacité productive :} industrie manufacturière limitée (environ 15\,\% du PIB), dépendance aux intrants importés (emballages, machines, produits chimiques) ; risque de n'être qu'une zone de transit plutôt qu'un pôle de transformation.
\end{itemize}

\textbf{Conclusion :} la ZLECAf est une \textbf{opportunité historique} pour le Cameroun de passer du statut de « pays charnière » à celui de pôle industriel et logistique de l'Afrique centrale. Mais cela suppose une \textbf{politique industrielle volontariste} (zones économiques spéciales de Kribi, Limbé et Douala réellement opérationnelles, énergie, formation technique), une \textbf{négociation ferme des règles d'origine} et une \textbf{régionalisation effective des infrastructures} (autoroute Yaoundé--Ngaoundéré, chemin de fer, barrage de Nachtigal). Sans cela, le Cameroun risque d'être une simple zone de transit pour les produits nigérians ou asiatiques.
\end{corrige}

\courssub{Synthèse : l'Afrique à la croisée des chemins}

\begin{center}
\begin{tabularx}{\linewidth}{|>{\bfseries}l|X|X|}\hline
\rowcolor{popblueL}\textbf{Dimension} & \textbf{Freins structurels} & \textbf{Leviers d'avenir} \\ \hline
\textbf{Économique} & Dépendance aux matières premières (malédiction), faible industrialisation (environ 11\,\% de manufacturier), extraversion commerciale, informel massif (environ 80\,\% de l'emploi). & ZLECAf (marché de 1,4 milliard d'habitants), diversification (agro-industrie, pharmacie, automobile), chaînes de valeur régionales, classe moyenne. \\ \hline
\textbf{Social et démographique} & Croissance démographique forte (environ 2,5\,\% par an), pression sur l'emploi, éducation et santé sous-financées, inégalités. & Dividende démographique (si éducation et emploi), urbanisation (marché), diaspora (transferts de l'ordre de 90 à 100 milliards de dollars par an, compétences). \\ \hline
\textbf{Institutionnel} & Gouvernance perfectible, instabilités (Sahel, Corne, Grands Lacs), frontières héritées de la colonisation, capacité fiscale faible (souvent moins de 15\,\% du PIB). & Union africaine (Agenda 2063), ZLECAf, réformes de l'État, numérique (administration en ligne, mobile money), jeunesse connectée et exigeante. \\ \hline
\textbf{Géographique} & Enclavement (16 pays sans littoral), climats fragiles (Sahel), maladies, réseaux de transport hérités de la colonie (radiaux vers les ports). & Corridors multimodaux, énergies renouvelables (solaire, hydroélectricité), ports en eau profonde (Kribi, Tanger Med, Lamu). \\ \hline
\end{tabularx}
\end{center}

\begin{propriete}[La condition sine qua non : la transformation structurelle]
Le décollage africain ne viendra pas de la seule croissance du PIB, mais de la \textbf{transformation structurelle} : le déplacement de la main-d'œuvre de l'agriculture vivrière et de l'informel à faible productivité vers l'industrie manufacturière et les services modernes à haute productivité. C'est la leçon de l'Asie (Japon, dragons, Chine, Vietnam). La ZLECAf est l'outil politique destiné à créer la \textbf{taille critique de marché} indispensable à l'industrialisation (économies d'échelle).
\end{propriete}

\begin{attention}[Erreur fréquente en dissertation : « L'Afrique ne se développe pas à cause de la corruption »]
La corruption est un symptôme et un frein, pas la cause unique. Plusieurs pays asiatiques (Chine, Indonésie, Vietnam) ont décollé \textbf{en dépit} d'une corruption réelle. La différence tient à l'\textbf{État développemental} : capacité à définir une stratégie industrielle, à discipliner les élites, à investir dans le capital humain et les infrastructures. Au Bac, analysez la \textbf{capacité de l'État} (extraction fiscale, régulation, planification) plutôt que de moraliser.
\end{attention}

Cette analyse de l'Afrique nous prépare à examiner la catégorie la plus fragile : les PPTE (Pays pauvres très endettés), où le poids de la dette vient s'ajouter aux handicaps structurels que nous venons de décrire.

\coursec{Les PPTE : le poids de la dette et de la pauvreté}

Nous venons de voir que les pays d'Afrique connaissent un décollage freiné : faible industrialisation, dépendance aux matières premières, marginalisation dans les flux mondiaux. Mais pourquoi ces difficultés s'enracinent-elles autant ? Une grande partie de la réponse tient en un mot : la \cle{dette}. Imagine un ménage camerounais qui emprunte pour payer les études de ses enfants, puis qui, frappé par une mauvaise récolte, doit emprunter à nouveau simplement pour rembourser le premier prêt. Très vite, tout son revenu passe dans les remboursements : plus rien pour investir, plus rien pour se soigner. À l'échelle d'un État, ce mécanisme porte un nom : le \cle{cercle vicieux de l'endettement}. C'est lui qui a conduit la communauté internationale à créer, en 1996, une catégorie particulière de pays : les \cle{PPTE}.

\courssub{Qu'est-ce qu'un PPTE ? Définition et localisation}

\begin{definition}[Les PPTE]
Les \cle{PPTE} (Pays Pauvres Très Endettés, en anglais \emph{HIPC : Heavily Indebted Poor Countries}) sont des pays à faible revenu dont l'endettement extérieur est devenu insoutenable, c'est-à-dire impossible à rembourser sans compromettre les dépenses essentielles (santé, éducation, infrastructures). Cette catégorie a été créée en \textbf{1996} par le \cle{FMI} (Fonds monétaire international) et la \cle{Banque mondiale}, dans le cadre de l'\emph{initiative PPTE}, afin d'organiser l'allègement de leur dette.
\end{definition}

L'initiative PPTE, lancée en 1996 puis renforcée en 1999, concerne environ \textbf{37 pays} (ordre de grandeur, la liste ayant légèrement évolué au fil des admissions). Leur localisation est très significative :

\begin{itemize}
\item l'immense majorité se trouve en \cle{Afrique subsaharienne} : Bénin, Burkina Faso, Cameroun, Côte d'Ivoire, Éthiopie, Ghana, Madagascar, Mali, Mozambique, Niger, République démocratique du Congo, Sénégal, Tanzanie, Tchad, Togo, Zambie, etc. ;
\item quelques pays se situent hors d'Afrique : \textbf{Haïti} (Caraïbes), la \textbf{Bolivie}, le \textbf{Honduras}, le \textbf{Nicaragua} et le \textbf{Guyana} (Amérique latine), ainsi que l'\textbf{Afghanistan} (Asie).
\end{itemize}

\begin{popfigure}
\includegraphics[width=\linewidth,height=7.5cm,keepaspectratio]{N15/ppte.png}
\legende{Les pays pauvres très endettés (PPTE) dans le monde, d'après les stades du programme de la Banque mondiale (les nuances de vert et de jaune indiquent l'avancement de chaque pays dans le programme ; gris : pays non concernés). La quasi-totalité se concentre en Afrique subsaharienne, dont le Cameroun, en position centrale sur le golfe de Guinée ; quelques cas existent en Amérique latine et dans les Caraïbes (Haïti, Bolivie, Honduras, Nicaragua, Guyana) et en Asie (Afghanistan). \textbf{À toi de faire} : retrouve sur la carte le Cameroun, le Sénégal, le Mozambique et le Nicaragua. \\ Source : Wikimedia Commons, Haytham65 (situation au 29 janvier 2021), CC BY-SA 4.0.}
\end{popfigure}

\begin{attention}[Ne pas confondre]
Tous les pays pauvres ne sont pas des PPTE, et tous les pays endettés ne sont pas pauvres ! Les États-Unis sont le pays le plus endetté du monde en valeur absolue, mais leur richesse leur permet de rembourser. Un PPTE, c'est un pays \emph{pauvre} dont la dette est \emph{insoutenable} par rapport à ses ressources. Ne confonds pas non plus PPTE et PMA (Pays les Moins Avancés), catégorie définie par l'ONU selon d'autres critères (revenu par habitant, faiblesse des ressources humaines, vulnérabilité économique), même si les deux listes se recoupent largement.
\end{attention}

\courssub{Le mécanisme de l'endettement : comment en est-on arrivé là ?}

Comprendre l'endettement des PPTE exige de remonter aux années 1960--1970, au moment des indépendances. Les jeunes États africains ont besoin de capitaux pour construire routes, écoles, barrages et usines : ils \textbf{empruntent} auprès des banques occidentales et des institutions internationales. Trois chocs successifs transforment alors ces emprunts en piège :

\begin{enumerate}
\item \textbf{Les chocs pétroliers (1973 et 1979)} : le prix du pétrole explose. Les pays importateurs de pétrole (la plupart des pays africains) voient leur facture énergétique s'envoler et doivent emprunter davantage. Dans le même temps, les taux d'intérêt mondiaux augmentent fortement au début des années 1980, ce qui alourdit mécaniquement le poids des dettes déjà contractées.
\item \textbf{La baisse des cours des matières premières} : dans les années 1980, les prix du café, du cacao, du coton ou du cuivre s'effondrent sur le marché mondial. Or ce sont les principales sources de devises de ces pays. Leurs recettes d'exportation diminuent alors que leurs dettes, libellées en dollars, restent intactes.
\item \textbf{L'effet boule de neige} : pour payer le \cle{service de la dette}, les États empruntent à nouveau, à des conditions souvent plus dures. La dette gonfle sans que l'économie ne se développe.
\end{enumerate}

\begin{definition}[Service de la dette]
Le \cle{service de la dette} est la somme annuelle qu'un pays doit verser à ses créanciers pour rembourser sa dette. Il comprend deux éléments : le \textbf{capital} (remboursement d'une fraction de la somme empruntée) et les \textbf{intérêts} (le prix de l'emprunt). On l'exprime souvent en pourcentage des recettes d'exportation ou du budget de l'État.
\end{definition}

\begin{popfigure}
\begin{tikzpicture}[
  box/.style={draw=popink, thick, rounded corners=4pt, fill=popblueL!60, text width=3.4cm, align=center, font=\small, inner sep=5pt},
  fleche/.style={-{Stealth[length=3mm]}, very thick, poporange}
]
\node[box] (a) at (0,3.2) {\textbf{1. Emprunt} pour financer le développement (routes, écoles, barrages)};
\node[box, fill=poporangeL!70] (b) at (5.4,0) {\textbf{2. Service de la dette} lourd : capital + intérêts à rembourser chaque année};
\node[box, fill=poppinkL!70] (c) at (0,-3.2) {\textbf{3. Sous-investissement} : moins de moyens pour la santé, l'éducation, les infrastructures};
\node[box, fill=poppurpleL!70] (d) at (-5.4,0) {\textbf{4. Pauvreté persistante} et recettes insuffisantes : l'État doit emprunter à nouveau};
\draw[fleche] (a.east) to[bend left=15] (b.north);
\draw[fleche] (b.south) to[bend left=15] (c.east);
\draw[fleche] (c.west) to[bend left=15] (d.south);
\draw[fleche] (d.north) to[bend left=15] (a.west);
\node[font=\small\itshape, popdark, align=center] at (0,0) {Cercle vicieux\\de l'endettement};
\end{tikzpicture}
\legende{Le cercle vicieux de l'endettement. Chaque étape provoque la suivante : la dette étouffe l'investissement, le sous-investissement entretient la pauvreté, et la pauvreté oblige à emprunter de nouveau. Les chocs extérieurs (chocs pétroliers, baisse des cours des matières premières) aggravent le mécanisme.}
\end{popfigure}

\begin{methode}[Construire un schéma de cercle vicieux]
Un schéma de cercle vicieux est un exercice classique au Baccalauréat. Voici la démarche :
\begin{enumerate}
\item \textbf{Identifier le phénomène initial} (ici : l'emprunt) et le phénomène final qui boucle le cycle (ici : la pauvreté qui impose un nouvel emprunt).
\item \textbf{Décomposer la chaîne causale} en étapes courtes, chacune formulée par un nom ou une proposition simple (pas de phrases longues dans les cases).
\item \textbf{Disposer les étapes en cercle} et les relier par des flèches orientées : chaque flèche signifie « entraîne » ou « provoque ».
\item \textbf{Vérifier la boucle} : la dernière flèche doit revenir au point de départ. Si le cycle ne se referme pas, ce n'est pas un cercle vicieux mais une simple chaîne de causes.
\item \textbf{Titrer et légender} le schéma : titre précis en haut, et si nécessaire une phrase de légende qui explicite le mécanisme.
\end{enumerate}
\end{methode}

\courssub{Les conséquences : une dette qui étouffe le développement}

Le service de la dette fonctionne comme une ponction prioritaire sur le budget de l'État : les créanciers internationaux doivent être payés avant tout. Dans de nombreux PPTE, au plus fort de la crise de la dette (années 1980--1990), le service de la dette absorbait une part considérable des recettes publiques — parfois l'équivalent, voire davantage, que les budgets cumulés de la santé et de l'éducation (ordres de grandeur variables selon les pays et les années).

Les conséquences concrètes sont dramatiques :

\begin{itemize}
\item \textbf{Santé} : moins d'hôpitaux, pénurie de médicaments et de personnel soignant, difficultés à lutter contre le paludisme, le VIH/sida ou la mortalité maternelle et infantile ;
\item \textbf{Éducation} : classes surchargées, enseignants insuffisants et mal payés, écoles délabrées, abandon des frais de scolarité compensé par des écoles payantes qui excluent les plus pauvres ;
\item \textbf{Infrastructures} : routes, ponts, réseaux d'eau et d'électricité non entretenus, ce qui renchérit les coûts de production et décourage les investisseurs ;
\item \textbf{Effet d'éviction} : l'État surendetté perd sa crédibilité financière et n'attire plus ni prêts nouveaux ni investissements privés, ce qui aggrave encore la situation.
\end{itemize}

Autrement dit, la dette ne se contente pas d'être le \emph{symptôme} du sous-développement : elle en devient une \emph{cause} active, en privant l'État des moyens d'investir dans son avenir.

\courssub{Les réponses internationales : PAS et initiative PPTE}

Face à la crise de la dette, les institutions financières internationales ont d'abord répondu par les \cle{PAS}, puis, devant leur échec relatif, par l'initiative PPTE.

\begin{definition}[Les PAS]
Les \cle{PAS} (Programmes d'Ajustement Structurel) sont des programmes de réformes économiques imposés par le FMI et la Banque mondiale, à partir des années 1980, en échange de nouveaux prêts. Ils reposent sur la \emph{stabilisation} puis l'\emph{ajustement} : réduction des dépenses publiques (gel des salaires, licenciements dans la fonction publique), dévaluation de la monnaie, privatisations des entreprises d'État, libéralisation du commerce et suppression des subventions (par exemple sur les prix des produits de première nécessité).
\end{definition}

Les PAS ont fait l'objet de \textbf{vives critiques}. En réduisant brutalement les dépenses sociales, ils ont aggravé la pauvreté et les inégalités : au Cameroun, les années 1987--1993, marquées par les plans d'ajustement, ont vu les salaires de la fonction publique fortement réduits (deux importantes révisions salariales à la baisse en 1993) et le chômage progresser. Les privatisations ont parfois bénéficié à des intérêts étrangers sans améliorer les services. Beaucoup d'économistes et d'ONG ont dénoncé des recettes uniformes (« \emph{one size fits all} »), appliquées sans tenir compte des réalités locales, et imposées de l'extérieur, ce qui pose un problème de souveraineté.

C'est pour corriger ces limites que le FMI et la Banque mondiale ont lancé en \textbf{1996} l'\textbf{initiative PPTE}, renforcée en 1999 puis complétée en 2005 par l'initiative d'allègement de la dette multilatérale (IADM). Son principe : accorder un \textbf{allègement}, voire une \textbf{annulation} partielle de la dette, aux pays qui remplissent des conditions de bonne gestion économique et qui consacrent les sommes dégagées à la lutte contre la pauvreté (santé, éducation, infrastructures rurales). Le processus comporte deux étapes : le \emph{point de décision} (le pays devient éligible et reçoit un premier allègement) puis le \emph{point d'achèvement} (l'allègement complet devient définitif et irréversible).

\begin{exemplebox}[Un effort sans précédent]
Grâce à l'initiative PPTE et à l'IADM, plus de \textbf{100 milliards de dollars} de dette ont été annulés au profit d'une trentaine de pays, dont la grande majorité en Afrique subsaharienne. Pour plusieurs pays, le service de la dette a été divisé par deux ou plus, libérant des ressources pour construire des écoles et des centres de santé. C'est l'une des plus importantes opérations d'effacement de dette de l'histoire.
\end{exemplebox}

\begin{attention}[L'annulation de la dette ne suffit pas]
Erreur fréquente dans les copies : affirmer qu'effacer la dette déclenche automatiquement le développement. C'est faux. L'allègement de la dette ne fait que \emph{libérer des marges de manœuvre} ; encore faut-il les utiliser efficacement. Sans \textbf{bonne gouvernance} (lutte contre la corruption, transparence budgétaire), sans \textbf{investissement productif} (agriculture, industrie, infrastructures) et sans stabilité politique, les sommes dégagées peuvent être dilapidées et le pays retomber dans l'endettement. La dette est un frein, mais sa suppression n'est pas un moteur : c'est une condition nécessaire, non suffisante.
\end{attention}

\courssub{Étude de cas : le Cameroun, un PPTE parmi les pays à revenus intermédiaires}

Le cas du Cameroun est particulièrement intéressant car il illustre la complexité des classifications. Avec un revenu par habitant supérieur à celui de la plupart de ses voisins, le Cameroun est classé parmi les \cle{pays à revenus intermédiaires} (tranche inférieure) par la Banque mondiale. Pourtant, l'ampleur de sa dette extérieure — héritée des grands travaux et des emprunts des années 1970--1980, puis aggravée par la crise économique de la fin des années 1980 — l'a conduit à être admis dans l'initiative PPTE.

\begin{itemize}
\item \textbf{2000} : le Cameroun atteint le \emph{point de décision} de l'initiative PPTE et bénéficie d'un premier allègement, assorti de l'élaboration d'un Document de stratégie de réduction de la pauvreté (DSRP).
\item \textbf{2006} : le pays atteint le \emph{point d'achèvement}. L'allègement devient définitif : plusieurs milliards de dollars de dette sont effacés (ordre de grandeur : environ 4 à 5 milliards USD au total, toutes initiatives confondues), ce qui ramène la dette extérieure à un niveau soutenable.
\item Les ressources libérées sont orientées vers les secteurs prioritaires : santé (lutte contre le VIH/sida et le paludisme), éducation de base, infrastructures rurales et entretien routier.
\end{itemize}

\begin{exoresolu}[Calculer l'effet d'un allègement de dette]
\textbf{Énoncé.} Un pays PPTE consacre chaque année 400 milliards de FCFA au service de sa dette. Après avoir atteint le point d'achèvement de l'initiative PPTE, ce service est réduit de 60 \%. Le pays décide d'affecter les trois quarts des sommes libérées à l'éducation et à la santé. 1) Calcule le nouveau service de la dette. 2) Calcule la somme annuelle libérée puis celle affectée à l'éducation et à la santé. 3) Explique en quelques lignes pourquoi cet allègement peut favoriser le développement, et à quelle condition.
\tcblower
\textbf{Solution.} 1) La réduction est de $400 \times \frac{60}{100} = 240$ milliards de FCFA. Le nouveau service de la dette est donc de :
\[ 400 - 240 = 160 \text{ milliards de FCFA par an.} \]
2) La somme libérée chaque année est de 240 milliards de FCFA. La part affectée à l'éducation et à la santé vaut :
\[ 240 \times \frac{3}{4} = 180 \text{ milliards de FCFA par an.} \]
3) Cet allègement libère des ressources budgétaires considérables (240 milliards de FCFA par an) qui peuvent financer la construction d'écoles et d'hôpitaux, le recrutement d'enseignants et de soignants : autant d'investissements dans le capital humain, condition du développement à long terme. Mais cet effet positif suppose une \textbf{bonne gouvernance} : si les fonds sont détournés ou mal gérés, l'allègement n'aura aucun effet durable, et le pays risque de s'endetter à nouveau.
\end{exoresolu}

\begin{savaistu}[Le jubilé de la dette]
L'annulation de la dette des pays pauvres a été portée dans l'opinion publique mondiale par le mouvement \emph{Jubilé 2000}, une coalition d'ONG et d'Églises qui réclamait l'effacement de la dette pour l'an 2000, en référence au jubilé biblique où les dettes étaient remises. Des personnalités comme le chanteur Bono ont fait pression sur les dirigeants du G8. Ce mouvement a fortement contribué au renforcement de l'initiative PPTE en 1999 et à l'initiative d'allègement multilatéral de 2005 : un bel exemple de la société civile influençant la gouvernance mondiale.
\end{savaistu}

\begin{aretenir}[L'essentiel de la section]
\begin{itemize}
\item Les \cle{PPTE} (environ 37 pays, surtout en Afrique subsaharienne, plus Haïti, la Bolivie, le Honduras, le Nicaragua, le Guyana et l'Afghanistan) sont des pays pauvres dont la dette extérieure est insoutenable ; la catégorie a été créée par le FMI et la Banque mondiale en 1996.
\item L'endettement résulte des emprunts des indépendances, aggravés par les chocs pétroliers, la hausse des taux d'intérêt et la baisse des cours des matières premières : c'est le \cle{cercle vicieux de l'endettement}.
\item Le \cle{service de la dette} (capital + intérêts) étouffe les budgets de santé, d'éducation et d'infrastructures.
\item Les réponses : les \cle{PAS} (ajustements imposés en échange de prêts, très critiqués pour leurs coûts sociaux) puis l'\textbf{initiative PPTE} (allègement/annulation de plus de 100 milliards USD), conditionnée à la bonne gestion et à la lutte contre la pauvreté.
\item Le \textbf{Cameroun}, pays à revenus intermédiaires classé PPTE, a atteint le point de décision en \textbf{2000} et le point d'achèvement en \textbf{2006}.
\item L'annulation de la dette est une condition nécessaire mais non suffisante du développement : la gouvernance et l'investissement productif restent décisifs.
\end{itemize}
\end{aretenir}

\begin{exercice}[Pour t'entraîner]
1) Définis en deux phrases les expressions « PPTE » et « service de la dette ». 2) Sur un fond de mappemonde, localise et nomme huit PPTE dont au moins deux hors d'Afrique. 3) Construis, à partir de la méthode vue en classe, un schéma de cercle vicieux montrant comment la baisse du cours du cacao peut conduire un État à un nouvel endettement. 4) Rédige un paragraphe argumenté (une quinzaine de lignes) répondant à la question : « L'annulation de la dette suffit-elle à faire décoller l'économie des PPTE ? »
\end{exercice}

\coursec{Les pays à revenus intermédiaires : entre émergence et stagnation}

Dans la section précédente, nous avons étudié les PPTE, ces pays pauvres très endettés dont le poids de la dette et de la pauvreté bloque tout décollage économique. Mais tous les pays en développement ne sont pas logés à la même enseigne : entre les PPTE et les puissances de la Triade se trouve un vaste groupe intermédiaire, celui des \cle{pays à revenu intermédiaire} (PRI). Le Cameroun en fait partie. Ces pays ont amorcé leur décollage, mais peinent à le transformer en envol véritable : ils sont à la fois en émergence et menacés de stagnation. Comment la Banque mondiale les définit-elle ? Quels progrès ont-ils accomplis ? Pourquoi restent-ils bloqués dans un « piège » ? Et quelle place ambiguë occupent-ils dans la mondialisation ?

\courssub{La définition des pays à revenus intermédiaires selon la Banque mondiale}

\textbf{Une intuition d'abord.} Imagine un marché de Mokolo à Yaoundé : tous les vendeurs ne sont ni misérables ni riches. Beaucoup gagnent de quoi vivre, épargner un peu, envoyer leurs enfants à l'école, mais sans jamais pouvoir ouvrir un grand magasin. Les pays à revenu intermédiaire ressemblent à ces vendeurs : ils ont dépassé la misère absolue, mais la richesse reste hors d'atteinte.

\begin{definition}[Pays à revenu intermédiaire (PRI)]
Un \cle{pays à revenu intermédiaire} est un pays dont le \cle{revenu national brut} (RNB) \cle{par habitant} se situe entre environ \num{1100} et \num{14000} dollars américains par an, selon les seuils fixés chaque année par la \cle{Banque mondiale}. On distingue :
\begin{itemize}
\item les \cle{PRI de la tranche inférieure} (ou « revenu intermédiaire inférieur ») : RNB/habitant compris entre environ \num{1100} et \num{4500} USD ;
\item les \cle{PRI de la tranche supérieure} (ou « revenu intermédiaire supérieur ») : RNB/habitant compris entre environ \num{4500} et \num{14000} USD.
\end{itemize}
Les seuils exacts sont révisés annuellement (ex. : exercice 2024) : il faut donc retenir des \textbf{ordres de grandeur}, non des chiffres absolus figés.
\end{definition}

\begin{exemplebox}[Des exemples à connaître pour le Bac]
\begin{itemize}
\item \textbf{PRI tranche inférieure} : le \cle{Cameroun}, la \cle{Côte d'Ivoire}, le \cle{Sénégal}, le \cle{Kenya}, le Nigeria, l'Inde.
\item \textbf{PRI tranche supérieure} : l'\cle{Afrique du Sud}, le \cle{Brésil}, le \cle{Maroc}, la Chine, la Turquie, la Thaïlande.
\end{itemize}
Le Cameroun, avec un RNB par habitant d'environ \num{1600} USD (ordre de grandeur, données récentes Banque mondiale), appartient aux PRI de la tranche inférieure. L'Afrique du Sud dépasse \num{6000} USD/habitant et le Brésil approche les \num{10000} USD/habitant : ils sont en haut de la catégorie.
\end{exemplebox}

\begin{popfigure}
\begin{center}
\begin{tikzpicture}
\begin{axis}[
    xbar, xmin=0, xmax=15000,
    width=0.95\linewidth, height=8.5cm,
    xlabel={RNB par habitant (USD, ordres de grandeur récents)},
    symbolic y coords={Sénégal, Kenya, Côte d'Ivoire, Cameroun, Maroc, Afrique du Sud, Brésil},
    ytick=data,
    nodes near coords, nodes near coords align={horizontal},
    every node near coord/.append style={font=\small},
    bar width=9pt,
    xmajorgrids=true, grid style={popline!50},
    legend style={at={(0.5,-0.15)}, anchor=north, legend columns=-1, font=\small},
]
\addplot[fill=poporange, draw=popdark] coordinates {
    (1600,Sénégal) (2100,Kenya) (2400,Côte d'Ivoire) (1600,Cameroun)
};
\addplot[fill=popblue, draw=popdark] coordinates {
    (4100,Maroc) (6500,Afrique du Sud) (10000,Brésil)
};
\addlegendentry{PRI tranche inférieure}
\addlegendentry{PRI tranche supérieure}
\draw[popdark, dashed, thick] (axis cs:4500,Sénégal) -- (axis cs:4500,Brésil);
\node[popdark, rotate=90, anchor=south, font=\scriptsize] at (axis cs:4500,Brésil) {Seuil ~4500 USD};
\end{axis}
\end{tikzpicture}
\end{center}
\legende{Diagramme en barres classant sept pays selon les catégories de revenu de la Banque mondiale (données récentes, ordres de grandeur arrondis). En orange : les PRI tranche inférieure (Cameroun, Côte d'Ivoire, Kenya, Sénégal) ; en bleu : les PRI tranche supérieure (Maroc, Afrique du Sud, Brésil). La ligne pointillée matérialise le seuil approximatif de \num{4500} USD/habitant entre les deux sous-catégories. Source : Banque mondiale, World Development Indicators.}
\end{popfigure}

\begin{methode}[Construire un diagramme en barres comparatif à partir d'un tableau statistique]
\begin{enumerate}
\item \textbf{Lire et trier le tableau} : identifier la variable étudiée (ici le RNB/habitant) et classer les pays dans l'ordre croissant ou décroissant.
\item \textbf{Choisir l'orientation} : barres horizontales si les noms de pays sont longs (plus lisibles), verticales sinon.
\item \textbf{Fixer l'échelle} : l'axe des valeurs part toujours de zéro et va jusqu'à une valeur légèrement supérieure au maximum du tableau.
\item \textbf{Tracer les barres} : longueurs strictement proportionnelles aux valeurs, largeur constante, espaces réguliers.
\item \textbf{Soigner la présentation} : titre précis (quoi, où, quand), axes gradués et légendés avec l'unité, source des données, éventuellement un code couleur pour distinguer des catégories.
\item \textbf{Rédiger une phrase d'analyse} : dégager le fait essentiel (hiérarchie, contraste, seuil).
\end{enumerate}
\end{methode}

\begin{attention}[Piège fréquent : confondre PIB et RNB, moyenne et réalité vécue]
Le PIB mesure la richesse produite sur le territoire ; le RNB ajoute les revenus nets reçus de l'étranger (transferts de la diaspora, revenus d'investissements). Au Bac, les deux indicateurs par habitant sont acceptés pour classer les pays, mais cite celui du document fourni. Surtout, le revenu \emph{par habitant} est une \textbf{moyenne} : il masque les inégalités internes. Un RNB/habitant de \num{6000} USD ne signifie pas que chaque Sud-Africain gagne \num{6000} USD ! Autre erreur : écrire que les PRI sont des « pays pauvres ». Non : ils ont déjà dépassé le stade des PMA et des PPTE.
\end{attention}

\courssub{Des progrès réels : industrialisation naissante, urbanisation et classe moyenne}

Les PRI ne sont pas des pays stagnants : depuis un demi-siècle, ils ont accompli des progrès sensibles qui les distinguent nettement des PPTE.

\textbf{Une industrialisation naissante.} Des filières industrielles se sont développées : agroalimentaire (brasseries du Cameroun -- SABC, transformation du cacao en Côte d'Ivoire -- Cémoi, Olam), textile, ciment (cimenteries de Douala et de Figuil -- Cimencam), assemblage automobile en Afrique du Sud (BMW, Ford, Volkswagen) et au Maroc (usines Renault à Tanger, Stellantis à Kénitra), aéronautique et agro-industrie au Brésil (Embraer, troisième constructeur aéronautique mondial). L'industrie reste toutefois modeste : elle représente souvent moins de 30 \% du PIB et peine à créer assez d'emplois formels.

\textbf{Une urbanisation rapide.} Les PRI connaissent un exode rural massif : plus de la moitié de leur population vit désormais en ville. Le Cameroun compte environ 58 \% d'urbains (ordre de grandeur, Banque mondiale) ; Douala dépasse 3,5 millions d'habitants (agglomération), Johannesburg et São Paulo figurent parmi les grandes métropoles mondiales. Ces villes concentrent les services modernes : banques, universités, hôpitaux, télécommunications.

\textbf{Une classe moyenne émergente.} Urbanisation et scolarisation font naître une \cle{classe moyenne} : fonctionnaires, commerçants, cadres, enseignants, qui consomment, épargnent et investissent dans l'éducation de leurs enfants.

\begin{savaistu}[La classe moyenne africaine, nouveau moteur de consommation]
Selon la Banque africaine de développement (BAD), la classe moyenne africaine — personnes vivant avec 2 à 20 USD par jour (parité de pouvoir d'achat) — est estimée à plus de \textbf{300 millions de personnes}, soit environ un tiers de la population du continent. Elle achète téléphones, téléviseurs, motos, fréquente les supermarchés (Douala compte plusieurs grandes surfaces : Mahima, Casino, U Market) et attire les investisseurs étrangers : opérateurs téléphoniques (MTN, Orange), groupes agroalimentaires et banques y voient un marché d'avenir.
\end{savaistu}

\courssub{Le piège du revenu intermédiaire : un décollage qui ne devient pas un envol}

Pourtant, malgré ces progrès, très peu de PRI parviennent à rejoindre le club des pays développés (à revenu élevé). Depuis 1960, seules quelques réussites (Corée du Sud, Singapour, Taïwan, et plus récemment le Chili) ont franchi le cap de manière durable. Les économistes parlent du \cle{piège du revenu intermédiaire}.

\begin{definition}[Le piège du revenu intermédiaire]
Le \cle{piège du revenu intermédiaire} désigne le blocage de la croissance d'un pays qui ne parvient \textbf{ni à concurrencer les pays à bas salaires} (ses salaires sont devenus trop élevés pour les productions simples et peu qualifiées) \textbf{ni à innover comme les pays riches} (son système éducatif, sa recherche et ses entreprises ne produisent pas encore de technologies de pointe). Coincé entre les deux, le pays stagne durablement dans la catégorie intermédiaire.
\end{definition}

\textbf{Le mécanisme, étape par étape.} Prenons l'exemple d'un pays qui s'est enrichi grâce au textile ou à l'assemblage :
\begin{enumerate}
\item Au départ, ses salaires très bas attirent les usines étrangères ; le pays s'industrialise et s'enrichit (croissance extensive).
\item En s'enrichissant, les salaires montent : les usines à bas coûts délocalisent alors vers des pays encore moins chers (PMA ou PPTE voisins, ou Asie du Sud-Est).
\item Pour rester compétitif, il faudrait monter en gamme : produire des biens sophistiqués, innover, investir dans le capital humain et la R\&D. Mais cela exige des ingénieurs, de la recherche, des infrastructures fiables, une bonne gouvernance, un climat des affaires stable — que le pays n'a pas encore pleinement acquis.
\item Résultat : la croissance ralentit, la productivité stagne, le pays reste « coincé au milieu ». C'est le cas du Brésil, dont l'industrie manufacturière recule face à la concurrence asiatique (désindustrialisation précoce), ou de l'Afrique du Sud, frappée par un chômage structurel massif (environ 32 \% de la population active selon la définition stricte, ordre de grandeur récent).
\end{enumerate}

\begin{exoresolu}[Expliquer pourquoi le Cameroun risque le piège du revenu intermédiaire]
Énoncé : à partir de vos connaissances, dégagez en une dizaine de lignes deux raisons pour lesquelles le Cameroun, PRI de la tranche inférieure, peine à progresser vers le statut de pays développé.
\tcblower
Solution détaillée : Premièrement, le Cameroun ne peut plus concurrencer les pays à très bas salaires (Bangladesh, Vietnam, Éthiopie) sur les productions simples : son industrie manufacturière reste étroite (agroalimentaire, ciment, bois de première transformation) et souffre du coût élevé de l'énergie (coupures fréquentes, prix du kWh industriel), des transports (état des routes, coûts logistiques Douala-Bertoua) et de la concurrence des produits importés d'Asie (textile, électronique, machines). Deuxièmement, il ne peut pas encore rivaliser avec les pays riches par l'innovation : la recherche publique est peu financée (moins de 0,5 \% du PIB), les formations techniques et supérieures restent insuffisantes et mal adaptées aux besoins des entreprises, et l'économie repose structurellement sur l'exportation de matières premières brutes (pétrole, cacao, bois, coton, aluminium) dont les prix sont fixés par les marchés mondiaux. Faute de transformation locale à forte valeur ajoutée et de montée en gamme technologique, la croissance reste dépendante des cours mondiaux et le pays risque de stagner durablement dans la catégorie des PRI de la tranche inférieure : c'est l'illustration du piège du revenu intermédiaire.
\end{exoresolu}

\courssub{De fortes inégalités internes : le dualisme économique et social}

La richesse moyenne des PRI cache de profonds contrastes internes : c'est le \cle{dualisme}.

\begin{definition}[Dualisme économique et social]
Le \cle{dualisme} désigne la coexistence, au sein d'un même pays, de deux espaces et de deux économies contrastés : d'un côté un secteur \cle{moderne} (villes dynamiques, industries formelles, banques, gratte-ciel des CBD, infrastructures de qualité), de l'autre un secteur \cle{traditionnel ou informel} (campagnes pauvres, agriculture vivrière à bas rendement, quartiers précaires, petits métiers non déclarés, sous-emploi).
\end{definition}

Dans les métropoles africaines, ce contraste saute aux yeux : à Douala, le quartier des affaires de Bonanjo aligne les tours de banques (Afriland First Bank, BICEC, UBA) et les hôtels internationaux, tandis qu'à quelques kilomètres, les quartiers précaires (zones marécageuses de la rive du Wouri comme Youpwé, Bépanda, habitats spontanés sur pentes instables) manquent d'eau potable, d'électricité fiable et d'assainissement. À Lagos (Nigeria), les tours bancaires de Victoria Island côtoient le bidonville flottant de Makoko. Dans les campagnes, l'agriculture vivrière à faible productivité maintient des millions de paysans dans la pauvreté monétaire, alimentant l'exode rural et la croissance des quartiers informels.

\begin{popfigure}
\includegraphics[width=\linewidth,height=10cm,keepaspectratio]{N15/douala_map.jpg}
\legende{Plan de Douala (fond OpenStreetMap), métropole portuaire sur l'estuaire du Wouri. On y repère, à l'ouest au bord du fleuve, Bonanjo et le quartier d'affaires, puis Bonapriso, Akwa et Bali, l'ancien et le nouvel aéroport, les grands axes (N3 vers Édéa et Yaoundé) et, plus à l'est, de vastes quartiers denses comme Bépanda, Ndokotti ou Makepe. L'élève y localise : 1. le CBD moderne de Bonanjo, qui concentre banques, administrations, port et hôtels ; 2. les quartiers résidentiels de la classe moyenne et aisée (Bonapriso, Bali, Akwa) ; 3. les quartiers populaires à habitat spontané (Youpwé, Bépanda\ldots), sièges de l'économie informelle. Le plan ne distingue pas ces types de quartiers : c'est l'élève qui les interprète. \\ Source : Wikimedia Commons, contributeurs d'OpenStreetMap, CC BY-SA 3.0.}
\end{popfigure}

\courssub{Une place ambiguë dans la mondialisation}

Les PRI occupent une position double, contradictoire, dans les échanges mondiaux.

\textbf{D'un côté, des fournisseurs de matières premières.} Comme les PPTE, ils exportent surtout des produits primaires : pétrole brut et cacao pour le Cameroun, café et cacao pour la Côte d'Ivoire, phosphates et agrumes pour le Maroc, minerais (or, platine, charbon) pour l'Afrique du Sud, soja, minerai de fer et pétrole pour le Brésil. Ils subissent donc la volatilité des cours mondiaux et la dépendance vis-à-vis des marchés de la Triade (USA, UE, Japon) et de la Chine (premier partenaire commercial de nombreux PRI africains et du Brésil).

\textbf{De l'autre, des marchés convoités.} Leur classe moyenne croissante et leur population jeune font des PRI de \cle{nouveaux marchés} pour les firmes de la Triade et des pays émergents : les multinationales y vendent voitures, téléphones, médicaments, produits alimentaires transformés, et y installent des filiales de distribution ou d'assemblage. Les PRI sont aussi des terrains de concurrence géoéconomique entre puissances : la Chine y investit massivement (infrastructures routes/ports/barrages via Exim Bank, mines, zones économiques spéciales), rivalisant avec les anciennes puissances coloniales (France, Royaume-Uni) et les nouveaux acteurs (Turquie, Inde, Brésil, pays du Golfe).

\begin{aretenir}[L'essentiel de la section]
\begin{itemize}
\item Les \cle{PRI} ont un RNB/habitant entre environ \num{1100} et \num{14000} USD (seuils Banque mondiale révisés annuellement) ; on distingue PRI tranche inférieure (Cameroun, Côte d'Ivoire, Sénégal, Kenya) et PRI tranche supérieure (Afrique du Sud, Brésil, Maroc, Chine).
\item Ils ont accompli de réels progrès : industrialisation naissante, urbanisation rapide, classe moyenne émergente (moteur de la consommation intérieure).
\item Ils risquent le \cle{piège du revenu intermédiaire} : ni assez compétitifs face aux pays à bas salaires, ni assez innovants face aux pays riches (faible R\&D, capital humain insuffisant, gouvernance perfectible).
\item Ils sont marqués par un fort \cle{dualisme} : villes modernes (CBD, quartiers résidentiels) face à campagnes pauvres et quartiers précaires informels.
\item Leur place dans la mondialisation est ambiguë : fournisseurs de matières premières \emph{et} nouveaux marchés convoités par la Triade et les émergents (concurrence d'influence).
\end{itemize}
\end{aretenir}

\begin{exercice}[Pour t'entraîner]
\begin{enumerate}
\item Définis en deux phrases le « piège du revenu intermédiaire » et illustre-le avec l'exemple du Brésil ou de l'Afrique du Sud.
\item À partir du diagramme en barres de cette leçon, rédige un commentaire de cinq lignes dégageant la hiérarchie des revenus entre les sept pays présentés et la fracture entre les deux tranches.
\item Sur un croquis de Douala inspiré de celui de la leçon, place et légende les trois espaces du dualisme urbain (titre, orientation, échelle indicative, légende numérotée).
\end{enumerate}
\end{exercice}

\begin{corrige}[Exercice]
1. Le piège du revenu intermédiaire est le blocage durable de la croissance d'un pays qui ne peut plus concurrencer les pays à bas salaires sur les productions simples, sans être capable d'innover comme les pays développés. Le Brésil l'illustre : ses salaires ont augmenté, son industrie manufacturière recule face à la concurrence asiatique (désindustrialisation précoce), mais sa recherche et ses technologies ne lui permettent pas encore de rivaliser durablement avec la Triade sur les biens à haute valeur ajoutée.
2. Le commentaire doit hiérarchiser (Brésil en tête \textasciitilde 10 000 USD, puis Afrique du Sud \textasciitilde 6 500 USD, Maroc \textasciitilde 4 100 USD, puis Côte d'Ivoire \textasciitilde 2 400 USD, Kenya \textasciitilde 2 100 USD, Cameroun et Sénégal \textasciitilde 1 600 USD), souligner le contraste net entre PRI tranche supérieure et inférieure (rapport d'environ 1 à 6 entre le Brésil et le Cameroun) et conclure sur la diversité interne du groupe des PRI.
3. Le croquis doit comporter un titre précis (ex. « Douala : dualisme urbain »), une orientation (flèche du nord), une échelle indicative (ex. 2 km) et une légende numérotée distinguant clairement le CBD moderne (1), les quartiers résidentiels aisés (2) et les quartiers précaires informels (3).
\end{corrige}

\begin{experience}[TP 5 guidé : Les foyers de la mondialisation — Carte mentale et flux]
\textbf{Objectif :} Sur un fond de carte mappemonde (fourni ou tracé schématiquement), représenter et localiser la Triade, les pays émergents (BRICS, PRI supérieurs) ainsi que les principaux flux financiers et de marchandises.

\textbf{Matériel :} Fond de carte mappemonde A3 (projection Robinson ou Mercator simplifiée), crayons de couleur (bleu, rouge, vert, orange), feutres fins, règle, atlas.

\textbf{Protocole (étapes) :}
\begin{enumerate}
\item \textbf{Repérage des pôles de puissance (Triade) :} En \textbf{bleu}, colorier/hachurer les trois pôles historiques : Amérique du Nord (USA, Canada), Union Européenne (UE), Japon. Légender : « Triade : pôles historiques de la mondialisation (PIB, commerce, finance, innovation) ».
\item \textbf{Repérage des nouveaux pôles (Émergents/PRI supérieurs) :} En \textbf{rouge}, localiser et nommer : Chine, Inde, Brésil, Russie, Afrique du Sud (BRICS), + Mexique, Indonésie, Turquie, Corée du Sud, Arabie Saoudite. Légender : « Pays émergents / PRI supérieurs : nouveaux moteurs de croissance ».
\item \textbf{Flux de marchandises (commerce) :} Au \textbf{feutre noir}, tracer 4 à 5 flèches épaisses (largeur proportionnelle aux volumes) pour les grands axes : Transpacifique (Chine $\leftrightarrow$ USA), Transatlantique (UE $\leftrightarrow$ USA), Europe-Asie (via Suez), Amérique du Sud-Chine (soja, minerai), Afrique-Chine (minerais, pétrole). Légender : « Principaux flux de marchandises (conteneurs, vrac) ».
\item \textbf{Flux financiers (IDE, portefeuille) :} Au \textbf{feutre vert}, tracer des flèches pointillées pour les Investissements Directs à l'Étranger (IDE) sortants de la Triade et de la Chine vers les PRI et PMA (Afrique, Asie du Sud-Est, Amérique Latine). Légender : « Flux d'IDE (usines, infrastructures, mines) ».
\item \textbf{Mise en page :} Titre centré en haut : « Les foyers et les flux de la mondialisation (année récente) ». Légende organisée en bas à gauche. Flèche du nord. Source : OMC, CNUCED, FMI.
\end{enumerate}

\textbf{Observations / Interprétation attendues :}
\begin{itemize}
\item La Triade reste le centre de gravité financier et technologique, mais son poids dans le commerce mondial diminue.
\item Les flux Sud-Sud (Chine-Afrique, Chine-Amérique du Sud, Inde-Afrique) s'intensifient.
\item Les PRI (surtout tranche supérieure) sont à la fois \textbf{origines} (IDE chinois, brésiliens, sud-africains) et \textbf{destinations} (IDE occidentaux, chinois) de flux.
\item L'Afrique (majorité PRI inférieurs + PPTE) reste périphérique : exportatrice de matières premières, importatrice de produits manufacturés, réceptrice nette d'IDE (surtout extractifs/infrastructures).
\end{itemize}
\end{experience}

\coursec{Synthèse et TP 5 : cartographier les foyers de la mondialisation}

Nous venons de voir que les pays à revenus intermédiaires se situent dans une position inconfortable : trop avancés pour bénéficier des aides réservées aux plus pauvres, trop fragiles pour rivaliser avec les puissances de la Triade. Cette situation intermédiaire n'a de sens que si l'on replace l'ensemble des pays étudiés dans cette leçon — pays du Golfe, pays d'Afrique, PPTE, pays à revenus intermédiaires — dans la \emph{hiérarchie mondiale des espaces}. C'est précisément l'objet de cette synthèse et du TP 5 : construire une carte du monde qui montre, d'un seul coup d'œil, qui domine la mondialisation, qui émerge, et qui peine à décoller.

\courssub{Bilan : une hiérarchie mondiale inégale}

\begin{definition}[Le système-monde]
Le \cle{système-monde} est un modèle d'analyse (théorisé par le sociologue Immanuel Wallerstein) qui décrit l'économie mondiale comme un ensemble hiérarchisé de trois types d'espaces : le \cle{centre} (ou pôle dominant), la \cle{semi-périphérie} et la \cle{périphérie}, liés entre eux par des échanges inégaux.
\end{definition}

L'intuition est simple : imagine un marché où trois catégories de vendeurs se côtoient. Les premiers fixent les prix, possèdent les capitaux et la technologie ; les seconds commencent à s'enrichir en vendant des produits manufacturés ; les troisièmes ne vendent que des produits bruts (cacao, coton, pétrole brut) dont les prix sont fixés ailleurs. À l'échelle de la planète, cela donne :

\begin{itemize}
\item \textbf{Le centre} : la \cle{Triade}, c'est-à-dire l'Amérique du Nord (États-Unis, Canada), l'Europe occidentale et l'Asie orientale (Japon, et désormais la Chine côtière et la Corée du Sud). Ces espaces concentrent la richesse, les sièges des firmes transnationales, les places financières (New York, Londres, Tokyo) et l'innovation.
\item \textbf{La semi-périphérie} : les \cle{pays émergents} (Chine dans son ensemble, Inde, Brésil, Afrique du Sud, Mexique, Turquie...) et les pays à revenus intermédiaires. Ils s'industrialisent, attirent des investissements, mais restent dépendants du centre pour les technologies et les capitaux.
\item \textbf{La périphérie} : les pays au décollage difficile — la plupart des pays d'Afrique subsaharienne, les \cle{PPTE} (Pays Pauvres Très Endettés), certains pays du Golfe malgré leur rente pétrolière. Ils fournissent des matières premières et captent peu d'investissements.
\end{itemize}

\begin{popfigure}
\begin{center}
\begin{tikzpicture}[scale=1]
% Cercles concentriques : du plus grand au plus petit
\fill[popgreenL] (0,0) circle (4.8);
\fill[poporangeL] (0,0) circle (3.2);
\fill[popblue!70!white] (0,0) circle (1.6);
% Labels
\node[font=\bfseries\small, text=popdark] at (0,0.55) {CENTRE};
\node[font=\scriptsize, text=popdark, align=center] at (0,-0.45) {Triade : Amérique du Nord,\\ Europe occidentale, Asie orientale};
\node[font=\bfseries\small, text=popdark] at (0,2.45) {SEMI-PÉRIPHÉRIE};
\node[font=\scriptsize, text=popdark, align=center] at (0,-2.45) {Pays émergents et à revenus\\ intermédiaires (Chine, Inde, Brésil...)};
\node[font=\bfseries\small, text=popdark] at (0,4.05) {PÉRIPHÉRIE};
\node[font=\scriptsize, text=popdark, align=center] at (0,-4.05) {Pays au décollage difficile : PPTE,\\ majeure partie de l'Afrique};
% Flux
\draw[-{Stealth[length=3mm]}, line width=2.2pt, poppurple] (2.6,1.8) -- (1.15,0.8);
\node[font=\scriptsize, text=poppurple, align=left] at (3.1,2.35) {IDE, technologies,\\ produits manufacturés};
\draw[-{Stealth[length=3mm]}, line width=1.2pt, popgreen!60!black] (-1.15,-0.8) -- (-2.6,-1.8);
\node[font=\scriptsize, text=popgreen!50!black, align=right] at (-3.1,-2.35) {Matières premières\\ (pétrole, cacao, coton...)};
\end{tikzpicture}
\end{center}
\legende{Schéma-bilan de la hiérarchie des espaces dans le système-monde. La flèche violette (épaisse) symbolise les flux dominants (capitaux, technologies) partant du centre ; la flèche verte (fine) symbolise les flux de matières premières venant de la périphérie. L'épaisseur des flèches traduit l'inégalité des échanges.}
\end{popfigure}

\begin{aretenir}[L'essentiel du bilan]
La mondialisation ne profite pas à tous de la même manière. Elle est polarisée par la Triade, intègre progressivement les pays émergents, et marginalise les pays au décollage difficile. La richesse d'un pays du Golfe (rente pétrolière) ne signifie pas un développement achevé : le \emph{décollage} suppose une économie diversifiée, une population formée et des institutions stables.
\end{aretenir}

\courssub{Une insertion inégale dans les flux mondiaux}

Pour comprendre la place des pays pauvres dans la mondialisation, observons deux types de flux : les \cle{IDE} (Investissements Directs à l'Étranger), c'est-à-dire les capitaux investis durablement par des firmes étrangères (usines, filiales, infrastructures), et les \cle{flux de marchandises} (exportations et importations de biens).

\begin{exemplebox}[L'Afrique, marginale dans les flux mondiaux]
L'Afrique ne capte qu'environ \textbf{3 à 4 \%} des flux mondiaux d'IDE (ordre de grandeur donné par la CNUCED ; les chiffres varient selon les années, autour de 40 à 80 milliards de dollars sur un total mondial dépassant souvent \num{1000} milliards). Sa part dans le commerce mondial de marchandises est également faible, de l'ordre de 2 à 3 \%. De plus, ces IDE se concentrent sur quelques pays et quelques secteurs : pétrole (Nigeria, Angola), mines (Afrique du Sud, RDC), et de plus en plus télécommunications et services. Le continent exporte surtout des produits \emph{bruts} dont les prix, fixés sur les bourses de Londres ou de New York, fluctuent fortement : c'est le problème de la \cle{détérioration des termes de l'échange}.
\end{exemplebox}

Le raisonnement à retenir est le suivant : un pays qui vend du cacao brut et importe du chocolat, des machines et des médicaments vend \emph{peu cher} ce qu'il produit et achète \emph{cher} ce dont il a besoin. Sans transformation locale ni capitaux, il reste en bas de la chaîne de valeur mondiale. C'est le cercle vicieux dans lequel sont pris la plupart des PPTE.

\courssub{TP 5 : cartographier les foyers de la mondialisation}

\textbf{Objectif du TP.} Sur un fond de carte mappemonde, tu vas représenter et localiser la Triade, les pays émergents et les pays au décollage difficile, puis figurer par des flèches proportionnelles les principaux flux financiers (IDE) et de marchandises. C'est exactement le type de production attendue au Baccalauréat dans les épreuves de croquis et de cartographie.

\begin{methode}[Réaliser une carte de flux en six étapes]
\begin{enumerate}
\item \textbf{Tracer le fond de carte} : contours simplifiés des continents, tropiques et équateur en pointillés, flèche du nord. Reste schématique : la lisibilité prime sur la précision.
\item \textbf{Localiser les espaces} : colorie la Triade (Amérique du Nord, Europe occidentale, Asie orientale) d'une première couleur, les pays émergents (Chine, Inde, Brésil, Afrique du Sud, Mexique, Turquie...) d'une deuxième, et les pays au décollage difficile (Afrique subsaharienne, PPTE) d'une troisième.
\item \textbf{Choisir les flux à représenter} : retiens 4 à 6 flux majeurs seulement (ex. : IDE États-Unis $\to$ Europe, Europe $\to$ Afrique, marchandises Chine $\to$ États-Unis, pétrole Golfe $\to$ Asie, matières premières Afrique $\to$ Europe et Chine).
\item \textbf{Proportionner les flèches} : l'épaisseur de chaque flèche doit être proportionnelle à la valeur du flux. Définis une échelle (ex. : 1 mm d'épaisseur $=$ 100 milliards de dollars) et respecte-la.
\item \textbf{Ordonner la légende} : présente d'abord les aires (couleurs), puis les flux (flèches), avec leur signification et l'échelle des épaisseurs.
\item \textbf{Titrer la carte} : le titre doit annoncer le sujet, l'espace et la date (ex. : « Les foyers de la mondialisation et les principaux flux mondiaux, vers 2020 »).
\end{enumerate}
\end{methode}

\begin{popfigure}
\includegraphics[width=\linewidth,height=9cm,keepaspectratio]{N15/revenus.png}
\legende{Fond de carte pour la mappemonde « Les foyers de la mondialisation et les principaux flux mondiaux, vers 2020 » (TP 5). Les groupes de revenu de la Banque mondiale en 2023 (carte en anglais ; légende en bas de la carte). Violet foncé : pays à faible revenu (la plus grande partie de l'Afrique subsaharienne, l'Afghanistan, le Yémen) ; rose : revenu intermédiaire inférieur (Inde, Nigeria, Cameroun, Maroc\ldots) ; vert clair : revenu intermédiaire supérieur (Chine, Brésil, Afrique du Sud, Mexique\ldots) ; vert foncé : revenu élevé (Amérique du Nord, Europe, Japon, Australie, Russie, pays du Golfe\ldots) ; hachures : pas de données. Elle aide à situer 1. la Triade (Amérique du Nord, Europe occidentale, Japon) ; 2. les pays émergents (Chine, Inde, Brésil, Afrique du Sud) ; 3. les pays au décollage difficile (Afrique subsaharienne). Les flux d'IDE et de marchandises, d'épaisseur proportionnelle, sont à ajouter. \\ Source : Wikimedia Commons, Our World in Data (données Banque mondiale, 2024), CC BY.}
\end{popfigure}

\begin{attention}[Erreurs fréquentes qui coûtent des points au Bac]
\begin{itemize}
\item \textbf{Oublier la légende ou le titre} : une carte sans légende ordonnée ni titre explicite est une carte non finie ; le correcteur ne peut pas la lire, et tu perds des points même si ton fond est correct.
\item \textbf{Des flèches non proportionnelles} : si tous tes flux ont la même épaisseur, tu ne montres pas l'inégalité des échanges, qui est pourtant l'idée centrale du TP.
\item \textbf{Surcharger la carte} : mieux vaut 5 flux bien choisis et lisibles que 15 flèches illisibles.
\item \textbf{Confondre les catégories} : la Chine est un pays émergent (semi-périphérie), pas un membre historique de la Triade ; le pétrole du Golfe ne suffit pas à classer ces pays parmi les pôles de la mondialisation.
\end{itemize}
\end{attention}

\begin{methode}[Rédiger la conclusion d'une carte ou d'un croquis : décrire, expliquer, nuancer]
\begin{enumerate}
\item \textbf{Décrire} : énonce ce que la carte montre (« La mondialisation est polarisée par trois pôles... les flux les plus épais relient... »). Cite des faits visibles, pas encore des explications.
\item \textbf{Expliquer} : mobilise tes connaissances pour rendre compte de ces faits (concentration des firmes transnationales et des places financières dans la Triade, spécialisation de la périphérie dans les matières premières, poids de la dette des PPTE).
\item \textbf{Nuancer / ouvrir} : montre les limites ou les évolutions (montée en puissance de la Chine, croissance de certains pays africains, initiatives d'annulation de la dette). Termine par une phrase d'ouverture.
\end{enumerate}
\end{methode}

\begin{exoresolu}[Rédiger la conclusion de la carte du TP 5]
Énoncé : à partir de la mappemonde que tu as construite, rédige une conclusion interprétative de 8 à 12 lignes.
\tcblower
\textbf{Solution détaillée (modèle de rédaction).}
\emph{Décrire} : la carte montre une mondialisation fortement polarisée. Trois pôles — l'Amérique du Nord, l'Europe occidentale et l'Asie orientale — concentrent l'essentiel des flux d'IDE et de marchandises : les flèches les plus épaisses relient ces pôles entre eux. Autour d'eux, les pays émergents (Chine, Inde, Brésil, Afrique du Sud) s'intègrent progressivement, tandis que les pays au décollage difficile, notamment en Afrique subsaharienne, ne reçoivent que des flux fins, orientés vers l'exportation de matières premières.
\emph{Expliquer} : cette hiérarchie s'explique par la concentration des firmes transnationales, des capitaux et des technologies dans la Triade, par la spécialisation de la périphérie dans les produits bruts aux prix instables, et par le poids de la dette qui étouffe les PPTE.
\emph{Nuancer} : la carte fige une réalité mouvante : la Chine est devenue un pôle majeur, et plusieurs pays africains connaissent une croissance soutenue. La mondialisation n'exclut donc pas définitivement la périphérie, mais elle l'y maintient tant que celle-ci ne transforme pas ses richesses sur place.
\end{exoresolu}

\courssub{Ouverture : quelles stratégies pour sortir du sous-développement ?}

La carte que tu viens de construire n'est pas une fatalité. Les économistes et les États africains débattent de plusieurs stratégies complémentaires :

\begin{itemize}
\item \textbf{L'industrialisation par transformation locale} : transformer le cacao en chocolat, le coton en textile, le pétrole brut en produits raffinés, afin de capter une plus grande part de la valeur ajoutée. C'est la voie suivie avec succès par les pays émergents d'Asie.
\item \textbf{L'intégration régionale} : regrouper les marchés nationaux trop étroits dans des ensembles régionaux (CEMAC et CEEAC en Afrique centrale, CEDEAO en Afrique de l'Ouest, et désormais la \cle{ZLECAf}, Zone de Libre-Échange Continentale Africaine, dont les échanges ont effectivement démarré en 2021) pour gagner en poids dans les négociations commerciales.
\item \textbf{La bonne gouvernance} : lutte contre la corruption, stabilité politique, transparence dans la gestion des ressources — conditions sans lesquelles ni les IDE ni l'aide internationale ne produisent d'effets durables.
\item \textbf{L'allègement de la dette et l'investissement humain} : l'initiative PPTE a montré que l'annulation de la dette doit s'accompagner d'investissements dans l'éducation et la santé, car le capital humain est le premier moteur du décollage.
\end{itemize}

\begin{savaistu}[Le port de Douala, porte d'entrée de l'Afrique centrale]
Le port autonome de Douala traite l'essentiel (environ 90 \%, ordre de grandeur couramment cité) du commerce extérieur du Cameroun : exportations de pétrole, de bois, de cacao, de café, et importations de biens d'équipement et de produits manufacturés. Il dessert aussi les pays enclavés voisins : une grande partie des marchandises destinées au \textbf{Tchad} et à la \textbf{Centrafrique} transite par Douala puis par les corridors routiers Douala--N'Djamena et Douala--Bangui. C'est un exemple parfait de la dépendance logistique de la périphérie : toute l'Afrique centrale s'insère dans la mondialisation par cet entonnoir portuaire quasi unique.
\end{savaistu}

\begin{exercice}[Pour t'entraîner]
Sur un fond de mappemonde vierge fourni par ton professeur : (1) localise et colorie la Triade, quatre pays émergents et l'ensemble des pays au décollage difficile ; (2) trace cinq flux majeurs (deux d'IDE, trois de marchandises) avec des épaisseurs proportionnelles ; (3) rédige la légende ordonnée, le titre, puis une conclusion de dix lignes suivant la méthode « décrire, expliquer, nuancer ».
\end{exercice}

\begin{corrige}[Exercice]
\textbf{(1) Aires} : Triade en bleu (Amérique du Nord, Europe occidentale, Japon/Asie orientale) ; émergents en orange (Chine, Inde, Brésil, Afrique du Sud) ; décollage difficile en vert (Afrique subsaharienne, PPTE, Golfe). \textbf{(2) Flux} : IDE États-Unis $\leftrightarrow$ Europe (flèche très épaisse), IDE Europe $\to$ Afrique (fine) ; marchandises Chine $\to$ Amérique du Nord et Europe (épaisse), pétrole Golfe $\to$ Asie (moyenne), matières premières Afrique $\to$ Europe et Chine (fine). \textbf{(3) Légende} : aires numérotées 1 à 3, flux 4 et 5 avec échelle d'épaisseur ; titre daté ; conclusion en trois temps conforme au modèle de l'exercice résolu ci-dessus.
\end{corrige}

\begin{aretenir}[Ce qu'il faut retenir de la leçon]
Le système-monde est hiérarchisé : un centre (la Triade) qui domine, une semi-périphérie (émergents et revenus intermédiaires) qui rattrape, une périphérie (PPTE, majeure partie de l'Afrique) qui peine à décoller malgré parfois une rente pétrolière (Golfe). L'Afrique ne capte que 3 à 4 \% des IDE mondiaux et vend des produits bruts. Savoir cartographier cette hiérarchie — fond de carte, flèches proportionnelles, légende ordonnée, titre daté, conclusion en trois temps — est une compétence directement évaluée au Baccalauréat.
\end{aretenir}

\newpage

\coursec{Activité d'intégration}

\coursec{Leçon 15 : Les pays au décollage économique difficile}

\courssub{TP 5 : Les foyers de la mondialisation — Cartographier les pôles et les flux}

\courssub{Objectifs de l'activité}

À l'issue de ce travail pratique, tu seras capable de :
\begin{itemize}
\item \cle{classer} des pays en catégories (Triade, pays émergents, pays à revenus intermédiaires, PPTE) à partir d'indicateurs économiques et sociaux (PIB par habitant, IDH, endettement) ;
\item \cle{reporter} sur un fond de carte mappemonde vierge ces catégories à l'aide d'un code couleur cohérent ;
\item \cle{tracer} les principaux flux de marchandises et d'investissements directs à l'étranger (IDE) avec des flèches d'épaisseur proportionnelle ;
\item \cle{rédiger} une légende complète et une conclusion argumentée montrant la place marginale des pays au décollage difficile dans la mondialisation ;
\item \cle{présenter} oralement ton croquis et défendre tes choix de représentation.
\end{itemize}

\courssub{Situation}

\textbf{Au lycée de Nkongsamba, un défi cartographique.} Dans le cadre de la semaine de la géographie, le proviseur du lycée de Nkongsamba (région du Littoral) lance un concours entre les classes de Terminale : chaque groupe doit produire une mappemonde intitulée « \emph{Les foyers et les flux de la mondialisation : un monde inégal} », qui sera affichée dans la salle polyvalente. Ton groupe reçoit le dossier documentaire suivant.

\textbf{Document 1. Les IDE dans le monde (ordres de grandeur, années récentes).} Les flux mondiaux d'investissements directs à l'étranger représentent environ \num{1300} à \num{1500} milliards de dollars par an. Les États-Unis, la Chine (avec Hong Kong), Singapour et les pays d'Europe de l'Ouest captent l'essentiel de ces flux. L'Afrique dans son ensemble ne reçoit qu'environ 3 à 5 \% des IDE mondiaux, concentrés sur quelques pays (Égypte, Afrique du Sud, Nigeria, Maroc) et sur les secteurs pétrolier et minier.

\textbf{Document 2. Principaux exportateurs de marchandises (ordres de grandeur).} La Chine est le premier exportateur mondial de marchandises (plus de \num{3000} milliards de dollars), suivie des États-Unis et de l'Allemagne. Les pays du Golfe (Arabie Saoudite, Émirats arabes unis, Qatar, Koweït) figurent parmi les grands exportateurs grâce aux hydrocarbures : le pétrole du Golfe s'écoule massivement vers l'Asie (Chine, Japon, Inde, Corée du Sud). L'Afrique exporte surtout des matières premières brutes (pétrole du Nigeria et de l'Angola, cacao de Côte d'Ivoire et du Ghana, cuivre de Zambie et de RDC, pétrole et bois du Cameroun) vers l'Union européenne et la Chine.

\textbf{Document 3. Indicateurs de dix pays (valeurs approchées, à titre pédagogique).}

\begin{center}
\begin{tabularx}{\linewidth}{|l|X|X|X|}
\hline
\rowcolor{popblueL} \textbf{Pays} & \textbf{PIB/hab. (USD)} & \textbf{IDH} & \textbf{Situation} \\
\hline
États-Unis & \num{70000} & 0,92 & Cœur de la Triade \\
\hline
Allemagne & \num{46000} & 0,94 & Triade (UE) \\
\hline
Japon & \num{40000} & 0,92 & Triade (Asie) \\
\hline
Chine & \num{12000} & 0,77 & Émergent, 1$^{\text{er}}$ exportateur \\
\hline
Brésil & \num{9000} & 0,75 & Émergent (BRICS) \\
\hline
Émirats arabes unis & \num{43000} & 0,91 & Rente pétrolière, Golfe \\
\hline
Gabon & \num{9000} & 0,71 & Revenu intermédiaire (pétrole) \\
\hline
Cameroun & \num{1600} & 0,58 & Revenu intermédiaire inférieur \\
\hline
Tchad & \num{700} & 0,39 & PPTE, pétrole récent \\
\hline
Haïti & \num{1700} & 0,54 & PPTE \\
\hline
\end{tabularx}
\end{center}

\textbf{Document 4. Rappel.} Les PPTE (pays pauvres très endettés), une quarantaine dans le monde dont plus de 30 en Afrique subsaharienne, cumulent un faible revenu, un IDH faible et un endettement extérieur lourd ; ils bénéficient d'initiatives d'allègement de la dette (initiative PPTE de 1996, renforcée en 1999, initiative IADM de 2005).

\courssub{Consignes}

Par groupes de quatre, et en t'appuyant sur les quatre documents :

\begin{enumerate}
\item \textbf{Classer.} Répartis les dix pays du document 3 en quatre catégories : Triade, pays émergents, pays à revenus intermédiaires (PRI), PPTE. Justifie chaque classement en citant au moins deux indicateurs. Que remarques-tu pour les Émirats arabes unis ?
\item \textbf{Reporter.} Sur le fond de carte mappemonde vierge fourni, colorie les pays (ou les ensembles continentaux) selon un code couleur que tu choisiras : une couleur par catégorie. Ajoute le Cameroun et ses voisins.
\item \textbf{Tracer les flux.} Représente par des flèches d'épaisseur proportionnelle : (a) le flux pétrolier du Golfe vers l'Asie orientale ; (b) les flux de matières premières de l'Afrique vers l'Europe et la Chine ; (c) les grands flux d'IDE entre les pôles de la Triade et vers la Chine. Indique le sens du flux par la pointe de la flèche.
\item \textbf{Légender.} Rédige une légende complète et organisée (d'abord les catégories de pays, puis les flux), avec un titre de carte et l'orientation.
\item \textbf{Conclure.} Rédige une conclusion d'une dizaine de lignes montrant que la mondialisation est polarisée par quelques foyers et que les pays au décollage difficile (PPTE, nombre de pays africains) y occupent une place marginale, malgré leurs ressources.
\item \textbf{Restituer.} Présente ton croquis en cinq minutes ; la classe et l'enseignant corrigent collectivement au tableau.
\end{enumerate}

\courssub{Ressources à mobiliser}

\begin{aretenir}[Ce qu'il faut avoir en tête avant de commencer]
\begin{itemize}
\item La \cle{Triade} : trois pôles qui concentrent l'essentiel des richesses, des échanges et des IDE — Amérique du Nord (États-Unis, Canada), Europe de l'Ouest (Union européenne), Asie orientale (Japon, puis Chine littorale, Corée du Sud).
\item Les \cle{pays émergents} : Chine, Inde, Brésil, Russie, Afrique du Sud (BRICS), Mexique, Turquie... : industrialisation rapide, insertion croissante dans les échanges mondiaux.
\item Les \cle{pays du Golfe} : décollage fondé sur la rente pétrolière et gazière, PIB par habitant élevé mais économie peu diversifiée.
\item Les \cle{PRI} : PIB par habitant situé entre environ \num{1000} et \num{12000} dollars (seuils de la Banque mondiale, à titre indicatif) ; le Cameroun est un pays à revenu intermédiaire inférieur.
\item Les \cle{PPTE} : pauvreté, IDH faible, dette extérieure écrasante ; marginalisés dans les flux mondiaux.
\item Savoir-faire cartographiques : choix d'un code couleur limité (4 à 5 teintes), flèches proportionnelles, légende hiérarchisée, titre, orientation.
\end{itemize}
\end{aretenir}

\courssub{Démarche et corrigé commenté}

\begin{exoresolu}[TP 5 : La mappemonde des foyers et flux de la mondialisation]
Énoncé : classer les dix pays, construire la mappemonde (catégories + flux), rédiger légende et conclusion.
\tcblower
\textbf{Étape 1 — Le classement justifié.}

\emph{Triade} : États-Unis, Allemagne, Japon. Justification : PIB par habitant très élevé (\num{40000} à \num{70000} dollars) et IDH supérieur à 0,9 ; ce sont les pôles qui émettent et reçoivent l'essentiel des IDE.

\emph{Pays émergents} : Chine, Brésil. Justification : PIB par habitant moyen (\num{9000} à \num{12000} dollars), IDH autour de 0,75-0,77, mais poids considérable dans le commerce mondial (la Chine est le premier exportateur mondial) et forte croissance.

\emph{Pays à revenus intermédiaires} : Gabon (environ \num{9000} dollars grâce au pétrole, mais IDH de 0,71 seulement), Cameroun (environ \num{1600} dollars, IDH 0,58 : revenu intermédiaire \emph{inférieur}).

\emph{PPTE} : Tchad (environ \num{700} dollars, IDH 0,39, endettement lourd malgré le pétrole exploité depuis 2003), Haïti (IDH 0,54, pauvreté et dette).

\emph{Cas particulier} : les Émirats arabes unis ont un PIB par habitant et un IDH de niveau « Triade » (\num{43000} dollars, 0,91), mais leur richesse repose presque entièrement sur la \cle{rente pétrolière} et non sur une économie diversifiée : on les classe donc à part, dans les « pays du Golfe au décollage fondé sur la rente ». C'est la subtilité attendue : les indicateurs seuls ne suffisent pas, il faut connaître la nature de l'économie.

\textbf{Étape 2 — Le report sur la carte.} Code couleur proposé : bleu foncé pour la Triade, orange pour les émergents, vert pour les PRI, rouge pour les PPTE, violet pour les pays du Golfe. On colorie aussi les ensembles : Europe de l'Ouest et Amérique du Nord en bleu, la majeure partie de l'Afrique subsaharienne en rouge (PPTE) avec quelques taches vertes (Cameroun, Côte d'Ivoire, Sénégal...), le Maghreb et l'Afrique du Sud en vert ou orange selon le cas.

\textbf{Étape 3 — Les flux.} Trois familles de flèches : une flèche épaisse du Golfe persique vers le Japon, la Chine et l'Inde (pétrole) ; des flèches moyennes de l'Afrique de l'Ouest et centrale vers l'Europe et la Chine (matières premières) ; des flèches très épaisses formant un triangle entre Amérique du Nord, Europe et Asie orientale (IDE intra-Triade), avec une branche vers la Chine. Aucune flèche d'IDE significative vers les PPTE : c'est précisément ce vide qu'il faut faire apparaître.

\textbf{Étape 4 — La légende.} Titre : « Les foyers et les flux de la mondialisation ». Légende en deux volets : d'abord les cinq couleurs de pays, ensuite les trois types de flèches avec leur signification. Flèche du nord et échelle indicative.

\textbf{Étape 5 — Conclusion rédigée (modèle).} « La carte montre un monde profondément inégal. La richesse, les échanges et les investissements se concentrent sur quelques foyers : la Triade (Amérique du Nord, Europe de l'Ouest, Asie orientale), rejointe par les pays émergents, la Chine en tête. Les pays du Golfe, riches de leur rente pétrolière, alimentent l'Asie en énergie sans diversifier leur économie. À l'opposé, les PPTE et une grande partie de l'Afrique, dont le Cameroun voisin des PPTE tchadien et centrafricain, restent à la marge : ils fournissent des matières premières brutes à bas prix mais ne reçoivent que 3 à 5 \% des IDE mondiaux. Leur décollage économique, freiné par la dette, la pauvreté et la faiblesse des industries, illustre une mondialisation qui intègre inégalement les territoires. »
\end{exoresolu}

\begin{popfigure}
\includegraphics[width=\linewidth,height=9cm,keepaspectratio]{N15/revenus.png}
\legende{Fond de carte pour le croquis attendu. Les groupes de revenu de la Banque mondiale en 2023 (carte en anglais ; légende en bas de la carte). Violet foncé : pays à faible revenu (la plus grande partie de l'Afrique subsaharienne, l'Afghanistan, le Yémen) ; rose : revenu intermédiaire inférieur (Inde, Nigeria, Cameroun, Maroc\ldots) ; vert clair : revenu intermédiaire supérieur (Chine, Brésil, Afrique du Sud, Mexique\ldots) ; vert foncé : revenu élevé (Amérique du Nord, Europe, Japon, Australie, Russie, pays du Golfe\ldots) ; hachures : pas de données. L'élève s'en sert pour distinguer la Triade et les pays riches (vert foncé), les pays émergents et intermédiaires (verts clairs et roses) et les pays les plus pauvres (violet), puis y ajoute lui-même les flux (IDE, pétrole du Golfe, matières premières d'Afrique). \\ Source : Wikimedia Commons, Our World in Data (données Banque mondiale, 2024), CC BY.}
\end{popfigure}

\begin{attention}[Pièges fréquents dans ce TP]
\begin{itemize}
\item Classer les Émirats dans la Triade parce que leur PIB par habitant est élevé : erreur, leur richesse est une \cle{rente}, pas le signe d'une économie diversifiée intégrée comme celle des pôles.
\item Oublier la pointe des flèches (sens du flux) ou tracer des flèches d'épaisseur identique alors que l'énoncé exige des flèches \cle{proportionnelles}.
\item Multiplier les couleurs (plus de cinq) : la carte devient illisible.
\item Confondre « pays émergents » et « pays à revenus intermédiaires » : tous les émergents sont des PRI, mais tous les PRI ne sont pas émergents (le Cameroun n'est pas encore un pays émergent).
\item Rédiger une légende sans titre de carte ni hiérarchie (pays d'abord, flux ensuite).
\item Ne pas localiser le Cameroun et ses voisins (Tchad, RCA, Nigeria, Gabon, Congo, Guinée équatoriale) alors que la consigne le demande explicitement.
\end{itemize}
\end{attention}

\courssub{Bilan de l'activité}

Ce TP t'a permis de passer du savoir au savoir-faire : en croisant PIB par habitant, IDH et nature des économies, tu as classé des pays en quatre catégories et découvert un cas limite (les Émirats, riches mais rentiers). La mappemonde construite met en évidence trois enseignements majeurs de la leçon : la \cle{polarisation} de la mondialisation autour de la Triade et des émergents ; le \cle{décollage fondé sur la rente} des pays du Golfe, intégrés par leurs hydrocarbures mais vulnérables ; et la \cle{marginalisation} des PPTE et de nombreux pays africains, fournisseurs de matières premières brutes mais exclus des grands flux d'investissement. Tu as enfin mobilisé les gestes essentiels du cartographe : choisir un code couleur, proportionner les flèches, hiérarchiser une légende et tirer d'une carte une conclusion géographique argumentée — compétences directement exigibles au Baccalauréat.

\newpage

\coursec{Méthodes et exercices résolus}

\courssub{Méthode 1 : Observer, localiser et décrire un phénomène géographique (pays en difficulté)}

\begin{methode}[Décrire la situation d'un pays au décollage difficile]
\begin{enumerate}
\item \textbf{Localiser} précisément : continent, sous-région, voisins, position maritime ou enclavement (ex. : le Tchad, pays enclavé d'Afrique centrale ; le Qatar, presqu'île du golfe Arabo-Persique).
\item \textbf{Identifier la catégorie} : pays du Golfe rentier, pays d'Afrique au décollage freiné, PPTE (Pays Pauvres Très Endettés), pays à revenus intermédiaires (PRI). Un pays peut cumuler plusieurs statuts (ex. : le Cameroun est à la fois PPTE et PRI).
\item \textbf{Décrire avec des chiffres} : IDH, PIB par habitant, part de l'agriculture dans l'emploi, taux d'endettement, espérance de vie. Toujours donner un ordre de grandeur et sa source approximative (PNUD, Banque mondiale).
\item \textbf{Hiérarchiser} : distinguer les atouts (ressources, position, population) des contraintes (dette, enclavement, dépendance aux matières premières, instabilité).
\item \textbf{Conclure} par une phrase bilan qui répond explicitement à la question posée.
\end{enumerate}
\textbf{Réflexe} : une description géographique va toujours du général (localisation) vers le particulier (indicateurs), jamais l'inverse.
\end{methode}

\begin{definition}[Les quatre catégories à connaître]
\begin{itemize}
\item \cle{Pays du Golfe} : monarchies riveraines du golfe Arabo-Persique (Arabie saoudite, Koweït, Bahreïn, Qatar, Émirats arabes unis, Oman), enrichies par la rente pétrolière et gazière.
\item \cle{PPTE} (Pays Pauvres Très Endettés) : catégorie définie par la Banque mondiale et le FMI (1996) selon le poids insoutenable de la dette extérieure ; environ 37 pays concernés, dont une trentaine en Afrique (Cameroun, Tchad, Niger...), mais aussi Haïti ou la Bolivie.
\item \cle{Pays à revenus intermédiaires} (PRI) : catégorie de la Banque mondiale fondée sur le revenu national brut par habitant (ex. : Cameroun, Côte d'Ivoire, Malaisie, Maroc).
\item \cle{PMA} (Pays les Moins Avancés) : catégorie de l'ONU combinant pauvreté, faiblesse des ressources humaines et vulnérabilité économique.
\end{itemize}
\end{definition}

\begin{exoresolu}[Décrire la situation du Tchad]
\textbf{Énoncé} : À partir de vos connaissances, décrivez la situation géographique et économique du Tchad et montrez qu'il s'agit d'un pays au décollage économique difficile. (10 lignes maximum)
\tcblower
\textbf{Solution détaillée.}

\emph{Étape 1 — Localisation.} Le Tchad est un pays d'Afrique centrale, enclavé (sans accès à la mer), situé au sud de la Libye, à l'est du Niger et du Nigeria, au nord du Cameroun et de la Centrafrique. Sa capitale est N'Djamena.

\emph{Étape 2 — Catégorie.} Le Tchad appartient au groupe des \cle{PPTE} (Pays Pauvres Très Endettés) et figure parmi les Pays les Moins Avancés (PMA).

\emph{Étape 3 — Indicateurs.} Son IDH est parmi les plus faibles du monde (autour de 0,4, ordre de grandeur PNUD) ; le PIB par habitant est de l'ordre de quelques centaines de dollars ; l'agriculture et l'élevage occupent la grande majorité de la population active ; l'espérance de vie est voisine de 55 ans.

\emph{Étape 4 — Atouts et contraintes.} Atout : le pétrole exploité depuis 2003 à Doba, exporté par l'oléoduc Tchad--Cameroun (environ \SI{1070}{km}) jusqu'au terminal de Kribi. Contraintes : enclavement qui renchérit les transports, dépendance à une seule ressource, sécheresses récurrentes, insécurité régionale (région du lac Tchad, bande sahélienne), poids de la dette.

\emph{Étape 5 — Conclusion.} Malgré la rente pétrolière, le Tchad reste un pays au décollage difficile : la combinaison de l'enclavement, de la pauvreté et de la dépendance pétrolière bloque son insertion bénéfique dans la mondialisation.
\end{exoresolu}

\courssub{Méthode 2 : Lire et interpréter un tableau statistique (IDH, PIB/habitant, dette)}

\begin{methode}[Exploiter un tableau de données comparatives]
\begin{enumerate}
\item \textbf{Lire le titre, l'unité, la source et la date} avant tout commentaire : un PIB en dollars courants ne se compare pas à un PIB en parité de pouvoir d'achat (PPA).
\item \textbf{Repérer les valeurs extrêmes} (maximum, minimum) et les citer avec leurs chiffres exacts.
\item \textbf{Calculer les écarts} : rapport (max/min) et taux de variation :
\[ \text{Taux de variation (\%)} = \frac{V_f - V_i}{V_i} \times 100 \]
\item \textbf{Classer} les pays en groupes cohérents (ex. : Golfe / Afrique / PRI) et vérifier si la classification correspond aux catégories du cours.
\item \textbf{Interpréter} : relier chaque chiffre à un mécanisme géographique (rente pétrolière, endettement, diversification industrielle).
\end{enumerate}
\textbf{Formule à retenir} : un écart s'exprime en fois (rapport) ou en points de pourcentage, jamais en confondant les deux.
\end{methode}

\begin{exoresolu}[Comparer des PIB par habitant]
\textbf{Énoncé} : On donne les PIB par habitant (PPA, dollars courants, ordres de grandeur Banque mondiale) : Qatar : 90\,000 ; Malaisie : 30\,000 ; Cameroun : 4\,000 ; Tchad : 1\,600. 1) Calculez l'écart entre le Qatar et le Tchad. 2) Calculez le taux de variation entre la Malaisie et le Cameroun (de la Malaisie vers le Cameroun). 3) Que révèlent ces écarts ?
\tcblower
\textbf{Solution détaillée.}

\emph{1) Écart Qatar--Tchad.} Rapport :
\[ \frac{90\,000}{1\,600} = 56,25 \]
Le PIB par habitant du Qatar est donc environ \textbf{56 fois} celui du Tchad. En valeur absolue : $90\,000 - 1\,600 = 88\,400$ dollars d'écart.

\emph{2) Taux de variation Malaisie $\to$ Cameroun.}
\[ \frac{4\,000 - 30\,000}{30\,000} \times 100 = \frac{-26\,000}{30\,000} \times 100 \approx -86,7\,\% \]
Le PIB par habitant camerounais est inférieur d'environ \textbf{87 \%} à celui de la Malaisie (autre formulation correcte : il ne représente que $4\,000/30\,000 \approx 13\,\%$ de celui de la Malaisie).

\emph{3) Interprétation.} Ces écarts révèlent la \cle{hiérarchie des territoires} dans la mondialisation : le Qatar illustre l'enrichissement rapide par la rente pétrolière et gazière (pays du Golfe) ; la Malaisie, pays à revenus intermédiaires, a réussi une diversification industrielle (électronique, transformation de l'huile de palme) ; le Cameroun, bien que PRI, reste freiné par la dépendance aux matières premières brutes ; le Tchad, PPTE enclavé, cumule pauvreté et dette. La mondialisation produit donc des intégrations très inégales.

\emph{Conclusion.} Les chiffres confirment que le « décollage » n'est pas un processus uniforme : entre les rentiers du Golfe et les PPTE africains, l'écart de richesse par habitant dépasse un facteur 50.
\end{exoresolu}

\courssub{Méthode 3 : Construire un diagramme en barres à partir d'un tableau}

\begin{methode}[Tracer un diagramme en barres lisible]
\begin{enumerate}
\item \textbf{Choisir l'échelle} : l'axe vertical doit partir de 0 et contenir la valeur maximale avec une marge ; graduations régulières.
\item \textbf{Ordonner les barres} (croissant ou décroissant) sauf si l'ordre du tableau est imposé.
\item \textbf{Tracer} des barres de largeur égale, espacées régulièrement, avec les valeurs inscrites au sommet.
\item \textbf{Légender} : titre complet (quoi, où, quand), nom et unité de chaque axe, source des données.
\item \textbf{Exploiter} : le graphique sert à appuyer le commentaire, pas à le remplacer.
\end{enumerate}
\textbf{Réflexe} : une barre sans valeur inscrite ni axe gradué est sanctionnée au Bac.
\end{methode}

\begin{exoresolu}[Représenter l'IDH de quatre pays]
\textbf{Énoncé} : Représentez par un diagramme en barres les IDH suivants (PNUD, ordres de grandeur) : Émirats arabes unis : 0,91 ; Malaisie : 0,80 ; Cameroun : 0,58 ; Tchad : 0,39. Puis dégagez deux enseignements.
\tcblower
\textbf{Solution détaillée.}

\emph{Étape 1 — Échelle.} Valeur maximale : 0,91. On gradue l'axe vertical de 0 à 1 avec un pas de 0,2.

\emph{Étape 2 — Tracé} (barres ordonnées de façon décroissante) :

\begin{popfigure}
\begin{tikzpicture}[x=1.9cm,y=5.5cm]
\draw[->] (0,0) -- (4.6,0) node[below left,font=\small]{Pays};
\draw[->] (0,0) -- (0,1.05) node[above,font=\small]{IDH};
\foreach \y in {0,0.2,0.4,0.6,0.8,1.0}{
  \draw[popline] (0,\y) -- (4.4,\y);
  \node[left,font=\scriptsize] at (0,\y) {\pgfmathprintnumber[fixed,precision=1]{\y}};
}
\fill[popblue]   (0.3,0) rectangle (1.0,0.91); \node[above,font=\scriptsize] at (0.65,0.91) {0,91}; \node[below,font=\scriptsize,align=center] at (0.65,0) {EAU};
\fill[popgreen]  (1.3,0) rectangle (2.0,0.80); \node[above,font=\scriptsize] at (1.65,0.80) {0,80}; \node[below,font=\scriptsize,align=center] at (1.65,0) {Malaisie};
\fill[poporange] (2.3,0) rectangle (3.0,0.58); \node[above,font=\scriptsize] at (2.65,0.58) {0,58}; \node[below,font=\scriptsize,align=center] at (2.65,0) {Cameroun};
\fill[poppink]   (3.3,0) rectangle (4.0,0.39); \node[above,font=\scriptsize] at (3.65,0.39) {0,39}; \node[below,font=\scriptsize,align=center] at (3.65,0) {Tchad};
\end{tikzpicture}
\legende{IDH de quatre pays au décollage contrasté (PNUD, ordres de grandeur). Échelle : 1 cm $\approx$ 0,18 point d'IDH.}
\end{popfigure}

\emph{Étape 3 — Enseignements.}
\begin{itemize}
\item Premier enseignement : le développement humain suit la richesse mais avec des nuances : les Émirats (IDH très élevé, $> 0,9$) ont transformé leur rente pétrolière en infrastructures, santé et éducation.
\item Deuxième enseignement : l'écart Cameroun--Tchad ($0,58 - 0,39 = 0,19$ point) montre que même parmi les pays africains, les situations sont contrastées : le Cameroun, pays à revenus intermédiaires, devance nettement le Tchad, PPTE enclavé.
\end{itemize}

\emph{Conclusion.} Le diagramme visualise la gradation des niveaux de développement : pays du Golfe $>$ PRI asiatique $>$ PRI africain $>$ PPTE.
\end{exoresolu}

\courssub{Méthode 4 : Construire et légender un croquis (pays du Golfe et rente pétrolière)}

\begin{methode}[Réaliser un croquis de synthèse géographique]
\begin{enumerate}
\item \textbf{Tracer le fond} : contours simplifiés mais reconnaissables (péninsule arabique, golfe Arabo-Persique), à main levée ou au crayon, sans détail inutile.
\item \textbf{Placer l'essentiel d'abord} : pays, capitales, champs pétroliers, ports et oléoducs.
\item \textbf{Utiliser des figurés simples et constants} : surface colorée = pays producteur ; point noir = ville ; étoile = capitale ; flèche = flux d'exportation ; trait plein = oléoduc.
\item \textbf{Rédiger la légende} en dessous, numérotée ou à puces, dans l'ordre d'apparition des figurés : tout figuré sur la carte doit figurer dans la légende, et réciproquement.
\item \textbf{Ajouter titre, orientation (flèche du nord) et échelle indicative.}
\end{enumerate}
\textbf{Réflexe} : la légende est un texte rédigé, pas une simple liste de symboles : « les flèches orange représentent les exportations de pétrole vers... ».
\end{methode}

\begin{exoresolu}[Croquis : la rente pétrolière du golfe Arabo-Persique]
\textbf{Énoncé} : Construisez un croquis intitulé « Le golfe Arabo-Persique, pôle pétrolier mondial » montrant les pays producteurs, les principales voies d'exportation et le goulot stratégique d'Ormuz.
\tcblower
\textbf{Solution détaillée.}

\emph{Étape 1 — Organisation de l'espace.} Le golfe est une mer bordée à l'ouest et au sud par la péninsule arabique (Arabie saoudite, Koweït, Bahreïn, Qatar, Émirats arabes unis, Oman) et au nord-est par l'Irak et l'Iran. La seule sortie maritime est le \cle{détroit d'Ormuz} (environ \SI{50}{km} de large au plus étroit), par lequel transite environ un cinquième de la consommation mondiale de pétrole (ordre de grandeur).

\emph{Étape 2 — Croquis.}

\begin{popfigure}
\includegraphics[width=\linewidth,height=9cm,keepaspectratio]{N15/hormuz.png}
\legende{Le détroit d'Ormuz, passage obligé des exportations d'hydrocarbures du golfe Persique (carte en anglais). Les flèches violettes montrent les deux couloirs de navigation (un par sens) qui relient le golfe Persique au golfe d'Oman, entre l'Iran au nord et la péninsule d'Oman (Musandam) et les Émirats arabes unis au sud ; Dubaï, Sharjah et Bandar Abbas sont localisés. La carte ne montre pas les flux de pétrole : ils sont à ajouter au croquis (flèches d'épaisseur proportionnelle vers l'Asie, l'Europe et l'Amérique du Nord). Échelle : voir la barre graduée en bas à droite. \\ Source : Wikimedia Commons, Goran tek-en, CC BY-SA 4.0.}
\end{popfigure}

\emph{Étape 3 — Analyse.} Le croquis montre la \cle{concentration spatiale} de la rente : quelques États de petite taille (Qatar, Koweït, Émirats) contrôlent d'immenses réserves, et leur richesse dépend d'un unique passage maritime, Ormuz, ce qui explique la valeur stratégique mondiale du golfe.

\emph{Conclusion.} Le décollage des pays du Golfe repose sur la rente pétrolière et gazière exportée par voie maritime ; leur vulnérabilité tient à la dépendance à cette ressource non renouvelable et au goulot d'Ormuz.
\end{exoresolu}

\courssub{Méthode 5 : Commenter un document (carte, photographie, statistiques) au Bac}

\begin{methode}[Rédiger un commentaire de document en trois parties]
\begin{enumerate}
\item \textbf{Présentation} (2-3 lignes) : nature du document, titre, source, date, ce qu'il représente.
\item \textbf{Description organisée} : décrire ce qu'on voit, du plus général au plus détaillé, en citant les éléments du document (chiffres, lieux, figurés) — sans encore expliquer.
\item \textbf{Interprétation} : mobiliser le cours pour expliquer (mécanismes : rente, dette, émergence, marginalisation) et répondre au problème posé.
\item \textbf{Conclusion brève} : une phrase bilan + éventuellement une ouverture.
\end{enumerate}
\textbf{Réflexe} : interdiction de paraphraser le document ; chaque affirmation d'interprétation doit s'appuyer sur un élément précis du document (« on observe que... ce qui s'explique par... »).
\end{methode}

\begin{exoresolu}[Commentaire d'une photographie de Dubaï]
\textbf{Énoncé} : Commentez une photographie du front de mer de Dubaï (Émirats arabes unis) montrant des gratte-ciel (dont la tour Burj Khalifa, 828 m), un port de plaisance et des chantiers, avec en arrière-plan le désert. Problème : comment la rente pétrolière transforme-t-elle les paysages des pays du Golfe ?
\tcblower
\textbf{Solution détaillée.}

\emph{1) Présentation.} Il s'agit d'une photographie aérienne récente du front de mer de Dubaï, l'un des sept émirats des Émirats arabes unis, montrant un paysage urbain ultramoderne édifié au bord du golfe Arabo-Persique.

\emph{2) Description.} Au premier plan, une marina avec des yachts ; au centre, une concentration exceptionnelle de gratte-ciel de verre et d'acier, dominée par la Burj Khalifa (828 m, plus haute tour du monde) ; à l'arrière-plan, le désert aride rappelle le milieu naturel hostile. Des grues témoignent de chantiers permanents.

\emph{3) Interprétation.} Ce paysage s'explique par la \cle{rente pétrolière et gazière} : les hydrocarbures ont fourni des capitaux considérables, investis dans l'immobilier de prestige, les infrastructures (port de Jebel Ali, premier port conteneur du Moyen-Orient ; aéroport international parmi les plus fréquentés du monde) et les services (finance, tourisme, commerce). Dubaï illustre la stratégie de \cle{diversification} des économies du Golfe : anticiper l'après-pétrole en devenant une place tertiaire mondiale, à l'image de Singapour. Mais ce développement repose sur une main-d'œuvre immigrée très nombreuse (les étrangers forment la grande majorité de la population des Émirats) et sur une consommation d'eau et d'énergie considérable (dessalement de l'eau de mer).

\emph{4) Conclusion.} La rente a métamorphosé un port de pêche de perles en métropole mondiale en un demi-siècle ; ce modèle spectaculaire reste dépendant des cours des hydrocarbures et pose la question de sa durabilité.
\end{exoresolu}

\courssub{Méthode 6 : TP 5 — Cartographier les foyers de la mondialisation sur une mappemonde}

\begin{methode}[Réaliser la carte « Triade, émergents et flux » (TP 5)]
\begin{enumerate}
\item \textbf{Reporter le fond} : contours simplifiés des continents sur la mappemonde vierge, avec l'équateur et les principaux parallèles comme repères.
\item \textbf{Colorier la Triade} : Amérique du Nord (États-Unis, Canada), Union européenne, Japon — une même couleur, car ce sont les trois pôles historiques.
\item \textbf{Colorier les émergents} d'une seconde couleur : Chine, Inde, Brésil, Russie, Afrique du Sud (BRICS), Mexique, Turquie, Malaisie...
\item \textbf{Marquer les pays au décollage difficile} (PPTE, Afrique subsaharienne) d'une troisième couleur ou d'un figuré.
\item \textbf{Tracer les flux} : flèches d'épaisseur proportionnelle — flux financiers (IDE) entre les pôles de la Triade et vers la Chine ; flux de marchandises (conteneurs) Asie $\to$ Europe et Asie $\to$ Amérique du Nord ; flux de matières premières Afrique/Moyen-Orient $\to$ pôles.
\item \textbf{Légender} chaque figuré, ajouter titre, orientation et source.
\end{enumerate}
\textbf{Réflexe} : les flux les plus épais relient les pôles entre eux — c'est la preuve visuelle que la mondialisation profite d'abord aux territoires déjà intégrés.
\end{methode}

\begin{exoresolu}[TP 5 : la mappemonde des foyers de la mondialisation]
\textbf{Énoncé} : Sur le fond de mappemonde fourni, représentez et localisez la Triade, les pays émergents et les principaux flux financiers et de marchandises. Rédigez la légende et tirez une conclusion de quatre lignes.
\tcblower
\textbf{Solution détaillée.}

\emph{Étape 1 — Schéma directeur.} Avant de tracer, on organise mentalement : trois pôles (Triade), un arc émergent (Asie orientale, Amérique du Sud, Afrique du Sud), une marge (Afrique subsaharienne hors Afrique du Sud, PPTE).

\emph{Étape 2 — Carte (version schématique).}

\begin{popfigure}
\includegraphics[width=\linewidth,height=9cm,keepaspectratio]{N15/revenus.png}
\legende{Fond de carte pour la mappemonde des foyers de la mondialisation (TP 5). Les groupes de revenu de la Banque mondiale en 2023 (carte en anglais ; légende en bas de la carte). Violet foncé : pays à faible revenu (la plus grande partie de l'Afrique subsaharienne, l'Afghanistan, le Yémen) ; rose : revenu intermédiaire inférieur (Inde, Nigeria, Cameroun, Maroc\ldots) ; vert clair : revenu intermédiaire supérieur (Chine, Brésil, Afrique du Sud, Mexique\ldots) ; vert foncé : revenu élevé (Amérique du Nord, Europe, Japon, Australie, Russie, pays du Golfe\ldots) ; hachures : pas de données. Sur cette base, on repère la Triade (Amérique du Nord, Europe, Japon), les pays émergents et les pays au décollage difficile (PPTE) ; les flux financiers (IDE), de marchandises et pétroliers du Golfe sont à tracer par l'élève (flèches d'épaisseur proportionnelle), car la carte ne les représente pas. \\ Source : Wikimedia Commons, Our World in Data (données Banque mondiale, 2024), CC BY.}
\end{popfigure}

\emph{Étape 3 — Légende rédigée (à recopier sous la carte).} La couleur bleue désigne les trois pôles de la Triade ; le vert, les pays émergents (BRICS et assimilés) ; le rose, les pays au décollage difficile dont les PPTE d'Afrique ; les flèches violettes, les flux financiers ; les flèches orange, les flux de marchandises ; les flèches dorées, les exportations d'hydrocarbures du Golfe.

\emph{Étape 4 — Conclusion.} La carte révèle une mondialisation \cle{hiérarchisée} : les flux les plus intenses connectent les pôles de la Triade et, de plus en plus, la Chine ; les pays émergents s'insèrent comme ateliers et marchés ; l'Afrique subsaharienne, surtout fournisseuse de matières premières brutes, reste à la marge des flux à forte valeur ajoutée. La mondialisation est donc sélective : elle intègre inégalement les territoires.
\end{exoresolu}

\courssub{Méthode 7 : Construire un raisonnement de synthèse (étude de cas comparée)}

\begin{methode}[Comparer deux trajectoires de développement]
\begin{enumerate}
\item \textbf{Définir les critères} de comparaison avant de rédiger : ressources, insertion dans les échanges, niveau de vie, diversification, contraintes.
\item \textbf{Construire un tableau comparatif} pour organiser les idées (même en brouillon).
\item \textbf{Dégager les facteurs explicatifs} : dotation en ressources, gouvernance, position géographique (littoral/enclavement), politique économique (ouverture, industrialisation), rôle des institutions internationales (FMI, Banque mondiale, initiative PPTE).
\item \textbf{Nuancer} : éviter les jugements définitifs ; un pays du Golfe riche reste vulnérable, un PPTE peut connaître des avancées (allègement de dette, croissance).
\item \textbf{Conclure} en répondant à la problématique par une typologie argumentée.
\end{enumerate}
\end{methode}

\begin{exoresolu}[Comparer le Qatar et le Niger]
\textbf{Énoncé} : À l'aide d'un tableau puis d'un paragraphe argumenté, comparez les trajectoires du Qatar et du Niger et expliquez pourquoi l'un a « décollé » et pas l'autre.
\tcblower
\textbf{Solution détaillée.}

\emph{Étape 1 — Tableau comparatif} (ordres de grandeur) :

\begin{center}
\begin{tabularx}{\linewidth}{|l|X|X|}
\hline
\rowcolor{popblueL} \textbf{Critère} & \textbf{Qatar} & \textbf{Niger} \\
\hline
Position & Littoral du golfe Arabo-Persique & Enclavé au Sahel (Afrique de l'Ouest) \\
\hline
Ressource clé & Gaz naturel (parmi les premières réserves mondiales) et pétrole & Uranium, pétrole modeste, agriculture de subsistance \\
\hline
PIB/habitant (PPA) & Très élevé (environ 90\,000 dollars, parmi les premiers mondiaux) & Très faible (environ 1\,200 dollars, parmi les derniers mondiaux) \\
\hline
Catégorie & Pays du Golfe à revenu élevé & PPTE / PMA \\
\hline
Population & Environ 3 millions, forte majorité d'immigrés & Environ 25 millions, très jeune, forte natalité \\
\hline
Contrainte majeure & Dépendance aux hydrocarbures, main-d'œuvre étrangère & Enclavement, sécheresses, insécurité, dette \\
\hline
\end{tabularx}
\end{center}

\emph{Étape 2 — Paragraphe argumenté.} Le Qatar a transformé une rente gazière exceptionnelle en développement : exportations de GNL (gaz naturel liquéfié) vers l'Asie et l'Europe, investissements dans les infrastructures (port de Hamad, aéroport international de Doha, Coupe du monde 2022), fonds souverain placé dans le monde entier. Sa petite population permet une redistribution massive : PIB par habitant parmi les plus élevés du monde. Le Niger, au contraire, cumule les handicaps : enclavement qui renchérit chaque tonne exportée (l'essentiel transite par le port de Cotonou, à plus de \SI{1000}{km}), climat sahélien qui fragilise l'agriculture, croissance démographique rapide qui absorbe les progrès économiques, poids de la dette qui l'a placé dans l'initiative PPTE. Sa rente uranifère, longtemps captée par des intérêts extérieurs, n'a pas irrigué l'économie nationale.

\emph{Étape 3 — Conclusion.} La comparaison montre que le décollage ne dépend pas seulement des ressources : il résulte de la combinaison d'une rente bien gérée, d'une position maritime ouverte et d'une démographie maîtrisée — conditions réunies au Qatar, absentes au Niger.
\end{exoresolu}

\begin{attention}[Les 5 erreurs qui coûtent des points au Bac]
\begin{enumerate}
\item \textbf{Confondre les catégories} : un pays émergent (Chine, Malaisie) n'est pas un pays à revenus intermédiaires « bloqué », et un PPTE n'est pas synonyme de « pays africain » (Haïti, la Bolivie ont été PPTE). Vérifie toujours la définition avant de classer : PPTE = critère d'endettement (Banque mondiale/FMI), PRI = critère de revenu par habitant.
\item \textbf{Décrire sans localiser ni chiffrer} : un commentaire qui ne cite ni un lieu précis, ni un chiffre du document, ni une date est hors sujet. Le correcteur attend au moins deux données exactes tirées du document.
\item \textbf{Se tromper dans le calcul du taux de variation} : diviser par la valeur finale au lieu de la valeur initiale. Rappel : $(V_f - V_i)/V_i \times 100$. Et ne jamais confondre « 3 fois plus » (rapport 3) avec « 3 \% de plus ».
\item \textbf{Oublier la légende ou la laisser incomplète} : au Bac comme au TP 5, tout figuré utilisé sur le croquis ou la carte doit apparaître dans la légende, avec titre, orientation et échelle indicative. Une carte sans légende perd jusqu'à la moitié des points de la question.
\item \textbf{Affirmer que la rente pétrolière suffit au développement} : c'est le piège classique. Le Nigeria ou le Tchad prouvent qu'une rente mal gérée ne développe pas (malédiction des ressources) ; il faut toujours évoquer la gestion de la rente, la diversification et la gouvernance dans la conclusion.
\end{enumerate}
\end{attention}

\newpage

\coursec{Exercices}

\courssub{Je vérifie mes connaissances}

\begin{exercice}[QCM : Catégories de pays en difficulté]
Parmi les affirmations suivantes, une seule est exacte. La recopier en la complétant.
\begin{enumerate}
\item Les PPTE (Pays Pauvres Très Endettés) sont majoritairement situés en \emph{...} (Afrique subsaharienne / Amérique latine / Asie du Sud-Est).
\item Un pays à revenus intermédiaires (PRI) se caractérise par un RNB/habitant \emph{...} (inférieur à 1 000 \$ / compris entre environ 1 100 \$ et 13 200 \$ / supérieur à 20 000 \$).
\item Le décollage économique des pays du Golfe (Arabie saoudite, Émirats arabes unis, Qatar) repose principalement sur \emph{...} (l'agriculture intensive / la rente pétrolière et gazière / l'industrie automobile).
\item L'initiative PPTE, lancée en 1996 par la Banque mondiale et le FMI, vise à \emph{...} (augmenter les taux d'intérêt / annuler ou réduire la dette extérieure / imposer des barrières douanières).
\end{enumerate}
\end{exercice}

\begin{attention}[Piège fréquent au Bac]
Ne confonds pas \cle{PIB} et \cle{RNB} : la Banque mondiale classe les pays selon le \textbf{revenu national brut (RNB) par habitant}, et non selon le PIB par habitant. Dans les exercices, on utilise souvent le PIB/habitant par approximation : signale-le dans ta copie, cela montre ta rigueur. Autre piège : les PPTE ne sont pas \emph{exclusivement} africains (Haïti, la Bolivie ou le Nicaragua en ont fait partie), même si la grande majorité des 39 pays concernés se trouve en Afrique subsaharienne.
\end{attention}

\begin{exercice}[Vrai ou Faux justifié]
Dire si les affirmations sont vraies ou fausses et justifier brièvement votre réponse en vous appuyant sur des exemples précis (pays, chiffres, organisations).
\begin{enumerate}
\item « Tous les pays d'Afrique sont classés parmi les PPTE. »
\item « Les pays à revenus intermédiaires (comme le Cameroun, le Ghana ou la Côte d'Ivoire) ont tous réussi leur émergence économique. »
\item « La diversification de l'économie est un défi majeur pour les monarchies du Golfe malgré leurs revenus pétroliers. »
\item « L'Indice de Développement Humain (IDH) des PPTE est systématiquement supérieur à 0,7. »
\end{enumerate}
\end{exercice}

\begin{exercice}[Définitions et concepts clés]
Définir précisément les termes suivants en une ou deux phrases, en donnant un exemple de pays pour chacun :
\begin{enumerate}
\item PPTE (Pays Pauvres Très Endettés)
\item Pays à revenus intermédiaires (PRI) — distinguer tranche inférieure et tranche supérieure.
\item Rente pétrolière
\item Initiative PPTE (initiative lancée en 1996, renforcée en 1999, complétée par l'initiative IADM d'annulation multilatérale de la dette en 2005)
\item Décollage économique (au sens de Rostow, appliqué au contexte actuel)
\end{enumerate}
\end{exercice}

\begin{exercice}[Calculs directs et indicateurs]
On donne les données suivantes pour deux pays fictifs en 2022 :
\begin{center}
\begin{tabular}{|l|c|c|}\hline
\textbf{Indicateur} & \textbf{Pays X} & \textbf{Pays Y} \\ \hline
PIB (milliards \$ courants) & 45 & 12 \\
Population (millions d'habitants) & 25 & 8 \\
Dette extérieure totale (milliards \$) & 18 & 9 \\
Exportations de biens et services (\% du PIB) & 35\% & 20\% \\ \hline
\end{tabular}
\end{center}
\begin{enumerate}
\item Calculer le PIB par habitant (en \$) pour chaque pays.
\item Calculer le ratio Dette extérieure / PIB (en \%) pour chaque pays.
\item Classer ces deux pays dans les catégories de la Banque mondiale (revenu faible, intermédiaire tranche inférieure, intermédiaire tranche supérieure, élevé) en utilisant les seuils approximatifs de l'exercice 2022 : faible < 1 085 \$ ; intermédiaire inférieur de 1 086 à 4 255 \$ ; intermédiaire supérieur de 4 256 à 13 205 \$ ; élevé > 13 205 \$.
\item Quel pays présente la situation la plus préoccupante au regard de l'endettement ? Justifier.
\end{enumerate}
\end{exercice}

\courssub{Je m'entraîne}

\begin{exercice}[Classement et analyse de pays africains]
À partir du tableau ci-dessous (ordres de grandeur 2022--2023, à manier avec prudence car les sources varient), répondre aux questions.
\begin{center}
\begin{tabular}{|l|c|c|c|l|}\hline
\textbf{Pays} & \textbf{PIB/hab (\$)} & \textbf{IDH (2022)} & \textbf{Dette ext./PIB} & \textbf{Principales exportations} \\ \hline
Cameroun & 1 600 & 0,587 & 45\% & Pétrole, bois, cacao \\
Niger & 590 & 0,400 & 55\% & Uranium, or, bétail \\
Sénégal & 1 600 & 0,517 & 70\% & Phosphates, poisson, arachide \\
Gabon & 8 500 & 0,706 & 60\% & Pétrole, manganèse, bois \\
Éthiopie & 1 000 & 0,492 & 25\% & Café, or, fleurs \\ \hline
\end{tabular}
\end{center}
\begin{enumerate}
\item Classer ces cinq pays en : PPTE, Pays à revenus intermédiaires (préciser la tranche), Pays les moins avancés (PMA).
\item Identifier le pays qui a le meilleur IDH et celui qui a le plus faible PIB/hab. Expliquer le décalage possible entre PIB/hab et IDH pour le Gabon.
\item Pour le Cameroun et le Sénégal (PIB/hab proche), comparer le niveau d'endettement et proposer une hypothèse sur l'usage de la dette.
\item En quoi l'exemple de l'Éthiopie (croissance rapide, dette extérieure modérée, IDH faible) illustre-t-il les limites du PIB/hab comme indicateur unique de développement ?
\end{enumerate}
\end{exercice}

\begin{exercice}[Le Golfe : rente et diversification — Analyse de documents]
\textbf{Document 1 :} Part des hydrocarbures dans les exportations de marchandises (ordres de grandeur 2022) : Arabie saoudite environ 75\%, Émirats arabes unis environ 30\%, Qatar environ 85\%, Koweït environ 90\%.
\textbf{Document 2 :} Projets de diversification : Vision 2030 de l'Arabie saoudite (ville nouvelle de Neom, tourisme, industrie, finance) ; Dubaï (hub logistique avec le port de Jebel Ali, tourisme, immobilier, finance).
\begin{enumerate}
\item Construire un diagramme en barres comparant la part des hydrocarbures dans les exportations des quatre pays.
\item Expliquer pourquoi le Qatar et le Koweït apparaissent plus vulnérables que les Émirats face à une baisse durable des cours du pétrole et du gaz.
\item Citer deux secteurs de diversification privilégiés par l'Arabie saoudite et les Émirats. Pourquoi ces secteurs sont-ils stratégiques dans la mondialisation ?
\item Montrer que, malgré la diversification, ces pays restent des foyers énergétiques de la mondialisation (flux sortants d'hydrocarbures, flux entrants de capitaux et de main-d'œuvre).
\end{enumerate}
\end{exercice}

\begin{exercice}[Croquis : Le Golfe persique, foyer pétrolier de la mondialisation]
Sur un fond de carte schématique de la péninsule Arabique et du Golfe, réaliser le croquis légendé : « Le Golfe persique : un foyer de la mondialisation fondé sur la rente hydrocarbure ».
\begin{enumerate}
\item Tracer et nommer les 6 pays du CCG (Conseil de coopération du Golfe) : Arabie saoudite, Émirats arabes unis, Qatar, Koweït, Bahreïn, Oman.
\item Localiser les principaux gisements (Ghawar en Arabie saoudite, Burgan au Koweït, North Field/South Pars partagé entre le Qatar et l'Iran) par un symbole adapté.
\item Représenter les terminaux d'exportation (Ras Tanura et Jubail en Arabie saoudite, Fujairah aux Émirats, île de Kharg en Iran) et les détroits stratégiques (Ormuz, Bab el-Mandeb).
\item Figurer par des flèches d'épaisseur proportionnelle les principaux flux de pétrole et de GNL vers l'Asie (Chine, Inde, Japon, Corée du Sud) et, plus faibles, vers l'Europe et l'Amérique du Nord.
\item Insérer les grands ports-hubs (Jebel Ali/Dubaï, Khalifa/Abou Dhabi, Hamad/Doha) et les aéroports-hubs (Dubaï, Doha, Riyad) comme nœuds tertiaires de la mondialisation.
\item Rédiger une légende organisée en 3 items : Gisements et production ; Flux d'hydrocarbures ; Nœuds logistiques et financiers.
\end{enumerate}
\end{exercice}

\begin{methode}[Réussir un croquis géographique au Bac]
\begin{enumerate}
\item Tracer d'abord le fond (contours simplifiés, échelle indicative, flèche du nord) au crayon léger.
\item Placer les éléments du plus grand au plus petit : espaces, puis flux (flèches proportionnelles), puis points (villes, ports, gisements).
\item Rédiger la légende à côté, organisée en items numérotés correspondant aux symboles : un symbole = une ligne de légende.
\item Donner un titre précis au croquis (lieu + thème + idée directrice).
\end{enumerate}
\end{methode}

\begin{exercice}[Le Cameroun : pays à revenus intermédiaires, décollage difficile]
À partir de vos connaissances et des données ci-dessous, rédiger un paragraphe argumenté d'environ 15 lignes.
\textbf{Données :} PIB/hab environ 1 600 \$ (PRI, tranche inférieure) ; IDH 0,587 (développement humain moyen, bas de la catégorie) ; taux de pauvreté environ 37\% (seuil national) ; dette publique environ 45\% du PIB (risque de surendettement jugé élevé mais soutenable par le FMI) ; structure des exportations : pétrole brut environ 40\%, bois environ 15\%, cacao-café-coton environ 15\%, aluminium environ 10\% ; projets structurants : barrage de Nachtigal (420 MW), port en eau profonde de Kribi, autoroute Yaoundé--Douala, adhésion à la ZLECAf.
\begin{enumerate}
\item Justifier le classement du Cameroun dans les PRI (tranche inférieure).
\item Identifier deux facteurs structurels qui freinent le décollage (dépendance aux matières premières, poids du secteur informel, déficit d'infrastructures, gouvernance).
\item Montrer comment l'intégration régionale (CEMAC, ZLECAf) et les grands chantiers d'infrastructures constituent des leviers potentiels.
\item Conclure sur le défi de la transformation structurelle (industrialisation, capital humain) pour passer de la croissance à l'émergence (Vision 2035).
\end{enumerate}
\end{exercice}

\courssub{Je me prépare au Bac}

\begin{exercice}[Sujet type Bac : Étude de documents — Différenciation des pays en difficulté]
\textbf{Document 1 : Tableau comparatif (ordres de grandeur 2022)}
\begin{center}
\begin{tabular}{|l|c|c|c|l|}\hline
\textbf{Pays} & \textbf{PIB/hab (\$)} & \textbf{IDH} & \textbf{Dette ext./PIB} & \textbf{Statut FMI/BM} \\ \hline
Tchad & 700 & 0,394 & 50\% & PPTE (point d'achèvement atteint) \\
Côte d'Ivoire & 2 500 & 0,550 & 55\% & PRI (tranche inférieure) \\
Nigéria & 2 100 & 0,548 & 30\% & PRI (tranche inférieure) \\
Arabie saoudite & 30 000 & 0,875 & 25\% & Revenu élevé \\
Bangladesh & 2 700 & 0,670 & 20\% & PRI (tranche inférieure) \\ \hline
\end{tabular}
\end{center}
\textbf{Document 2 :} Extrait adapté d'un rapport de la Banque mondiale (2023) : « Les PPTE restent vulnérables aux chocs climatiques et aux fluctuations des cours des matières premières. L'allègement de la dette libère des ressources pour la santé et l'éducation, mais la soutenabilité budgétaire nécessite une mobilisation accrue des recettes intérieures. »
\textbf{Document 3 :} Schéma (description) : « Pays du Golfe : rente pétrolière $\rightarrow$ fonds souverains (SWF) $\rightarrow$ diversification (tourisme, finance, logistique, industrie) $\rightarrow$ projection mondiale (Neom, Expo 2020 Dubaï, Coupe du monde Qatar 2022). »

\begin{enumerate}
\item \textbf{À partir du Document 1,} classer les cinq pays dans les catégories du programme (PPTE, PRI, pays du Golfe à revenu élevé). Justifier chaque classement par deux indicateurs.
\item \textbf{À partir des Documents 1 et 2,} expliquer pourquoi le Tchad, avec un ratio Dette/PIB proche de celui de la Côte d'Ivoire, relève du statut PPTE alors que la Côte d'Ivoire en est sortie.
\item \textbf{À partir du Document 3,} montrer que les pays du Golfe utilisent leur rente pour s'insérer activement dans la mondialisation, et ne se contentent pas d'une rente de situation.
\item \textbf{Synthèse :} « La catégorie de "pays à décollage difficile" masque des réalités très contrastées. » Rédiger un développement construit d'environ 20 lignes en vous appuyant sur l'ensemble des documents et vos connaissances pour discuter cette affirmation (opposition PPTE/PRI, rôle de la rente, trajectoires de diversification).
\end{enumerate}
\end{exercice}

\begin{exercice}[Sujet type Bac : Croquis et commentaire — Les foyers de la mondialisation (TP 5)]
\textbf{Consigne :} Sur un fond de mappemonde fourni, réaliser le croquis suivant : « Les foyers de la mondialisation : Triade, pays émergents, flux majeurs ».
\begin{enumerate}
\item \textbf{Construction du croquis (12 points) :}
\begin{itemize}
\item Délimiter et nommer les trois pôles de la Triade (Amérique du Nord, Europe occidentale, Asie orientale : Japon, Corée du Sud, Taïwan, façade est de la Chine).
\item Localiser et nommer 5 grands pays émergents (BRICS et assimilés) : Chine, Inde, Brésil, Russie, Afrique du Sud (ou Indonésie, Mexique).
\item Représenter par des flèches d'épaisseur variable les 4 principaux flux de marchandises : transpacifique (Asie--Amérique du Nord), transatlantique (Europe--Amérique du Nord), Europe--Asie (par Suez), intra-asiatique.
\item Figurer 2 grands flux financiers (IDE) : sortants de la Triade vers les émergents ; investissements croisés Sud-Sud (Chine vers l'Afrique, fonds du Golfe vers le monde).
\item Localiser 3 grands ports mondiaux (Shanghai, Singapour, Rotterdam, Los Angeles/Long Beach, Dubaï/Jebel Ali, Busan) et 2 aéroports-hubs (Dubaï, Londres, Hong Kong, Singapour).
\end{itemize}
\item \textbf{Légende organisée (4 points) :} rédiger la légende complète en 3 items : Pôles de puissance ; Espaces émergents ; Flux et nœuds de la mondialisation.
\item \textbf{Commentaire (4 points) :} en 10 lignes, expliquer comment ce croquis met en évidence la multipolarisation de la mondialisation et le rôle charnière des pays émergents.
\end{enumerate}
\end{exercice}

\begin{exercice}[Sujet type Bac : Développement construit — Le Cameroun face aux défis de l'émergence]
\textbf{Sujet :} « Le Cameroun, pays à revenus intermédiaires de la tranche inférieure, peine à enclencher un décollage économique durable. » Discutez cette affirmation en montrant les atouts, les freins structurels et les leviers d'action publique.

\textbf{Plan suggéré (à suivre pour la rédaction) :}
\begin{enumerate}
\item \textbf{Introduction :} accroche (Vision 2035, statut de PRI), définition du « décollage difficile », annonce du plan (atouts / freins / leviers).
\item \textbf{I. Des atouts réels mais sous-exploités (le pays « Afrique en miniature »).}
\begin{itemize}
\item Ressources diversifiées (agricoles, minières, énergétiques, forestières, touristiques).
\item Position stratégique (golfe de Guinée, porte d'entrée de l'Afrique centrale, corridor Douala--N'Djaména et Douala--Bangui).
\item Démographie jeune, marché intérieur en expansion, accès à la ZLECAf.
\end{itemize}
\item \textbf{II. Des freins structurels persistants (le « piège du revenu intermédiaire »).}
\begin{itemize}
\item Dépendance aux matières premières (pétrole, bois, cacao) : vulnérabilité aux chocs externes, faible valeur ajoutée locale.
\item Déficit d'infrastructures (énergie, transport, numérique) et coûts de production élevés.
\item Gouvernance : corruption, climat des affaires perfectible, secteur informel dominant (environ 90\% de l'emploi), évasion fiscale.
\item Capital humain : inadéquation formation-emploi, santé et éducation à renforcer (IDH 0,587).
\end{itemize}
\item \textbf{III. Leviers d'action et pistes de transformation structurelle.}
\begin{itemize}
\item Industrialisation par substitution aux importations et transformation locale (cacao, bois, coton, alumine).
\item Infrastructures structurantes (barrages de Nachtigal et Memve'ele, port de Kribi, autoroutes, numérique) et partenariats public-privé.
\item Amélioration du climat des affaires, lutte contre la corruption, mobilisation des recettes fiscales (pression fiscale inférieure à 15\% du PIB).
\item Intégration régionale effective (CEMAC, ZLECAf) : marché élargi, chaînes de valeur régionales.
\item Capital humain : enseignement technique et professionnel, santé, protection sociale.
\end{itemize}
\item \textbf{Conclusion :} bilan nuancé : le décollage est possible (exemples du Vietnam ou du Bangladesh il y a vingt ans) mais suppose une rupture dans la gouvernance et le modèle de croissance (inclusif, industriel, durable). Ouverture sur le rôle de la diaspora et des financements verts.
\end{enumerate}
\textbf{Consignes de rédaction :} 1,5 à 2 pages. Utiliser des données chiffrées précises (PIB/hab, IDH, parts des exportations, dette, population). S'appuyer sur des exemples concrets (usines de transformation du cacao, barrage de Nachtigal, ZLECAf, Stratégie nationale de développement SND30, 2020--2030).
\end{exercice}

\begin{aretenir}[L'essentiel de la leçon]
\begin{itemize}
\item Les \cle{PPTE} (environ 39 pays, surtout en Afrique subsaharienne) cumulent pauvreté, IDH faible et dette extérieure insoutenable ; l'initiative PPTE (1996, renforcée 1999, complétée par l'IADM en 2005) allège leur dette contre des réformes et des investissements sociaux.
\item Les \cle{PRI} (RNB/hab entre environ 1 100 et 13 200 \$) forment un groupe hétérogène : certains émergent (Bangladesh, Vietnam), d'autres stagnent dans le « piège du revenu intermédiaire » (Cameroun, Kenya).
\item Les \cle{pays du Golfe} ont décollé grâce à la rente pétrolière et gazière, qu'ils convertissent en diversification (logistique, finance, tourisme) via leurs fonds souverains : ce sont des foyers de la mondialisation, mais vulnérables à la transition énergétique.
\item En Afrique, le décollage est freiné par la dépendance aux matières premières, le déficit d'infrastructures et de capital humain ; la ZLECAf ouvre une perspective d'intégration.
\end{itemize}
\end{aretenir}

\newpage

\coursec{Corrigés des exercices}

\coursec{Leçon 15 : Les pays au décollage économique difficile — Exercices corrigés et analyses}

\begin{aretenir}[Seuils de classification de la Banque mondiale — Repères chiffrés]
La Banque mondiale révise chaque année (juillet) les seuils de RNB par habitant (méthode Atlas) pour classer les pays. Les valeurs ci-dessous sont des \textbf{ordres de grandeur} (exercice 2023-2024) à citer avec « environ » au Bac :
\begin{itemize}
\item \textbf{Revenu faible :} RNB/hab $\le$ 1\,135 \$ (ex. : Niger, Tchad, Éthiopie).
\item \textbf{Revenu intermédiaire tranche inférieure :} 1\,136 \$ -- 4\,465 \$ (ex. : Cameroun, Sénégal, Côte d'Ivoire, Bangladesh, Nigeria).
\item \textbf{Revenu intermédiaire tranche supérieure :} 4\,466 \$ -- 13\,845 \$ (ex. : Gabon, Chine, Afrique du Sud, Brésil).
\item \textbf{Revenu élevé :} $>$ 13\,845 \$ (ex. : Arabie saoudite, Qatar, Émirats, pays de l'OCDE).
\end{itemize}
\textbf{Attention :} les seuils 2022 (1\,086 / 4\,255 / 13\,205 \$) figurent encore dans certains manuels ; indique toujours l'année de référence ou écris « environ ».
\end{aretenir}

\begin{corrige}[Exercice 1 — QCM : Catégories de pays en difficulté]
\begin{enumerate}
\item Les PPTE sont majoritairement situés en \textbf{Afrique subsaharienne}. Sur les 39 pays éligibles à l'initiative PPTE, plus de 30 se trouvent en Afrique subsaharienne (Tchad, Niger, Mali, Burkina Faso, Éthiopie, RDC\ldots). Quelques exceptions existent : Haïti, la Bolivie, le Nicaragua, le Honduras, la Guyane\ldots
\item Un PRI se caractérise par un RNB/habitant \textbf{compris entre environ 1\,100 \$ et 13\,800 \$} (seuils 2023-2024 : tranche inférieure 1\,136--4\,465 \$, tranche supérieure 4\,466--13\,845 \$).
\item Le décollage des pays du Golfe repose principalement sur \textbf{la rente pétrolière et gazière}. L'Arabie saoudite, le Qatar (premier exportateur mondial de GNL avec l'Australie et les États-Unis) et le Koweït tirent l'essentiel de leurs recettes d'exportation des hydrocarbures.
\item L'initiative PPTE vise à \textbf{annuler ou réduire la dette extérieure}. Lancée en 1996 par la Banque mondiale et le FMI, renforcée en 1999, elle allège la dette des pays éligibles en contrepartie de réformes et de dépenses sociales (santé, éducation).
\end{enumerate}
\textbf{Point de vigilance :} au Bac, cite toujours les seuils avec l'indication « environ » ou « ordre de grandeur », car la Banque mondiale les révise chaque année.
\end{corrige}

\begin{corrige}[Exercice 2 — Vrai ou Faux justifié]
\begin{enumerate}
\item \textbf{Faux.} Tous les pays africains ne sont pas des PPTE : l'Afrique du Sud, le Maroc, l'Algérie, l'Égypte ou le Gabon n'en font pas partie. Le Gabon est même un PRI de tranche supérieure (RNB/hab environ 8\,500 \$). Environ 39 pays sont concernés par l'initiative PPTE, dont la majorité en Afrique subsaharienne, mais pas la totalité des 54 États africains.
\item \textbf{Faux.} Les PRI n'ont pas tous émergé : le Cameroun (RNB/hab environ 1\,600 \$, IDH 0,587) stagne dans le « piège du revenu intermédiaire », tandis que le Bangladesh ou le Vietnam, partis d'un niveau comparable, connaissent une émergence rapide grâce à l'industrie manufacturière (textile). Le Ghana et la Côte d'Ivoire restent dépendants du cacao et de l'or.
\item \textbf{Vrai.} Malgré la rente, la diversification est vitale : les hydrocarbures représentent encore environ 75\% des exportations saoudiennes et 85\% des exportations qataries. La transition énergétique mondiale menace ce modèle, d'où la Vision 2030 (Neom, tourisme) et le modèle de Dubaï (logistique, finance).
\item \textbf{Faux.} L'IDH des PPTE est généralement \emph{inférieur} à 0,7 : Tchad 0,394, Niger 0,400, Éthiopie 0,492 (ordres de grandeur 2022). Un IDH supérieur à 0,7 correspond au « développement humain élevé », incompatible avec le statut de pays pauvre très endetté.
\end{enumerate}
\textbf{Point de vigilance :} une justification au Bac exige toujours au moins un exemple chiffré ou un pays nommé ; une réponse « vrai/faux » seule ne rapporte aucun point.
\end{corrige}

\begin{corrige}[Exercice 3 — Définitions et concepts clés]
\begin{enumerate}
\item \textbf{PPTE (Pays Pauvres Très Endettés) :} pays à faible revenu dont la dette extérieure est jugée insoutenable par la Banque mondiale et le FMI, et qui bénéficient à ce titre d'un allègement de dette conditionné à des réformes. Exemple : le Tchad, le Niger.
\item \textbf{Pays à revenus intermédiaires (PRI) :} pays dont le RNB par habitant se situe entre environ 1\,100 \$ et 13\,800 \$ (seuils Banque mondiale 2023-2024). On distingue la \emph{tranche inférieure} (environ 1\,136--4\,465 \$ ; exemple : le Cameroun, environ 1\,600 \$) et la \emph{tranche supérieure} (environ 4\,466--13\,845 \$ ; exemple : le Gabon, environ 8\,500 \$, ou la Chine).
\item \textbf{Rente pétrolière :} revenu tiré de l'extraction et de l'exportation d'hydrocarbures, largement indépendant de l'effort productif national, et capté essentiellement par l'État (ou les compagnies nationales comme Saudi Aramco). Exemple : l'Arabie saoudite, où le pétrole fournit l'essentiel des recettes budgétaires et d'exportation.
\item \textbf{Initiative PPTE :} dispositif lancé en 1996 par la Banque mondiale et le FMI, renforcé en 1999, visant à réduire la dette extérieure des pays pauvres jusqu'à un niveau « soutenable » (point de décision, puis point d'achèvement), complété en 2005 par l'initiative IADM (MDRI) qui annule la dette multilatérale restante. Exemple : le Cameroun a atteint le point d'achèvement en 2006.
\item \textbf{Décollage économique (Rostow) :} phase de l'évolution d'une économie (théorie de W. Rostow, 1960) où l'investissement productif dépasse durablement environ 10\% du revenu national et où une industrie motrice entraîne une croissance auto-entretenue. Appliqué aujourd'hui, il désigne le passage d'une économie de rente ou de subsistance à une économie industrialisée et diversifiée. Exemple : le « décollage » réussi de la Corée du Sud dans les années 1960--1980, attendu mais difficile pour le Cameroun.
\end{enumerate}
\end{corrige}

\begin{corrige}[Exercice 4 — Calculs directs et indicateurs]
\begin{enumerate}
\item \textbf{PIB par habitant.} On applique : PIB/hab $=$ PIB $\div$ population.
\begin{align*}
\text{Pays X : } & \frac{45 \text{ milliards \$}}{25 \text{ millions hab}} = \frac{45\,000\,000\,000}{25\,000\,000} = 1\,800 \text{ \$/hab} \\
\text{Pays Y : } & \frac{12 \text{ milliards \$}}{8 \text{ millions hab}} = \frac{12\,000\,000\,000}{8\,000\,000} = 1\,500 \text{ \$/hab}
\end{align*}
\item \textbf{Ratio Dette extérieure / PIB.}
\begin{align*}
\text{Pays X : } & \frac{18}{45} = 0{,}40 \text{ soit } 40\,\% \\
\text{Pays Y : } & \frac{9}{12} = 0{,}75 \text{ soit } 75\,\%
\end{align*}
\item \textbf{Classement.} Pays X : 1\,800 \$/hab $\in$ [1\,136 ; 4\,465] $\Rightarrow$ \textbf{revenu intermédiaire, tranche inférieure}. Pays Y : 1\,500 \$/hab $\in$ [1\,136 ; 4\,465] $\Rightarrow$ \textbf{revenu intermédiaire, tranche inférieure}. Les deux pays appartiennent à la même catégorie malgré des situations d'endettement très différentes.
\item \textbf{Situation la plus préoccupante : le Pays Y.} Trois arguments : (a) son ratio Dette/PIB de 75\% dépasse largement le seuil de vigilance de 40\% couramment retenu pour les pays en développement ; (b) il est presque deux fois plus endetté que le Pays X (75\% contre 40\%) alors que son PIB/hab est plus faible (1\,500 \$ contre 1\,800 \$) ; (c) ses exportations ne représentent que 20\% du PIB (contre 35\% pour X) : il génère donc moins de devises pour servir une dette extérieure libellée en monnaies étrangères. Le risque de surendettement y est maximal.
\end{enumerate}
\textbf{Point de vigilance :} attention aux unités : convertir les milliards et les millions avant de diviser, ou vérifier que le rapport milliards/millions donne bien des milliers de dollars ($45/25 = 1{,}8$ millier de dollars $= 1\,800$ \$).
\end{corrige}

\begin{corrige}[Exercice 5 — Classement et analyse de pays africains]
\begin{enumerate}
\item \textbf{Classement :}
\begin{itemize}
\item \textbf{PPTE et PMA :} Niger (590 \$/hab, IDH 0,400) et Éthiopie (1\,000 \$/hab, IDH 0,492) — revenu faible, développement humain faible, tous deux sur la liste des PMA de l'ONU et éligibles à l'initiative PPTE.
\item \textbf{PRI, tranche inférieure :} Cameroun et Sénégal (environ 1\,600 \$/hab, entre 1\,136 et 4\,465 \$). Ces deux pays ont bénéficié de l'initiative PPTE (points d'achèvement atteints en 2006 pour le Cameroun, 2004 pour le Sénégal) tout en étant aujourd'hui classés PRI.
\item \textbf{PRI, tranche supérieure :} Gabon (environ 8\,500 \$/hab, entre 4\,466 et 13\,845 \$).
\end{itemize}
\item \textbf{Meilleur IDH :} Gabon (0,706) ; \textbf{plus faible PIB/hab :} Niger (590 \$). \textbf{Décalage Gabon :} le PIB/hab élevé repose sur la rente pétrolière et minière (pétrole, manganèse), concentrée entre peu de mains et peu créatrice d'emplois. Le PIB/hab est une \emph{moyenne} qui ne dit rien de la redistribution : avec des inégalités fortes et des services de santé et d'éducation perfectibles, l'IDH (0,706) reste inférieur à ce que laisserait attendre un revenu de 8\,500 \$.
\item \textbf{Cameroun vs Sénégal :} à richesse par habitant quasi identique (1\,600 \$), le Sénégal est beaucoup plus endetté (70\% du PIB contre 45\%). \textbf{Hypothèse :} le Sénégal a massivement emprunté pour financer les infrastructures du Plan Sénégal Émergent (autoroutes à péage, aéroport international Blaise-Diagne de Diass, train express régional Dakar--Diamniadio), en pariant que la croissance future générée par ces investissements permettra le remboursement — stratégie risquée mais cohérente si les projets sont rentables.
\item \textbf{Leçon de l'Éthiopie :} avec une croissance rapide (souvent proche de 8--10\% par an dans les années 2010) et une dette extérieure modérée (25\% du PIB), l'Éthiopie semblerait « réussir » si l'on ne regardait que ces indicateurs. Or, avec plus de 120 millions d'habitants, le PIB/hab reste faible (1\,000 \$) et l'IDH (0,492) traduit des retards majeurs en santé, éducation et niveau de vie, aggravés par les conflits. Le PIB/hab ne mesure ni la redistribution, ni le développement humain, ni la stabilité : il ne peut être un indicateur unique de développement.
\end{enumerate}
\end{corrige}

\begin{corrige}[Exercice 6 — Le Golfe : rente et diversification]
\begin{enumerate}
\item \textbf{Diagramme en barres} (hauteurs proportionnelles aux pourcentages) :

\begin{popfigure}
\begin{tikzpicture}
\begin{axis}[
  ybar, bar width=0.8cm, width=0.9\linewidth, height=6cm,
  ymin=0, ymax=100,
  ylabel={Part des hydrocarbures (\% des exportations)},
  symbolic x coords={ArabieSaoudite, Emirats, Qatar, Koweit},
  xtick=data,
  xticklabels={Arabie saoudite, Émirats, Qatar, Koweït},
  x tick label style={font=\small},
  nodes near coords, nodes near coords align={vertical},
]
\addplot[fill=popblue] coordinates {(ArabieSaoudite,75) (Emirats,30) (Qatar,85) (Koweit,90)};
\end{axis}
\end{tikzpicture}
\legende{Part des hydrocarbures dans les exportations de marchandises de quatre pays du Golfe (ordres de grandeur 2022, sources : OPEP, Banque mondiale).}
\end{popfigure}

\item \textbf{Vulnérabilité comparée :} le Qatar (85\%) et le Koweït (90\%) dépendent quasi exclusivement des hydrocarbures pour leurs recettes d'exportation : une baisse durable des cours du pétrole et du gaz amputerait directement leurs recettes budgétaires et leurs excédents commerciaux, sans secteur de substitution. Les Émirats (30\%), grâce à Dubaï (port de Jebel Ali, tourisme, finance, immobilier, compagnie Emirates), disposent de sources de revenus non pétrolières qui amortissent le choc.
\item \textbf{Secteurs de diversification :} Arabie saoudite : tourisme et industrie/finance (Vision 2030, ville nouvelle de Neom) ; Émirats : logistique portuaire et aérienne (Jebel Ali, Emirates, Etihad) et finance. Ces secteurs sont stratégiques car ils insèrent ces pays dans les \emph{flux tertiaires} de la mondialisation (capitaux, marchandises conteneurisées, passagers), moins exposés à l'épuisement des ressources et à la transition énergétique que la seule rente.
\item \textbf{Foyers énergétiques malgré tout :} ces pays restent au cœur des flux énergétiques mondiaux : flux sortants de pétrole et de GNL vers l'Asie (Chine, Inde, Japon, Corée du Sud) par le détroit d'Ormuz ; flux entrants de capitaux (IDE attirés par la stabilité et les grands projets) et de main-d'œuvre (les migrants représentent environ 85\% de la population du Qatar et des Émirats). La diversification elle-même est financée par la rente via les fonds souverains (PIF saoudien, QIA qatari, ADIA émirien).
\end{enumerate}
\textbf{Point de vigilance :} ne pas écrire que les Émirats sont « sortis du pétrole » : Abou Dhabi reste pétrolier ; c'est la fédération, tirée par Dubaï, qui est diversifiée.
\end{corrige}

\begin{corrige}[Exercice 7 — Croquis : Le Golfe persique, foyer pétrolier]
\textbf{Corrigé-type (éléments attendus) :}

\begin{popfigure}
\includegraphics[width=\linewidth,height=9.5cm,keepaspectratio]{N15/golfe_petrole.jpg}
\legende{Fond de carte pour le croquis du Golfe persique (carte en anglais). Elle fournit les repères du croquis attendu : 1. les gisements (champs de pétrole en vert, champs de gaz en rose) ; 2. les oléoducs et gazoducs ; 3. les villes et capitales (Riyad, Koweït, Doha, Abou Dhabi, Mascate) et la façade maritime sur le golfe Persique, le golfe d'Oman et la mer Rouge. Les flèches de flux d'hydrocarbures proportionnelles (Asie $\gg$ Europe $>$ Amérique du Nord) et le détroit d'Ormuz sont à ajouter par l'élève. \\ Source : Wikimedia Commons, Energy Information Administration (États-Unis), domaine public.}
\end{popfigure}

\textbf{Barème indicatif (sur 20) :} fond de carte lisible avec nord et titre : 3 pts ; 6 pays du CCG correctement placés : 3 pts ; 3 gisements avec symbole : 3 pts ; terminaux et détroits (Ormuz, Bab el-Mandeb) : 4 pts ; flux proportionnels (Asie dominante) : 4 pts ; légende organisée en 3 items : 3 pts.

\textbf{Points de vigilance :} (a) le Qatar est une presqu'île rattachée à l'Arabie saoudite, pas une île ; (b) le détroit d'Ormuz est le passage obligé des tankers : environ un cinquième du pétrole mondial y transite ; (c) l'épaisseur des flèches doit traduire la hiérarchie des flux (Asie $\gg$ Europe $>$ Amérique du Nord) ; (d) un symbole non légendé ne rapporte rien.
\end{corrige}

\begin{corrige}[Exercice 8 — Le Cameroun : paragraphe argumenté]
\textbf{Proposition de paragraphe (environ 15 lignes) :}

Le Cameroun est un pays à revenus intermédiaires de la tranche inférieure : avec un PIB par habitant d'environ 1\,600 \$, il se situe dans la fourchette 1\,136--4\,465 \$ définie par la Banque mondiale, et son IDH de 0,587 le place en bas de la catégorie du développement humain moyen. Ce classement masque un décollage difficile, freiné par des facteurs structurels. D'abord, la dépendance aux matières premières : le pétrole brut fournit environ 40\% des exportations, auxquels s'ajoutent le bois (15\%), le cacao-café-coton (15\%) et l'aluminium (10\%), produits peu transformés localement et exposés aux fluctuations des cours mondiaux. Ensuite, le déficit d'infrastructures et le poids du secteur informel (environ 90\% de l'emploi) renchérissent les coûts de production et limitent la mobilisation fiscale, alors que la dette publique atteint déjà environ 45\% du PIB et que 37\% de la population vit sous le seuil national de pauvreté. Des leviers existent cependant : les grands chantiers structurants (barrage de Nachtigal, 420 MW ; port en eau profonde de Kribi ; autoroute Yaoundé--Douala) peuvent lever les goulets d'étranglement énergétiques et logistiques, tandis que l'intégration régionale (CEMAC, et surtout la ZLECAf dont le Cameroun est partie prenante) élargit les débouchés vers un marché africain de 1,4 milliard de consommateurs. Le défi est donc celui de la transformation structurelle : industrialiser (transformation locale du cacao, du bois, du coton), investir dans le capital humain et améliorer la gouvernance, afin de convertir la croissance en émergence effective à l'horizon de la Vision 2035.

\textbf{Barème indicatif (sur 10) :} justification du classement PRI avec seuils : 2 pts ; deux freins structurels identifiés et chiffrés : 3 pts ; leviers (infrastructures + intégration régionale) explicités : 3 pts ; conclusion sur la transformation structurelle : 2 pts.

\textbf{Points de vigilance :} un paragraphe argumenté n'est pas une liste : les idées doivent être enchaînées par des connecteurs logiques ; chaque affirmation doit être appuyée par un chiffre ou un exemple du sujet.
\end{corrige}

\begin{corrige}[Exercice 9 — Sujet type Bac : Différenciation des pays en difficulté]
\begin{enumerate}
\item \textbf{Classement (Document 1) :}
\begin{itemize}
\item \textbf{PPTE : Tchad} — PIB/hab de 700 \$ (revenu faible, sous le seuil de 1\,135 \$) et IDH de 0,394 (développement humain faible) ; le tableau confirme le statut PPTE avec point d'achèvement atteint.
\item \textbf{PRI tranche inférieure : Côte d'Ivoire, Nigéria, Bangladesh} — PIB/hab compris entre 1\,136 et 4\,465 \$ (2\,500, 2\,100 et 2\,700 \$) et IDH de développement humain moyen (0,550 ; 0,548 ; 0,670).
\item \textbf{Pays du Golfe à revenu élevé : Arabie saoudite} — PIB/hab d'environ 30\,000 \$ (supérieur au seuil de 13\,845 \$) et IDH très élevé (0,875), richesse fondée sur la rente pétrolière.
\end{itemize}
\item \textbf{Tchad vs Côte d'Ivoire :} à ratio Dette/PIB comparable (50\% contre 55\%), c'est la \emph{capacité de remboursement} qui diffère. Le Tchad a un revenu par habitant trois fois plus faible (700 \$ contre 2\,500 \$), une assiette fiscale étroite et des exportations quasi mono-produit (pétrole) : le Document 2 rappelle que les PPTE restent « vulnérables aux chocs climatiques et aux fluctuations des cours des matières premières ». La Côte d'Ivoire, économie plus diversifiée (cacao, café, services, industrie naissante) et plus riche, a pu atteindre le point d'achèvement de l'initiative PPTE (2012) et « sortir » du statut, sa dette étant devenue soutenable. Le Document 2 souligne la condition : « la mobilisation accrue des recettes intérieures ».
\item \textbf{Insertion active dans la mondialisation (Document 3) :} les pays du Golfe ne se contentent pas d'exporter du brut (rente de situation) : ils recyclent la rente via leurs \emph{fonds souverains} (SWF) qui investissent dans le monde entier, et financent la diversification (tourisme, finance, logistique, industrie). Les grands événements et projets (Neom, Expo 2020 Dubaï, Coupe du monde Qatar 2022) sont des instruments de projection mondiale et d'attractivité (capitaux, touristes, main-d'œuvre qualifiée). Ils deviennent ainsi des \emph{pôles} de la mondialisation, et non de simples réservoirs énergétiques.
\item \textbf{Synthèse — proposition de développement construit (environ 20 lignes) :}

L'expression « pays à décollage difficile » rassemble sous une même étiquette des situations profondément contrastées, que les documents invitent à distinguer. D'un côté, les PPTE comme le Tchad cumulent les handicaps : revenu faible (700 \$/hab), IDH très bas (0,394), dépendance à une matière première et dette insoutenable qui n'a pu être résorbée que par l'initiative PPTE ; le Document 2 montre que l'allègement de dette y libère des ressources pour la santé et l'éducation, condition première du développement humain. De l'autre, les PRI comme la Côte d'Ivoire, le Nigéria ou le Bangladesh disposent d'un revenu intermédiaire et d'économies en voie de diversification : le Bangladesh (IDH 0,670, dette 20\% du PIB) illustre une trajectoire d'émergence par l'industrie textile exportatrice, tandis que d'autres PRI stagnent dans le « piège du revenu intermédiaire ». À l'opposé, les pays du Golfe, classés « revenu élevé » (Arabie saoudite : 30\,000 \$/hab, IDH 0,875), ont décollé grâce à la rente pétrolière ; mais le Document 3 révèle que leur défi est symétrique : convertir une rente épuisable en économie diversifiée avant la transition énergétique mondiale. Ainsi, la « difficulté » du décollage revêt des formes différentes : sortir de la pauvreté pour les PPTE, échapper à la stagnation pour les PRI, échapper à la rente pour les pétromonarchies. La catégorie est donc utile pédagogiquement, mais elle masque une hiérarchie de situations et de trajectoires qu'indicateurs chiffrés (PIB/hab, IDH, dette) et analyse des stratégies (diversification, insertion dans les flux mondiaux) permettent seuls de restituer.

\textbf{Barème indicatif (sur 20) :} question 1 : 4 pts (classement + 2 indicateurs par pays) ; question 2 : 4 pts ; question 3 : 4 pts ; synthèse : 8 pts (introduction et problématique 1 pt, exploitation des trois documents 3 pts, connaissances personnelles 2 pts, organisation et conclusion 2 pts).

\textbf{Points de vigilance :} (a) chaque classement doit être justifié par \emph{deux} indicateurs chiffrés tirés du tableau ; (b) dans la synthèse, citer explicitement les documents (« le Document 2 montre que\ldots ») ; (c) ne pas opposer les pays de façon manichéenne : le Bangladesh montre qu'un PRI peut émerger, ce qui nuance le propos.
\end{enumerate}
\end{corrige}

\begin{corrige}[Exercice 10 — Sujet type Bac : Croquis « Les foyers de la mondialisation » (TP 5)]
\textbf{1. Éléments attendus sur le croquis :}

\begin{popfigure}
\includegraphics[width=\linewidth,height=9cm,keepaspectratio]{N15/shipping_routes.jpg}
\legende{Fond de carte du croquis-type : les flux de marchandises (carte de la densité du trafic maritime commercial : les continents apparaissent en noir ; plus le bleu est foncé, plus le trafic est dense). On y lit les grandes routes de la mondialisation : Atlantique Nord, Pacifique, route Europe--Asie par Suez et le détroit de Malacca, route du Cap, et la concentration du trafic autour des pôles de puissance (Amérique du Nord, Europe, Asie orientale). Sur cette base, le croquis attendu situe : 1. les pôles de puissance (Triade) ; 2. les espaces émergents (BRICS) ; 3. les ports-hubs et les flux financiers, que la carte ne montre pas. \\ Source : Wikimedia Commons, T. Hengl, CC BY-SA 3.0.}
\end{popfigure}

\textbf{2. Légende organisée (modèle) :}
\begin{itemize}
\item \textbf{Item 1 — Pôles de puissance :} les trois pôles de la Triade (Amérique du Nord ; Europe occidentale ; Asie orientale : Japon, Corée du Sud, Taïwan, façade est chinoise) — cadres en pointillés.
\item \textbf{Item 2 — Espaces émergents :} BRICS et assimilés (Chine, Inde, Brésil, Russie, Afrique du Sud) — points orange.
\item \textbf{Item 3 — Flux et nœuds :} flux majeurs de marchandises (flèches pleines d'épaisseur proportionnelle : transpacifique, transatlantique, Asie--Europe par Suez, intra-asiatique) ; flux financiers/IDE (flèches en pointillés : Triade $\rightarrow$ émergents ; Sud-Sud, Chine $\rightarrow$ Afrique, fonds du Golfe) ; nœuds portuaires et aéroportuaires (Shanghai, Singapour, Rotterdam, Los Angeles, Jebel Ali/Dubaï ; hubs de Dubaï, Londres, Hong Kong).
\end{itemize}

\textbf{3. Commentaire (proposition, 10 lignes) :}

Ce croquis met en évidence la \emph{multipolarisation} de la mondialisation : à côté des trois pôles historiques de la Triade, qui concentrent encore l'essentiel de la richesse, des IDE et des flux transatlantique et transpacifique, s'affirment des puissances émergentes — la Chine d'abord, premier exportateur mondial de marchandises, mais aussi l'Inde, le Brésil, la Russie et l'Afrique du Sud réunies dans les BRICS. Ces pays ne sont plus de simples périphéries fournisseuses de matières premières : ils sont devenus des ateliers, des marchés et des investisseurs, comme le montre le flux d'IDE Sud-Sud de la Chine vers l'Afrique (infrastructures, mines) et les placements des fonds souverains du Golfe dans le monde entier. Les grands nœuds logistiques (Shanghai, Singapour, Dubaï) confirment ce basculement du centre de gravité des échanges vers l'Asie. La mondialisation apparaît ainsi comme un système multipolaire et réticulaire, où les émergents jouent un rôle charnière entre le Nord et le Sud, même si la Triade conserve le commandement financier et technologique.

\textbf{Barème indicatif (sur 20) :} croquis : 12 pts (Triade 3 pts ; 5 émergents 2 pts ; 4 flux de marchandises proportionnels 4 pts ; 2 flux financiers 2 pts ; ports et aéroports 1 pt) ; légende en 3 items : 4 pts ; commentaire : 4 pts.

\textbf{Points de vigilance :} (a) l'épaisseur des flèches doit être hiérarchisée (transpacifique $>$ Asie--Europe $>$ transatlantique $>$ intra-asiatique, ordres de grandeur admis) ; (b) ne pas placer la Chine entière dans la Triade : seule sa façade orientale participe du pôle asiatique, la Chine étant classée parmi les émergents ; (c) tout symbole doit figurer dans la légende ; (d) le commentaire doit mobiliser le vocabulaire du cours : Triade, émergents, multipolarisation, IDE, nœuds, flux.
\end{corrige}

\begin{corrige}[Exercice 11 — Sujet type Bac : Développement construit — Le Cameroun face aux défis de l'émergence]
\textbf{Proposition de corrigé rédigé (plan en trois parties) :}

\textbf{Introduction.} Classé parmi les pays à revenus intermédiaires de la tranche inférieure (PIB/hab d'environ 1\,600 \$), doté d'un IDH de 0,587 et engagé dans la Vision 2035 qui ambitionne d'en faire un pays émergent, le Cameroun illustre la figure du « décollage difficile » : une croissance réelle mais insuffisante pour transformer en profondeur l'économie et réduire une pauvreté qui touche encore environ 37\% de la population. Comment un pays aux atouts reconnus peine-t-il à enclencher une dynamique d'émergence durable ? Nous verrons successivement ses atouts sous-exploités, les freins structurels qui le maintiennent dans le « piège du revenu intermédiaire », puis les leviers d'action publique susceptibles d'enclencher la transformation structurelle.

\textbf{I. Des atouts réels mais sous-exploités.} « Afrique en miniature », le Cameroun cumule des ressources exceptionnellement diversifiées : agricoles (cacao — 5\textsuperscript{e} producteur mondial —, café, coton, bananes, palmier à huile), forestières (deuxième masse forestière d'Afrique), minières (bauxite, fer, cobalt, rutile encore peu exploités), énergétiques (pétrole, gaz, immense potentiel hydroélectrique) et touristiques. Sa position est stratégique : ouvert sur le golfe de Guinée, il est la porte d'entrée de l'Afrique centrale, avec les corridors Douala--N'Djaména et Douala--Bangui qui desservent le Tchad et la Centrafrique enclavés. Enfin, sa démographie jeune (près de 28 millions d'habitants, dont la majorité a moins de 25 ans) et son adhésion à la ZLECAf ouvrent l'accès à un marché continental de 1,4 milliard de consommateurs. Mais ces atouts restent un potentiel : la richesse repose sur l'exportation de produits bruts, sans transformation locale significative.

\textbf{II. Des freins structurels persistants : le piège du revenu intermédiaire.} Le premier frein est la dépendance aux matières premières : le pétrole brut fournit environ 40\% des exportations, le bois 15\%, le cacao-café-coton 15\%, l'aluminium 10\% — des produits à faible valeur ajoutée locale, dont les cours fluctuants exposent l'économie aux chocs externes. Le deuxième frein est le déficit d'infrastructures : l'énergie reste insuffisante et coûteuse malgré les barrages de Memve'ele et de Nachtigal (420 MW), le réseau routier et ferroviaire est dégradé, et les coûts de production et de transport découragent l'investissement industriel. Le troisième frein tient à la gouvernance : corruption persistante, climat des affaires perfectible, évasion fiscale, et surtout un secteur informel qui emploie environ 90\% des actifs et échappe à l'impôt — d'où une pression fiscale inférieure à 15\% du PIB, insuffisante pour financer les services publics. Enfin, le capital humain reste fragile : IDH de 0,587, inadéquation entre formation et emploi, systèmes de santé et d'éducation à renforcer. La dette publique (environ 45\% du PIB, risque de surendettement jugé élevé mais soutenable par le FMI) réduit en outre la marge de manœuvre budgétaire.

\textbf{III. Leviers d'action et pistes de transformation structurelle.} La Stratégie nationale de développement SND30 (2020--2030) fixe le cap : industrialisation par substitution aux importations et transformation locale des matières premières — usines de transformation du cacao (Atlantic Cocoa à Kribi), bois de deuxième et troisième transformation, filière coton-textile, projet d'alumine. Les infrastructures structurantes doivent être achevées et entretenues : barrages de Nachtigal et Memve'ele, extension du port en eau profonde de Kribi (déjà tête de pont des exportations), autoroute Yaoundé--Douala, couverture numérique, en s'appuyant sur les partenariats public-privé. L'amélioration de la gouvernance est décisive : lutte effective contre la corruption, sécurisation du climat des affaires, élargissement de l'assiette fiscale par la formalisation de l'informel. L'intégration régionale (CEMAC réformée, ZLECAf) peut élargir les débouchés et insérer le pays dans des chaînes de valeur africaines. Enfin, l'investissement dans le capital humain — enseignement technique et professionnel, santé, protection sociale — conditionne la productivité de long terme.

\textbf{Conclusion.} Le décollage camerounais n'est pas une fatalité différée : le Bangladesh ou le Vietnam, au niveau de revenu du Cameroun il y a vingt-cinq ans, ont émergé par l'industrie exportatrice et l'investissement massif dans l'éducation. Mais l'émergence suppose une rupture : passer d'une économie de rente et d'exportation de bruts à un modèle industriel, inclusif et durable, ce qui exige d'abord un progrès de gouvernance. La diaspora, les financements verts et la ZLECAf peuvent y contribuer, à condition que l'État transforme les atouts en avantages compétitifs.

\textbf{Barème indicatif (sur 20) :} introduction (accroche, définitions, problématique, plan) : 3 pts ; partie I (atouts, au moins 3 arguments illustrés) : 4 pts ; partie II (freins, au moins 3 arguments chiffrés) : 5 pts ; partie III (leviers, au moins 3 pistes avec exemples) : 5 pts ; conclusion nuancée avec ouverture : 2 pts ; qualité de la langue et de l'organisation : 1 pt.

\textbf{Points de vigilance :} (a) respecter la consigne de discussion : il faut \emph{nuancer} l'affirmation, pas seulement l'illustrer ; (b) mobiliser les données du cours (PIB/hab, IDH, parts des exportations, dette, SND30, ZLECAf) ; (c) éviter le catalogue : chaque partie doit dégager une idée directrice annoncée en tête de paragraphe ; (d) les comparaisons internationales (Bangladesh, Vietnam) valorisent la copie.
\end{corrige}

\newpage

\coursec{Fiche bilan}

\begin{popfigure}
\centering
\begin{tikzpicture}[mindmap, grow cyclic, every node/.style={concept, text=white, font=\small\bfseries},
root concept/.append style={concept color=popink, font=\large\bfseries},
level 1/.append style={level distance=4.6cm, sibling angle=60},
level 2/.append style={level distance=2.6cm, sibling angle=45, font=\scriptsize\bfseries}]
\node[root concept]{Pays au décollage économique difficile}
  child[concept color=popgold]{ node {Pays du Golfe}
    child { node {Rente pétrolière} }
    child { node {Dépendance aux hydrocarbures} }
  }
  child[concept color=popgreen]{ node {Pays d'Afrique}
    child { node {Décollage freiné} }
    child { node {Matières premières brutes} }
  }
  child[concept color=poporange]{ node {PPTE}
    child { node {Dette insoutenable} }
    child { node {Aide et allègement} }
  }
  child[concept color=popblue]{ node {Revenus intermédiaires}
    child { node {Émergence inégale} }
    child { node {Piège du revenu intermédiaire} }
  }
  child[concept color=poppurple]{ node {Mondialisation}
    child { node {Triade et émergents} }
    child { node {Marginalisation des Suds} }
  }
  child[concept color=poppink]{ node {Indicateurs}
    child { node {IDH, PIB/hab.} }
    child { node {Dette, pauvreté} }
  };
\end{tikzpicture}
\legende{Carte mentale de la leçon : les catégories de pays en difficulté dans la mondialisation, leurs caractéristiques et les indicateurs qui permettent de les identifier.}
\end{popfigure}

\begin{bilan}
\textbf{Définitions clés à connaître par cœur}
\begin{itemize}
\item \cle{Décollage économique} : passage d'une économie traditionnelle à une croissance auto-entretenue (industrialisation, hausse durable du revenu). Il est dit \emph{difficile} lorsque des obstacles structurels (dette, dépendance, pauvreté) le freinent.
\item \cle{PPTE} (Pays Pauvres Très Endettés) : pays à faible revenu dont la dette extérieure est insoutenable ; environ 37 pays concernés, dont plus de 30 en Afrique subsaharienne (Cameroun admis en 2000, point d'achèvement en 2006).
\item \cle{Pays à revenus intermédiaires} (PRI) : pays dont le PIB par habitant se situe entre environ \num{1000} et \num{12000} dollars (seuils Banque mondiale, ordres de grandeur) ; certains émergent (Chine, Brésil), d'autres stagnent.
\item \cle{Rente pétrolière} : revenu tiré de l'exportation des hydrocarbures, sans transformation locale ; source de richesse mais aussi de dépendance (économie \emph{rentière}).
\item \cle{Triade} : les trois foyers dominants de la mondialisation : Amérique du Nord, Europe occidentale, Asie orientale (Japon, Chine, Corée du Sud).
\end{itemize}

\textbf{Repères chiffrés et classements (ordres de grandeur à retenir)}
\begin{center}
\begin{tabularx}{\linewidth}{|l|X|X|}
\hline
\rowcolor{popblueL} \textbf{Catégorie} & \textbf{Caractéristiques} & \textbf{Exemples} \\
\hline
Pays du Golfe & Fort PIB/hab. grâce au pétrole ; dépendance aux hydrocarbures et à la main-d'œuvre étrangère & Arabie saoudite, Koweït, Qatar, Émirats arabes unis \\
\hline
Pays d'Afrique & Faible industrialisation, exportation de produits bruts, IDH souvent faible & Cameroun, Tchad, Niger, RDC \\
\hline
PPTE & Dette extérieure insoutenable, pauvreté massive, aide internationale & Cameroun, Mali, Haïti, Bolivie \\
\hline
Revenus intermédiaires & Industrialisation en cours, inégalités fortes, risque de stagnation & Brésil, Turquie, Afrique du Sud, Cameroun (PRI inférieur) \\
\hline
\end{tabularx}
\end{center}

\textbf{Méthodes express}
\begin{itemize}
\item \emph{Commenter un tableau statistique} : titre $\rightarrow$ unités $\rightarrow$ valeurs extrêmes $\rightarrow$ calculs (taux de variation, parts en \%) $\rightarrow$ interprétation géographique.
\item \emph{Construire un croquis de flux} : fond de carte simplifié $\rightarrow$ flèches d'épaisseur proportionnelle aux flux $\rightarrow$ légende numérotée $\rightarrow$ titre et orientation.
\item \emph{Classer un pays} : croiser PIB/habitant, IDH, structure des exportations et poids de la dette avant de conclure.
\end{itemize}
\end{bilan}

\begin{aretenir}[Checklist avant le Bac]
\begin{itemize}
\item[$\square$] Je sais définir précisément : décollage économique, PPTE, pays à revenus intermédiaires, rente pétrolière.
\item[$\square$] Je connais les critères d'admission au programme PPTE et la place du Cameroun (initiative PPTE, 2000--2006).
\item[$\square$] Je sais expliquer pourquoi la rente pétrolière des pays du Golfe est une richesse fragile.
\item[$\square$] Je sais citer au moins trois freins au décollage des pays africains (dette, exportations brutes, instabilité, faible industrialisation).
\item[$\square$] Je distingue pays émergents et pays à revenus intermédiaires stagnants, avec un exemple de chaque.
\item[$\square$] Je sais calculer une part en pourcentage et un taux de variation à partir d'un tableau statistique.
\item[$\square$] Je sais construire et légender un croquis de flux (flèches proportionnelles, titre, orientation, échelle indicative).
\item[$\square$] Je sais localiser sur une mappemonde la Triade, les pays émergents (BRICS) et les pays marginalisés.
\item[$\square$] Je connais les principaux flux mondiaux : financiers (IDE), marchandises (pétrole, matières premières), migratoires.
\item[$\square$] Je sais rédiger une conclusion nuancée : la mondialisation intègre inégalement les territoires.
\end{itemize}
\end{aretenir}

\findecours

\end{document}
